\documentclass[a4paper,11pt,reqno]{amsart}


\usepackage{amsmath,amsfonts,amsthm,amssymb}
\usepackage{feynmp}
%\usepackage[latin1]{inputenc}

% To compile the graphs: mpost roughevpics


%% %   Somewhat wider and taller page than in art12.sty
  \topmargin 0.0in  \headsep 0.4in  \textheight 9.0in
   \oddsidemargin 0.02in  \evensidemargin 0.05in
\textwidth 6.2in

\newcommand{\note}[1]{\marginpar{\framebox{\textbf{NB}}}\textbf{#1}}

\DeclareMathOperator{\Dom}{Dom}
\newcommand{\ol}[1]{\overline{#1}}
\newcommand{\ap}{A_{o}}
\newcommand{\cB}{\mathcal{B}}
\newcommand{\cZ}{\mathcal{Z}}
\newcommand{\cY}{\mathcal{Y}}

\newcommand{\cba}{\cB_\alpha}
\newcommand{\cBalpha}{\cB_\alpha}
\newcommand{\chs}{\hat\cac_{\mbox{\tiny HS},1}}
\newcommand{\cou}{\hat\cac_{\mbox{\tiny op},1}}
\newcommand{\dn}{\cs_{n}}
\newcommand{\dnp}{\Delta_{n+1}}
\newcommand{\ddo}{d^{o}}
\newcommand{\ehs}{\hat\ce_{\mbox{\tiny HS},1}}
\newcommand{\eou}{\hat\ce_{\mbox{\tiny op},1}}
\newcommand{\gm}{\gamma_{m}}
\newcommand{\gn}{\gamma_{n}}
\newcommand{\hca}{\hat {\mathcal C}}
\newcommand{\hdd}{\hat d}
\newcommand{\hd}{\hat \delta}
\newcommand{\he}{\hat {\mathcal E}}
\newcommand{\hep}{\hat \varepsilon}
\newcommand{\hl}{\hat \Lambda}
\newcommand{\hm}{\hat M}
\newcommand{\HS}{\mbox{\tiny HS}}
\newcommand{\hz}{\hat z}
\newcommand{\ka}{\kappa}
\newcommand{\ko}{\kappa_{0}}
\newcommand{\kn}{\kappa_{n}}
\newcommand{\kr}{\kappa_{r}}
\newcommand{\iou}{\int_{0}^{1}}
\newcommand{\iot}{\int_{0}^{t}}
\newcommand{\ist}{\int_{s}^{t}}
\newcommand{\nhs}{\|_{\mbox{\tiny HS}}}
\newcommand{\nop}{\|_{\mbox{\tiny op}}}
\newcommand{\ot}{[0,t]}
\newcommand{\ott}{[0,T]}
\newcommand{\ou}{[0,1]}
\newcommand{\rn}{r^{n}}
\newcommand{\rnp}{r^{n+1}}
\newcommand{\sst}{\Sigma^*}
\newcommand{\td}{\tilde \delta}
\newcommand{\tl}{\tilde \Lambda}
\newcommand{\tk}{\tilde K}
\newcommand{\tq}{\tilde{\mathcal Q}}
\newcommand{\1}{{\bf 1}}
\newcommand{\Norm}[2]{\mathcal{N}[{#1};{#2}]}

\DeclareMathOperator{\id}{\text{Id}}
\DeclareMathOperator{\expect}{{E}}
\DeclareMathOperator{\prob}{{\mathbb P}}
\DeclareMathOperator{\var}{Var}
\DeclareMathOperator{\Id}{Id}
\newcommand{\eps}{\varepsilon}
\newcommand{\latt}{{\mathbb{Z}^d}}
%\newcommand{\qed}{\hfill$\square$}
%\newcommand{\proof}{{\sc Proof.}\hspace{0.3cm}}
\makeatletter
\newcommand{\eqcolon}{\mathrel{\mathord{=}\raise.2\p@\hbox{:}}}
\newcommand{\coloneq}{\mathrel{\raise.2\p@\hbox{:}\mathord{=}}}
\makeatother
\newcommand{\der}{\delta}
\newcommand{\XX}{\mathbb{X}}
\newcommand{\RR}{\mathbb{R}}
%\DeclareMathOperator{\NN}{N}
\newcommand{\NN}{N}
%\newcommand{\LL}{\mathcal{L}}

\newcommand{\UU}{\mathcal{U}}
\newcommand{\QQ}{\mathcal{Q}}
\newcommand{\HH}{\mathcal{H}}
\newcommand{\CC}{\mathcal{C}}
\newcommand{\FF}{\mathcal{F}}
\newcommand{\BB}{\mathcal{B}}
\newcommand{\DD}{\mathcal{D}}
\newcommand{\ZZ}{\mathcal{Z}}
\newcommand{\LL}{\mathcal{L}}
\renewcommand{\SS}{\mathcal{S}}
\newcommand{\OC}{\mathcal{C}}
\newcommand{\OB}{\Omega \mathcal{B}}
\newcommand{\OD}{\Omega \mathcal{D}}
\newcommand{\ie}{\emph{i.e.}}
\newcommand{\norm}[1]{\lVert #1\rVert}
%\newcommand{\Norm}[1]{\left\lVert #1\right\rVert}
\newcommand{\abs}[1]{\lvert #1\rvert}
\newcommand{\Abs}[1]{\left\lvert #1\right\rvert}
\newcommand{\wt}{\widetilde}

%%%%%%%%%%%%%%%%%%%%%%%%%%%%%%%%%%%%%%
%%%%%%%%%%%%%%% Mathbb %%%%%%%%%%%%%%%%%%%
%%%%%%%%%%%%%%%%%%%%%%%%%%%%%%%%%%%%%%%%
\newcommand{\C}{\mathbb C}
\newcommand{\D}{\mathbb D}
\newcommand{\R}{\mathbb R}
\newcommand{\N}{\mathbb N}
\newcommand{\Q}{\mathbb Q}
\newcommand{\Z}{\mathbb Z}

%%%%%%%%%%%%%%%%%%%%%%%%%%%%%%%%%%%%%
%%%%%%%%%%%%%%% Calligraphic %%%%%%%%%%%%%%%%
%%%%%%%%%%%%%%%%%%%%%%%%%%%%%%%%%%%%%%
\newcommand{\cb}{\mathcal B}
\newcommand{\cac}{\mathcal C}
\newcommand{\hcac}{\hat {\mathcal C}}
\newcommand{\ce}{\mathcal E}
\newcommand{\cf}{\mathcal F}
\newcommand{\ch}{\mathcal H}
\newcommand{\ci}{\mathcal I}
\newcommand{\cj}{\mathcal J}
\newcommand{\cl}{\mathcal L}
\newcommand{\cn}{\mathcal N}
\newcommand{\cq}{\mathcal Q}
\newcommand{\cs}{\mathcal S}
\newcommand{\cw}{\mathcal W}
\newcommand{\cz}{\mathcal Z}
\newcommand{\crr}{\mathcal R}

%%%%%%%%%%%%%%%%%%%%%%%%%%%%%%%%%%%%%%%%%%
%%%%%%%%%%%%%%% Greek %%%%%%%%%%%%%%%%%%%%%%%%
%%%%%%%%%%%%%%%%%%%%%%%%%%%%%%%%%%%%%%%%%%%
\newcommand{\al}{\alpha}
\newcommand{\ep}{\varepsilon}
\newcommand{\io}{\iota}
\newcommand{\ga}{\gamma}
\newcommand{\la}{\lambda}
\newcommand{\laa}{\Lambda}
\newcommand{\om}{\omega}
\newcommand{\oom}{\Omega}
\newcommand{\ssi}{\Sigma}
\newcommand{\si}{\sigma}
\newcommand{\te}{\theta}
\newcommand{\vp}{\varphi}
\newcommand{\ze}{\zeta}

%%%%%%%%%%%%%%%%%%%%%%%%%%%%%%%%%%%%%%%%%%
%%%%%%%%%%%%%%% Brackets %%%%%%%%%%%%%%%%%%%%%%%
%%%%%%%%%%%%%%%%%%%%%%%%%%%%%%%%%%%%%%%%%%%
\newcommand{\lp}{\left(}
\newcommand{\rp}{\right)}
\newcommand{\lc}{\left[}
\newcommand{\rc}{\right]}
\newcommand{\lcl}{\left\{}
\newcommand{\rcl}{\right\}}
\newcommand{\lln}{\left|}
\newcommand{\rrn}{\right|}
\newcommand{\lla}{\left\langle}
\newcommand{\rra}{\right\rangle}
\newcommand{\lan}{\left\langle}

%%%%%%%%%%%%%%%%%%%%%%%%%%%%%%%%%%%%%%%%%%%
%%%%%%%%%%%%%%% Begin-End %%%%%%%%%%%%%%%%%%%%%%
%%%%%%%%%%%%%%%%%%%%%%%%%%%%%%%%%%%%%%%%%%%
\newcommand{\bean}{\begin{eqnarray*}}
\newcommand{\eean}{\end{eqnarray*}}
\newcommand{\ben}{\begin{enumerate}}
\newcommand{\een}{\end{enumerate}}
\newcommand{\beq}{\begin{equation}}
\newcommand{\eeq}{\end{equation}}



\newtheorem{theorem}{Theorem}[section]
\newtheorem{acknowledgement}[theorem]{Acknowledgement}
\newtheorem{algorithm}[theorem]{Algorithm}
\newtheorem{axiom}[theorem]{Axiom}
\newtheorem{case}[theorem]{Case}
\newtheorem{claim}[theorem]{Claim}
\newtheorem{conclusion}[theorem]{Conclusion}
\newtheorem{condition}[theorem]{Condition}
\newtheorem{conjecture}[theorem]{Conjecture}
\newtheorem{corollary}[theorem]{Corollary}
\newtheorem{criterion}[theorem]{Criterion}
\newtheorem{definition}[theorem]{Definition}
\newtheorem{example}[theorem]{Example}
\newtheorem{exercise}[theorem]{Exercise}
\newtheorem{hypothesis}{Hypothesis}
\newtheorem{lemma}[theorem]{Lemma}
\newtheorem{notation}[theorem]{Notation}
\newtheorem{problem}[theorem]{Problem}
\newtheorem{proposition}[theorem]{Proposition}
\newtheorem{question}{Question}
\newtheorem{reponse}{Reponse}
%\newtheorem{remark}[theorem]{Remark}
\newtheorem{solution}[theorem]{Solution}
\newtheorem{summary}[theorem]{Summary}
\theoremstyle{remark}
\newtheorem{remark}[theorem]{Remark}

\newenvironment{pseudoproof}{{\sc (Pseudo)Proof.}}{\hfill\qed?}




\newcommand{\orho}{\overline{\rho}}
\newcommand{\onu}{\overline{\nu}}
\newcommand{\omu}{\overline{\mu}}
\newcommand{\TT}{\mathcal{T}}
\newcommand{\cA}{\mathcal{A}}
%\newcommand{\troot}{\circ}




\usepackage{pstricks} 
\usepackage{pst-node}
\usepackage{pst-tree}

\newcommand{\tnode}{\TR{\raisebox{0.5pt}{\ensuremath{\bullet}}}} 
%\newcommand{\tnodel}[1]{\TR{\raisebox{0.5pt}{\footnotesize \ensuremath{#1}}}} 
%\newcommand{\tnodel}[1]{\TR{\raisebox{0.5pt}{\ensuremath{\bullet_{#1}}}}} 
\newcommand{\tnodel}[1]{\TR{\makebox[6pt][l]{\raisebox{0.5pt}{\ensuremath{\bullet{\text{\footnotesize
            #1}}}}}}} 
\newcommand{\tnodea}{\TR{\raisebox{0.5pt}{\ensuremath{\star}}}} 
\newcommand{\tnodeb}{\TR{\raisebox{0.5pt}{\ensuremath{\circ}}}} 

\newcommand{\troot}{\bullet}


\newcommand{\tsnode}{\TR{\raisebox{0.5pt}{\ensuremath{\scriptstyle\bullet}}}} 
%\newcommand{\tsroot}{\TR{\raisebox{0.5pt}{{ \ensuremath{\bullet}}}}} 
\newcommand{\tsroot}{\ensuremath{\scriptstyle\bullet}}
\newcommand{\aaabbb}{\pstree{\tsnode}{\pstree{\tsnode}{\tsnode}}} 
\newcommand{\aaaabbbb}{\pstree{\tsnode}{\pstree{\tsnode}{\pstree{\tsnode}{\tsnode}}}} 
\newcommand{\aabb}{\pstree{\tsnode}{\tsnode}} 
\newcommand{\aababb}{\pstree{\tsnode}{\tsnode \tsnode}}
\newcommand{\aaababbb}{\pstree{\tsnode}{\aababb}}
\newcommand{\aabababb}{\pstree{\tsnode}{\tsnode\tsnode\tsnode}}
\newcommand{\aaabbabb}{\pstree{\tsnode}{\tsnode \aabb}} 
\newcommand{\smalltrees}{\psset{levelsep=-5pt,nodesep=-2pt,treesep=1pt}}
\newcommand{\largetrees}{\psset{levelsep=-10pt,nodesep=-4pt,treesep=5pt}
}
\largetrees
\newcommand{\btnode}{\Tdot}
\newcommand{\btv}{\pstree{\btnode}{\btnode \btnode}}
\newcommand{\btvva}{\pstree{\btnode}{\btv \btnode}}
\newcommand{\btvvb}{\pstree{\btnode}{\btnode \btv}}
\newcommand{\btvvv}{\pstree{\btnode}{\btv \btv}}
\newcommand{\btvvva}{\pstree{\btnode}{\btvva \btnode}}
\newcommand{\btvvvb}{\pstree{\btnode}{\btvvb \btnode}}
\newcommand{\btvvvc}{\pstree{\btnode}{\btnode \btvva}}
\newcommand{\btvvvd}{\pstree{\btnode}{\btnode \btvvb}}

\newcommand{\btrees}{\psset{dotsize=1pt, levelsep=-2pt,nodesep=2pt,treesep=3pt}}
\newcommand{\Btrees}{\psset{dotsize=1pt, levelsep=-4pt,nodesep=4pt,treesep=6pt}}

\begin{document}
  \begin{fmffile}{roughevpics}

\title{Rough evolution equations}
\author{Massimiliano Gubinelli \and Samy Tindel}
\date{\today}
\keywords{Rough paths theory; Stochastic PDEs; Fractional Brownian motion.}

\subjclass[2000]{60H05, 60H07, 60G15}

\address{
{\it Massimiliano Gubinelli:}
{\rm Equipe de Probabilit\'es et Statistique,
Universit\'e Paris-Sud,
91405 Orsay, France}.
{\it Email: }{\tt Massimiliano.Gubinelli@math.u-psud.fr}. 
\newline
$\mbox{ }$\hspace{0.1cm}
{\it Samy Tindel:}
{\rm Institut {\'E}lie Cartan Nancy, Universit\'e de Nancy, B.P. 239,
54506 Vand{\oe}uvre-l{\`e}s-Nancy Cedex, France}.
{\it Email: }{\tt tindel@iecn.u-nancy.fr}
}

\begin{abstract}
We show how to generalize Lyons' rough paths theory in order
to give a pathwise meaning to some non-linear
infinite-dimensional
evolution equations associated to an
analytic semigroup and driven by an irregular noise. As an
illustration, we apply the theory to  a class of 1d
SPDEs driven by a space-time fractional Brownian motion.
\end{abstract}

\maketitle


\tableofcontents 

\section{Introduction}
This paper can be seen as part of an ongoing project whose aim is to give
a pathwise definition to stochastic PDEs. Indeed, the rough path
theory \cite{FV,Lej,LyonsBook,Lyons} and its variants
\cite{rough,FP} have now reached a certain level of maturity, leading
to a proper definition of differential equations driven by irregular
signals and in particular by a fractional Brownian motion \cite{CQ}.
Starting from this observation, we have tried in \cite{glt}
to define and solve the following general problem: let $\cb$ be a
separable Banach space, and $A:D(A)\to\cb$ the infinitesimal generator
of an analytical semigroup $\{S_t;t\ge 0  \}$ on $\cb$,
inducing the family $\{\cb_\al;\al\in\R\}$ with
$\cb_\al=D((-A)^{\al})$. Let also $f$ be a function from $\cb$
to $\cl(\cb_{-\al},\cb_{-\al})$ for a given $\al>0$ and $x$ a noisy
input, considered as a function from $\R_+$ to $\cb_{-\al}$. Then,
for $T>0$, consider the equation
\begin{equation}
  \label{eq:eq1}
 dy_t = A y_t dt + f(y_t) dx_t,
\quad t\in\ott,
\end{equation}
with an initial condition $y_0\in\cb$. The main example we have in
mind is the case of the 1-dimensional heat equation in $\ou$, namely
$\cb=L^2(\ou)$, $A=\Delta$ with Dirichlet
boundary conditions, the usual Sobolev spaces $\cb_\al=H_\al=W^{2\al,2}$,
and $x$ a fractional Brownian motion with Hurst parameter $H$ taking
values in $\cb_{-\al}$. Notice in particular that we wish to consider
a noise $x$ which is irregular in both time and space. Then, in
\cite{glt}, we gave a local existence and uniqueness result for
equation (\ref{eq:eq1}), by considering it in its mild form
\begin{equation}
  \label{eq:eq3b}
 y_t = S_t y_0 + \int_0^t S_{t  s} f(y_s) dx_s,
\end{equation}
where we let $S_{ts}=S_{t-s}$ and interpreting the integral in this mild formulation as a Young integral.
Once the equation is set under the form (\ref{eq:eq3b}), the main
problem one is faced with is to quantify the regularization of the
semi-group $S_{t  s}$ on the term $f(y_s) dx_s$, and then to
elaborate the right fixed point argument in order to solve the
equation. The general results of \cite{glt} could be applied
in the case of the stochastic heat equation driven by a fractional
Brownian motion with Hurst parameter $H>1/2$. They should be compared with the reference \cite{Maslo-Nualart}, where a non-linear fractional SPDE is solved thanks to some fractional calculus methods, but where $x$ is a smooth noise in space.


\vspace{0.2cm}

In the current article, we would like to go one step further with respect
to \cite{glt}, and set the basis of a real rough path expansion in
order to define and solve equation (\ref{eq:eq3b}), which would allow
to consider, in the case of the heat equation in $\ou$, a
fractional Brownian motion with Hurst parameter $H\le 1/2$. This task
is quite long and involved, but let us summarize at this point
some of the ideas we have followed:

\vspace{0.2cm}

\noindent
{\bf (1)}
We will recast equation (\ref{eq:eq3b}) in a suitable way for expansions
according to the following simple observation: we have tried to solve
our evolution equation by means of its infinite dimensional setting,
since it allows to consider $x$  and $y$ as functions of a unique
parameter $t\in\ott$, which makes its rough path type analysis easier
(see \cite{g2} and \cite{qt} for a multiparametric setting).
However, when we come to the applications to the heat equation, we will
consider the evolution equation in $\ott\times \ou$ under the form
\begin{equation}\label{mulf}
y(t,\xi)=
\iou G_t(\xi,\eta) y_0(\eta) d\eta
+\iot\iou G_{t-s}(\xi,\eta)\si(y_s(\eta))\, x(ds,d\eta),
\end{equation}
where $G$ stands for the fundamental solution to the heat equation,
$\si:\R\to\R$ is a regular function, and $x(ds,d\eta)$ is understood
as the distributional derivative of a real-valued continuous process
on $\ott\times \ou$. This definition of our equation is of course
equivalent to (\ref{eq:eq3b}) when $f$ is considered as the pointwise
non-linear operator $[f(y_t)](\xi)\equiv\si(y_t(\xi))$. Now, when
written under its multiparametric form (\ref{mulf}), the equation
is also equivalent to
$$
y(t,\xi)=
\iou G_t(\xi,\eta) y_0(\eta) d\eta
+\iot\iou G_{t-s}(\xi,\eta)\, x(ds,d\eta) \si(y_s(\eta)) ,
$$
and it happens that this simple reformulation is much more convenient
for our future expansions than the original one. When we go back to the
original infinite dimensional setting, we can recast (\ref{eq:eq3b})
into
\begin{equation}
  \label{eq:eq3}
 y_t = S_t y_0 + \int_0^t S_{t  s} \, dx_s f(y_s) ,
\end{equation}
where $f$ is now a smooth function from $\cb$ to $\cb$, and $x$ will
be understood as a H\"older-continuous process taking values in a space
of {\it deregularizing} operators from $\cb$ to a distributional space
$\cb_{-\zeta}$ for a certain $\zeta>0$. The product $dx_s f(y_s)$
will then be regularized again by the action of $S_{t  s}$, in
a way which will be quantified later on. Notice that the form
(\ref{eq:eq3}) of our evolution equation is a little unusual in the SPDE
theory, but makes sense in our context.

\vspace{0.2cm}

\noindent
{\bf (2)}
Instead of considering Riemann sums like in \cite{glt} or like in the
original Lyons' theory \cite{LyonsBook}, our analysis will be based
on the theory of {\it generalized differentials}, called $k$-increments,
contained in \cite{rough}. Roughly speaking, this theory is based on
the fact that an elementary operator, called $\delta$, can transform
an integral $\ist dg_u [h_u-h_s]$, seen as a function of the
variables $s$ and $t$, into a finite difference product
$(g_t-g_s)(h_t-h_s)$. Furthermore, under some additional regularity
properties on $g$ and $h$, the operator $\der$ can be inverted,
and its inverse $\laa$, called \emph{sewing map} (from~\cite{FP}), will be the building stone of our
extension of the notion of integral. Notice that, whenever $g$
and $h$ are H\"older-continuous with H\"older exponent $>1/2$,
this extension coincides with the usual Young integral.
When we consider an integral of the form $\ist dg_u \phi(g_u)$
for a H\"older-continuous function $g$ with H\"older exponent
in $(1/3,1/2]$ admitting a Levy area, our definition of integral
also coincides with Lyons' one, as shown in \cite{rough}. In fact,
if the usual rough path theory gives a richer point of view on  the
algebraic structure of the path $x$, it is worth mentioning that
our approach has at least two advantages:
\begin{enumerate}
\item
Once our unusual setting is assimilated, it becomes quite easy
to figure out how a given expansion in terms of $x$ can be leaded.
And indeed, it will become clear throughout the paper, that
the $k$-increments theory provides a tool allowing some natural
computations for our generalized integrals.
\item
The only step where a discretization procedure is needed is the construction
of the $\laa$ map alluded to above, and this avoids some of the cumbersome
calculations which are one of the main ingredients of the rough path theory.
\end{enumerate}
We hope that this paper will advocate for the use of the
$k$-increments theory, which obviously does not exclude the other approaches
\cite{LyonsBook,FP}.

\vspace{0.2cm}

\noindent
{\bf (3)}
The fact that we are dealing with an evolution problem will force us
to change some of the algebraic structure we will rely on, especially
if one wants to take advantage of the regularizing effect of $S_t$.
This will lead us to introduce an operator $a_{ts}=S_{t  s}-\id$ for
$t\ge s$, and a modified $\der$ operator, called $\hd$, defined
by $\hd=\der-a$. The whole increment theory will have to be build again
based on this modified operator, and we will see that it is really suitable
for the evolution setting induced by (\ref{eq:eq3}). In particular, we
will be able to define analogs of the Levy area and of the higher order iterated integrals,
which are of course harder to express than in the finite dimensional
case, but can be written, in the bilinear case
(that is $\si(r)=r$ in (\ref{mulf})), as
\beq\label{eq:intro-def-op-m}
X^2_{ts} = \int_s^t S_{t  u}dx_u \int_s^u S_{u  v}dx_v S_{v  s},
\qquad
X^3_{ts} = \int_s^t S_{t  u}dx_u X^2_{us}, \qquad \text{etc.}
\eeq
Obviously, a  convenient definition of iterated integrals is the key to
reach the case of a H\"older continuous noise of order $\le 1/2$.

\vspace{0.2cm}

\noindent
{\bf (4)}
The whole integration theory can be expressed in an abstract way, by
just supposing a certain set of assumptions on some incremental
operators like $X^2$ and $X^3$. However, we will try to check these assumptions in
some interesting cases, like the infinite-dimensional fractional
Brownian motion for our Young type integration, or the infinite-dimensional
Brownian motion for our step 3 expansion, based on $X^2$ and $X^3$.
Notice that the  rough expansions for the fractional
Brownian motion should be investigated in details too, but one is
faced with an additional problem in this situation: on one hand,
a Stratonovich type integration requires a lot of regularity
in space for the noise, due to the well-known presence of some
trace terms. On the other hand, the Skorokhod integral doesn't
fulfill the algebraic requirements we ask for our integral extension.
A discussion of these problems and some ideas to solve them
will be included at the end of
the paper, but for sake of conciseness, we will postpone
a complete development of this part to a subsequent paper,
and stick here to the Brownian case.




\vspace{0.2cm}



This paper is structured as follows. In Sect.~\ref{sec:one-dim} we
recall the basic setup of~\cite{rough} which allows to embed
the theory of rough paths in a theory of integration of
{\it generalized differentials}, called here
$k$-\emph{increments}. We wrote it with the aim of having a
self-contained and pedagogical introduction to the topic.
However we give also a new and very elementary proof of the
existence of the basic integration  map $\Lambda$ of~\cite{rough}.
In Sect.~\ref{sec:3} we
introduce and study a modified coboundary
induced by the operator $\hd$ on  the complex of increments, using
the additional data provided by an analytic semigroup $S$,
 in such a
way that the new complex can be shown to act simply on convolution
integrals of the form appearing in eq.~(\ref{eq:eq3}) and on their
iterated versions. This new complex maintain many of the properties of
the original complex (e.g. its cohomology is trivial) and it is shown
that when equipped with H\"older-like norms which measures ``smallness'' of
the increments, it admits a map, called $\hat \Lambda$ here, which is
the main tool for building an integration (or better, convolution) theory over
those $1$-increments which are good enough (again, in a suitable
sense, to be specified in due time).
A key feature of this perturbed complex is that, due to the
convolution with the semigroup $S$, ``space'' and ``time''
regularity of increments depends on each other: we can gain space
regularity by loosing some time regularity and vice-versa. This property will be
essential for the solution of the evolution problem by fixed-point arguments.
In Sect.~\ref{sec:young} we use the theory outlined in
Sect.~\ref{sect:complex} to define the convolution integral in the Young
sense and solve a class of  non-linear evolution problems, reobtaining
some results of the work~\cite{glt}. Notice that we will also improve some of 
our previous results contained in~\cite{glt}, in the sense that we will be able to construct global solutions to our evolution equations
in the Young context. 
In Sect.~\ref{sec:rough} we study the (bi)-linear evolution problem
\begin{equation}
  \label{eq:eq3lin}
 y_t = S_t y_0 + \int_0^t S_{t  s} dx_s y_s.
\end{equation}
 We
 will also introduce a notion of rough-path suitable for noises driving
 evolution equations. By exploiting this path-wise technique we are
 able to obtain automatically the flow semigroup of the equation and
 we will show how to express this semigroup as a convergent series of
 iterated-integrals which are the lift of the step-3 rough path used
 in the construction of the solution.
Then in Sect.~\ref{sec:rough-nonlin} we turn to a non-linear case of evolution
system, namely the case of the quadratic type equation
\begin{equation*}
%  \label{eq:eq3nonlin}
 y_t = S_t y_0 + \int_0^t S_{t  s} dx_s B(y_s\otimes y_s),
\end{equation*}
where $B$ stands for the pointwise multiplication of functions. This requires the additional
 careful introduction of a collection of a priori increments indexed by planar trees, and an 
 associate notion of controlled path.
Finally, all our results will be applied in the concrete case of the stochastic heat equation
on the circle, in a setting recalled at Sect.~\ref{sec:circle}. The case of a
fractional Brownian case is handled in the special situation of the Young theory, while we stick
to the example of an infinite-dimensional Brownian motion in the rougher situation.
We build the rough path associated to this latter noise and
provide concrete conditions where the theory outlined in the previous
sections can be fruitfully applied. A systematic study of the regularity properties of the incremental operators defined as $X^2$ or $X^3$ in (\ref{eq:intro-def-op-m}) will also be 
provided at Sections~\ref{sec:diagrammatica} and \ref{sec:complex-graphs}, thanks
to some Feynman diagrams techniques.



\section{Algebraic integration in one dimension}
\label{sec:one-dim}



The integration theory introduced in \cite{rough} is based on an
algebraic structure, which turns out to be useful for computational
purposes, but has also its own interest. Since this setting is quite
non-standard, compared with the one developed in~\cite{LyonsBook},
and since we will elaborate on it throughout the paper, we will recall
briefly here its main features. We also provide an elementary proof
of the existence of the $\laa$ map.

\subsection{Increments}\label{incr}

As mentioned in the introduction, the extended integral we deal
with is based on the notion of increment, together with an
elementary operator $\der$ acting on them. However, this simple
structure gives raise to a nice topological structure that we will
describe briefly here: first of all,  for an arbitrary real number
$T>0$, a vector space $V$, and an integer $k\ge 1$, we denote by
$\CC_k(V)$ the set of functions $g : [0,T]^{k} \to V$ such
that $g_{t_1 \cdots t_{k}} = 0$
whenever $t_i = t_{i+1}$ for some $i\le k-1$.
Such a function will be called a
\emph{$(k-1)$-increment}, and we will
set $\CC_*(V)=\cup_{k\ge 1}\CC_k(V)$. The operator $\der$
alluded to above can be seen as a coboundary operator acting on
$k$-increments, inducing a cochain complex $(\CC_*,\delta)$,
and is defined as follows on $\CC_k(V)$:
\begin{equation}
  \label{eq:coboundary}
\delta : \CC_k(V) \to \CC_{k+1}(V) \qquad
(\delta g)_{t_1 \cdots t_{k+1}} = \sum_{i=1}^{k+1} (-1)^i g_{t_1
  \cdots \hat t_i \cdots t_{k+1}} ,
\end{equation}
where $\hat t_i$ means that this particular argument is omitted.
Then a fundamental property of $\der$, which is easily verified,
is that
$\delta \delta = 0$, where $\delta \delta$ is considered as an operator
from $\CC_k(V)$ to $\CC_{k+2}(V)$.
 We will denote $\cZ\CC_k(V) = \CC_k(V) \cap \text{Ker}\delta$
and $\cB \CC_k(V) :=
\CC_k(V) \cap \text{Im}\delta$, respectively the spaces of
\emph{$k$-cocycles} and of \emph{$k$-coboundaries}, following standard
conventions of homological
algebra.

\vspace{0.2cm}

Some simple examples of actions of $\der$,
which will be the ones we will really use throughout the paper,
 are obtained by letting
$g\in\CC_1$ and $h\in\CC_2$. Then, for any $t,u,s\in\ott$, we have
\begin{equation}
\label{eq:simple_application}
  (\der g)_{ts} = g_t - g_s,
\quad\mbox{ and }\quad
(\der h)_{tus} = h_{ts}-h_{tu}-h_{us}.
\end{equation}
Furthermore, it is readily checked that
the complex $(\CC_*,\delta)$ is \emph{acyclic}, i.e.
$\cZ \CC_{k+1}(V) = \cB \CC_{k}(V)$ for any $k\ge 1$, or otherwise stated,
the sequence
\begin{equation}
\label{eq:exact_sequence}
 0 \rightarrow \RR \rightarrow \CC_1(V)
 \stackrel{\der}{\longrightarrow} \CC_2(V) \stackrel{\der}{\longrightarrow} \CC_3(V) \stackrel{\der}{\longrightarrow} \CC_4(V) \rightarrow \cdots
\end{equation}
is exact. In particular, the following basic property, which we
label  for further use, holds true:
\begin{lemma}\label{exd}
Let $k\ge 1$ and $h\in \cz\cac_{k+1}(V)$. Then there exists a (non unique)
$f\in\cac_{k}(V)$ such that $h=\der f$
\end{lemma}

\begin{proof}
This elementary proof is included in \cite{rough}, see also Prop.~\ref{acyc} below. Let us just mention that
$f_{t_1\ldots t_{k}}=h_{t_1\ldots t_{k}0}$ is a possible choice.
\end{proof}

\vspace{0.2cm}

\begin{remark}

Observe that Lemma \ref{exd} implies that all the elements
$h \in\CC_2(V)$ such that $\der h= 0$ can be written as $h = \der f$
for some (non unique) $f \in \CC_1(V)$. Thus we get a heuristic
interpretation of $\der |_{\CC_2(V)}$:  it measures how much a
given 1-increment  is far from being an {\it exact} increment of a
function (i.e. a finite difference).  
\end{remark}

\vspace{0.2cm}

Notice that our future discussions will mainly rely on
$k$-increments with $k \le 2$, for which we will use some
analytical assumptions. Namely,
we measure the size of these increments by H\"older norms
defined in the following way: for $f \in \CC_2(V)$ let
$$
\norm{f}_{\mu} \equiv
\sup_{s,t\in\ott}\frac{|f_{ts}|}{|t-s|^\mu},
\quad\mbox{and}\quad
\CC_1^\mu(V)=\lcl f \in \CC_2(V);\, \norm{f}_{\mu}<\infty  \rcl.
$$
In the same way, for $h \in \CC_3(V)$, set
\begin{eqnarray}
  \label{eq:normOCC2}
  \norm{h}_{\gamma,\rho} &=& \sup_{s,u,t\in\ott}
\frac{|h_{tus}|}{|u-s|^\gamma |t-u|^\rho}\\
%\quad\mbox{and}\quad
\|h\|_\mu &\equiv &
\inf\left \{\sum_i \|h_i\|_{\rho_i,\mu-\rho_i} ;\, h =
 \sum_i h_i,\, 0 < \rho_i < \mu \right\} ,\nonumber
\end{eqnarray}
where the last infimum is taken over all sequences $\{h_i \in \CC_3(V) \}$ such that $h
= \sum_i h_i$ and for all choices of the numbers $\rho_i \in (0,z)$.
Then  $\|\cdot\|_\mu$ is easily seen to be a norm on $\CC_3(V)$, and we set
$$
\CC_3^\mu(V):=\lcl h\in\CC_3(V);\, \|h\|_\mu<\infty \rcl.
$$
Eventually,
let $\CC_3^{1+}(V) = \cup_{\mu > 1} \CC_3^\mu(V)$,
and remark that the same kind of norms can be considered on the
spaces $\cZ \CC_3(V)$, leading to the definition of some spaces
$\cZ \CC_3^\mu(V)$ and $\cZ \CC_3^{1+}(V)$.

\vspace{0.2cm}

With these notations in mind,
the following proposition is a basic result which is at the core of
 our approach to path-wise integration:
\begin{proposition}[The sewing map $\Lambda$]
\label{prop:Lambda}
There exists a unique linear map $\Lambda: \cZ \CC^{1+}_3(V)
\to \CC_2^{1+}(V)$ (the \emph{sewing} map) such that
$$
\delta \Lambda  = \id_{\cZ \CC_3(V)}.
$$
Furthermore, for any $\mu > 1$,
this map is continuous from $\cZ \CC^{\mu}_3(V)$
to $\CC_2^{\mu}(V)$ and we have
\begin{equation}\label{ineqla}
\|\Lambda h\|_{\mu} \le \frac{1}{2^\mu-2} \|h\|_{\mu} ,\qquad h \in  \cZ \CC^{1+}_3(V).
\end{equation}
\end{proposition}

\begin{proof}
For sake of completeness, we include a proof of this result here, which is more elementary than the one provided in \cite{rough}, and
which will be generalized at Theorem \ref{th:conv-lambda}.
For notational sake, we will omit the dependence in $V$ in
our functional spaces, and write for instance $\cac_3$ instead
of $\cac_3(V)$.
Let then $h$ be an element of
$\cZ \CC_3^\mu \subset \cZ \CC_3^{1+}$ for some $\mu > 1$.

\vspace{0.2cm}

\noindent
{\it Step 1:}
Let us first prove the uniqueness of the 1-increment $M\in\CC_2^\mu$ such that
$\der M=h$. Indeed, let $M,\hm$ be two elements of $\CC_2^\mu$ satisfying
$\der M=\der\hm=H$ and set $Q = M -\hm$. Then $\der Q = 0$ and $Q \in
\CC_2^\mu$. Invoking Lemma \ref{exd}, there exists an element
$q \in \CC_1$ such that $Q = \der q$, but since $\mu >
1$, $q$ is a function on $[0,T]$ with zero derivative, i.e. a constant
and then $Q = 0$.

\vspace{0.2cm}


\noindent
{\it Step 2:}
Let us construct now a process $M\in\CC_2^\mu$, with $\mu>1$, satisfying
$\der M=h$. Since $\der h=0$, invoking again Lemma \ref{exd}, we know
that there exists a $B\in\CC_2$ such that $\der B=h$.
Pick $s,t\in\ott$, such that $s<t$ in order to fix ideas, and for $n\ge 0$, consider the
dyadic partition $\{\rn_i; i\le 2^n  \}$ of the interval $[s,t]$, where
\begin{equation}\label{defpartin}
\rn_i=s+\frac{(t-s)i}{2^n},
\quad\mbox{for}\quad i\le 2^n.
\end{equation}
Then, for $n\ge 0$  set
\begin{equation}\label{defmtsn}
M_{ts}^{n}=B_{ts}-\sum_{i=0}^{2^n-1} B_{\rn_{i+1},\rn_{i}}.
\end{equation}
Then it is readily checked that $M_{ts}^{0}=0$. Furthermore
 we have
\begin{equation*}
  \begin{split}
M_{ts}^{n}-M_{ts}^{n+1}
& =
\sum_{i=0}^{2^n-1}
\lp B_{\rnp_{2i+2},\rnp_{2i}}-B_{\rnp_{2i+1},\rnp_{2i}}
-B_{\rnp_{2i+2},\rnp_{2i+1}}\rp
\\& = \sum_{i=0}^{2^n-1} \lp \der B \rp_{\rnp_{2i+2},\rnp_{2i+1},\rnp_{2i}}
=\sum_{i=0}^{2^n-1} h_{\rnp_{2i+2},\rnp_{2i+1},\rnp_{2i}},
%\\&= \sum_{i=0}^{2^n-1} h_{\rnp_{2i+2},\rnp_{2i+1},\rnp_{2i}},
  \end{split}
\end{equation*}
and since $h\in\CC_3^\mu$ with $\mu>1$, we obtain
$$
\lln M_{ts}^{n}-M_{ts}^{n+1} \rrn
\le
\frac{\| h \|_{\mu}(t-s)^\mu}{2^{n(\mu-1)}},
$$
which yields that $M_{ts}\equiv\lim_{n\to\infty}M_{ts}^{n}$ exists, and
satisfies inequality (\ref{ineqla}).

\vspace{0.2cm}

\noindent
{\it Step 3:} Let us consider now a general sequence $\{\pi_n;\, n\ge 1\}$ of partitions
$\{r_0^n, r_1^n,\ldots,r_{k_n}^n,$ $r_{k_n+1}^n\}$
of $[s,t]$, with $s=r_0^n<r_1^n<\ldots<r_{k_n}^n<r_{k_n+1}^n=t$.
We assume that $\pi_{n}\subset\pi_{n+1}$, and $\lim_{n\to\infty}k_n=\infty$. Set
\begin{equation}\label{eq:def-mst-pi-n}
M_{ts}^{\pi_n}=B_{ts}-\sum_{l=0}^{k_n} B_{\rn_{l+1},\rn_{l}}.
\end{equation}
It is easily seen that there exists $1\le l\le k_n$ such that 
\begin{equation}
\label{eq:min-dist}
|r_{l+1}^n-r_{l-1}^n| \le \frac{2 |t-s|}{k_n}
\end{equation}
Pick now such an index $l$, and let us transform $\pi^n$ into $\hat\pi$, where
$$
\hat\pi= \lcl r_0^n, r_1^n,\ldots,r_{l-1}^n,r_{l+1}^n, \ldots,r_{k_n}^n, r_{k_n+1}^n \rcl.
$$
Then, as in the previous step,
$$
M_{ts}^{\hat\pi}=M_{ts}^{\pi_n} - (\der B)_{r_{l+1}^n, r_{l}^n,r_{l-1}^n}
= M_{st}^{\pi_n} - h_{r_{l+1}^n, r_{l}^n,r_{l-1}^n},
$$
using the definition of the space $\cac_3^\mu$ and the bound~(\ref{eq:min-dist}) we have
$$
\lln M_{ts}^{\hat\pi}-M_{ts}^{\pi_n} \rrn
\le 2^{\mu} \| h\|_{\mu} \lp \frac{t-s}{k_n} \rp^\mu .
$$
Repeating now this operation until we end up with the trivial partition
$\hat\pi_0\equiv\{s,t\}$, for which $M_{st}^{\hat\pi_0}=0$, we obtain:
$$
\lln M_{ts}^{\pi_n} \rrn \le 
2^{\mu} \| h\|_{\mu} |t-s|^\mu \sum_{j=1}^{k_n} j^{-\mu} \le
2^{\mu} \| h\|_{\mu} |t-s|^\mu \sum_{j=1}^{\infty} j^{-\mu} \equiv
c_{\mu,h} |t-s|^{\mu}.
$$
Hence, there exists a subsequence $\{\pi_m;\, m\ge 1\}$ of $\{\pi_n;\, n\ge 1\}$
such that $M_{ts}^{\pi_m}$ converges to an element $M_{ts}$, satisfying
$M_{ts}\le c_{\mu,h} |t-s|^{\mu}$. With the same considerations as in \cite{Lej},
it can also be checked that that the limit $M$ does not depend on the particular
sequence of partitions we have chosen, and thus coincides with the one constructed at Step 2.


\vspace{0.2cm}

\noindent
{\it Step 4:}
It remains to show that $\der M=h$. Consider then $0\le s<u<t\le T$, and two sequences
of partitions $\pi_{us}^n$ and $\pi_{tu}^n$ of $[s,u]$ and $[u,t]$ respectively, 
whose meshes tend to 0 as $n\to\infty$. Set also $\pi_{ts}^n=\pi_{tu}^n\cup\pi_{us}^n$.
From the previous step, one can construct easily some subsequences 
$\pi_{tu}^m,\pi_{us}^m,\pi_{ts}^m$, with $\pi_{ts}^m=\pi_{tu}^m\cup\pi_{us}^m$, such that
$$
\lim_{m\to\infty}M_{tu}^{\pi_{tu}^m}= M_{tu}, \quad
\lim_{m\to\infty}M_{us}^{\pi_{us}^m}= M_{us}, \quad
\lim_{m\to\infty}M_{ts}^{\pi_{ts}^m}= M_{ts}.
$$ 
Call now $k_{ts}^m$ (resp. $k_{tu}^m,\, k_{us}^m)$ the number of points of the partition
$\pi_{ts}^m$ (resp. $\pi_{tu}^m,\, \pi_{us}^m$). Then, 
a direct computation, using definition  (\ref{eq:def-mst-pi-n}), shows that
for any $0 \le i \le 2^n$ we have:
\begin{align*}
&M^{\pi_{ts}^m}_{ts} -
M^{\pi_{su}^m}_{tu} - M^{\pi_{ut}^m}_{us}  \\
=& (\der B)_{tus}- 
\lp \sum_{l=0}^{k_{ts}^m+k_{us}^m+1} B_{r_{l+1}^{m} r_{l}^{m}} 
- \sum_{l=0}^{k_{tu}^m} B_{r_{l+1}^{m} r_{l}^{m}}
- \sum_{l=k_{tu}^m+1}^{k_{tu}^m+k_{us}^m+1} B_{r_{l}^{m} r_{l+1}^{m}} \rp  
=(\der B)_{tus}=h_{tus}.
\end{align*}
Taking the limit $m \to \infty$ in the latter relation, we get $(\der M)_{tus}=h_{tus}$,
which ends the proof.


\end{proof}

\vspace{0.2cm}

We can now give an algorithm for a canonical decomposition of the
preimage of the space $\cZ \CC_3^{1+}(V)$, or in other words, of a function
$g\in\CC_2 (V)$ whose increment $\der g$ is smooth enough:
\begin{corollary}
Take an element $g\in\CC_2 (V)$, such that $\der g\in\CC_3^\mu(V)$ for
$\mu>1$. Then $g$ can be decomposed in a unique way as
$$
g=\der f+ \Lambda \delta g,
$$
where $f\in\CC_1(V)$.
\end{corollary}

\begin{proof}
Elementary, see~\cite{rough}.
\end{proof}

At this point the connection of the structure we introduced with
the problem of integration of irregular functions can be still quite
obscure to the non-initiated reader. However something interesting is
already going on and the previous corollary has a very nice
consequence which is the subject of the following property.

\begin{corollary}[Integration of small increments]
\label{cor:integration}
For any 1-increment $g\in\CC_2 (V)$, such that $\der g\in\CC_3^{1+}$,
set
$
\delta f = (\id-\Lambda \delta) g
$.
Then
$$
(\delta f)_{ts} = \lim_{|\Pi_{ts}| \to 0} \sum_{i=0}^n g_{t_{i+1}\, t_i},
$$
where the limit is over any partition $\Pi_{ts} = \{t_0=t,\dots,
t_n=s\}$ of $[t,s]$ whose mesh tends to zero. The
1-increment $\delta f$ is the indefinite integral of the 1-increment $g$.
\end{corollary}

\begin{proof}
Just consider the equation $g = \delta f + \Lambda \delta g$ and write
\begin{equation*}
  \begin{split}
S_\Pi & = \sum_{i=0}^n g_{t_{i+1}\, t_i}
%\\ &
=
\sum_{i=0}^n (\delta f)_{t_{i+1}\, t_i}
+
\sum_{i=0}^n (\Lambda \delta g)_{t_{i+1}\, t_i}
\\ & =
(\delta f)_{ts}
+
\sum_{i=0}^n (\Lambda \delta g)_{t_{i+1}\, t_i}.
      \end{split}
\end{equation*}
Then observe that, due to the fact that $\Lambda \delta g \in
\CC_3^{1+}(V)$, the last sum converges to zero.
\end{proof}

\subsection{Computations in $\CC_*$}\label{cpss}

For sake of simplicity, let us assume, until Section \ref{sec:3},
that $V=\R$, and set $\cac_k(\R)=\cac_k$. Then
the complex $(\CC_*,\delta)$ is an (associative, non-commutative)
graded algebra once endowed with the following product:
for  $g\in\CC_n $ and $h\in\CC_m $ let  $gh \in \CC_{n+m} $
the element defined by
\begin{equation}\label{cvpdt}
(gh)_{t_1,\dots,t_{m+n-1}}=
g_{t_1,\dots,t_{n}} h_{t_{n},\dots,t_{m+n-1}},
\quad
t_1,\dots,t_{m+n+1}\in\ott.
\end{equation}
In this context,
the coboundary $\delta$ act as a graded derivation with respect to the
algebra structure. In particular we have the following useful properties.

\begin{proposition}\label{difrul}
The following differentiation rules hold true:
\begin{enumerate}
\item
Let $g,h$ be two elements of $\CC_1 $. Then
\begin{equation}\label{difrulu}
\der (gh) = \der g\,  h + g\, \der h.
\end{equation}
\item
Let $g \in \CC_1 $ and  $h\in \CC_2 $. Then
$$
\der (gh) = \der g\, h + g \,\der h, \qquad
\der (hg) = \der h\, g  - h \,\der g.
$$
\end{enumerate}
\end{proposition}

\begin{proof}
We will just prove (\ref{difrulu}), the other relations being equally trivial:
if $g,h\in\CC_1 $, then
$$
\lc \der (gh) \rc_{ts}
= g_th_t-g_sh_s
=g_t\lp h_t-h_s \rp +\lp  g_t-g_s\rp h_s\\
=g_t \lp \der h \rp_{ts}+ \lp \der g \rp_{ts} h_s,
$$
which proves our claim.

\end{proof}

\vspace{0.2cm}

The iterated integrals of smooth functions on $\ott$ are obviously
particular cases of elements of $\CC$ which will be of interest for
us, and let us recall  some basic  rules for these objects:
consider $f,g\in\CC_1^\infty $, where $\CC_1^\infty $ is the set of
smooth functions from $\ott$ to $\R$. Then the integral $\int dg \,
f$, which will be denoted by
$\cj(dg \,  f)$, can be considered as an element of
$\CC_2^\infty$. That is, for $s,t\in\ott$, we set
$$
\cj_{ts}(dg \,  f)
=
\left(\int  dg f \right)_{ts} = \int_s^t  dg_u f_u.
$$
The multiple integrals can also be defined in the following way:
given a smooth element $h \in \CC_2^\infty$ and $s,t\in\ott$, we set
$$
\cj_{ts}(dg\, h )\equiv
\left(\int dg h \right)_{ts} = \int_s^t dg_u h_{us} .
$$
In particular, the double integral $\cj_{ts}( df^3df^2\,f^1)$ is defined, for
$f^1,f^2,f^3\in\cac_1^\infty$, as
$$
\cj_{ts}( df^3df^2\,f^1)
=\lp \int df^3df^2\,f^1  \rp_{ts}
= \int_s^t df_u^3 \,\cj_{us}\lp  df^2 \,  f^1 \rp .
$$
and if $f^1,\dots,f^{n+1}\in\cac_0^\infty$, we set
\begin{equation}\label{multintg}
\cj_{ts}(df^{n+1}df^n \cdots df^2 f^1)
=
\int_s^t  df_u^{n+1}\, \cj_{us}\lp df^n\cdots df^2 \,f^1\rp\,,
\end{equation}
which defines the iterated integrals of smooth functions recursively.

%Observe that a $n$th order integral $\cj(df^n\cdots df^2 df^1)$ could be defined along the same lines.

\medskip

The following relations between multiple integrals and the operator $\der$ will also be useful in the remainder of the paper:
\begin{proposition}\label{dissec}
Let $f,g$ be two elements of $\CC_1^\infty$. Then, recalling the convention
(\ref{cvpdt}), it holds that
$$
\der f = \cj( df), \qquad
\der\lp \cj(f dg)\rp = 0, \qquad
\der\lp \cj (dg df)\rp =  (\der g) (\der f) = \cj(dg) \cj(df),
$$
and, in general,
$$
 \der \lp \cj( df^n \cdots df^1)\rp  =
 \sum_{i=1}^{n-1}
\cj\lp df^n \cdots df^{i+1}\rp \cj\lp df^{i}\cdots df^1\rp.
$$
\end{proposition}

\begin{proof}
Here again, the proof is elementary, and we will just show the third of these relations: we have, for $s,t\in\ott,$
$$
\cj_{ts} (dg df)
= \int_s^t dg_u (f_u-f_s)
= \int_s^t dg_u f_u - K_{ts},
$$
with $K_{ts}=(g_t-g_s)f_s$. The first term of the right hand side is easily seen to be in $\cz\cac_2$. Thus
$$
\der\lp \cj (dg df)\rp_{tus} =
-\lp\der K\rp_{tus} = [g_t-g_u][f_u-f_s],
$$
which gives the announced result.

\end{proof}



\subsection{Dissection of an integral}
\label{sec:dissection}
The purpose of this section is not to provide an account on all the computations contained in \cite{rough}. However, we will go into some semi-heuristic considerations that, hopefully, will shed some light on the way we will solve rough PDEs later on: with the notations of Section \ref{cpss} in mind, we will try to give, intuitively speaking, a meaning to the integral
$\int \vp(x) dx=\cj(\vp(x) dx)$ for a non-smooth function $x\in\cac_1$.
Notice that, in the sequel, $x$ should be considered as a vector valued function, since the whole theory can be handled via the Doss-Soussman methodology in the real case. However, we will present the main ideas of the algorithm below as if $x$ were real valued, the generalization from $\R$ to $\R^n$ being just a matter of (cumbersome) notations.

\subsubsection{The Young case}

The first idea one can have in mind in order to define $\cj(dx \vp(x))$ is to
perform an expansion around the increment $dx$: indeed, in the smooth case,
we have
\begin{equation}
  \label{eq:dec-1}
\cj\lp dx \varphi(x) \rp = \der x \, \varphi(x)  + \cj\lp dx d\varphi(x)\rp.
\end{equation}
If we wish to extend the right hand side of (\ref{eq:dec-1}) to a non-smooth case, we see that the first term is harmless, since it is defined independently of the regularity of $x$, by
$$
\lc\der x\, \varphi(x) \rc_{ts}
=
\lc x_t-x_s \rc \vp(x_s)
=
\lc \der x \, \vp(x)  \rc_{ts},
\quad\mbox{ for }\quad s,t\in\ott.
$$
The last term of (\ref{eq:dec-1})
is more problematic and we proceed to its {\it dissection} by
the application of
$\der$: invoking Proposition \ref{dissec}, we get, in the smooth case, that
\begin{equation}
  \label{eq:dec-2}
\der\lp \cj (dx\, d\varphi(x) )\rp =  \der x \der (\varphi(x)),
\quad\mbox{ i.e. }
\lc \der\lp \cj (dx\, d\varphi(x) )\rp\rc_{tus}
= \lc \der x\rc_{tu} \lc\der (\varphi(x))\rc_{us}.
\end{equation}
Now the r.h.s. of (\ref{eq:dec-2}) is well defined independently of
the regularity of $x$.
Thus, if $\der x\,\der (\varphi(x)) \in\cac_3^{1+}$, which happens
when $x\in\cac_1^\al$ with $\al>\frac12$ and $\vp\in C^1(\R)$, then
Proposition \ref{prop:Lambda}
can be applied, and $\Lambda [\der x\,\der( \varphi(x))]$ is defined
unambiguously. Hence, owing to (\ref{eq:dec-2}), we set
$$
\cj (dx d\varphi(x)  )=
\laa\lp \der x\, \der(\varphi(x)) \rp
$$
and
\begin{equation}\label{defyg}
\cj (dx \varphi(x)  )
= \der x \,\vp(x)
+\laa\lp \der x\, \der(\varphi(x)) \rp = (\id-\Lambda \delta) [\delta x \varphi(x)],
\end{equation}
where the last equality is due to Proposition \ref{difrul} and to the
fact that $\der\der x=0$.
Notice once again that this construction is valid whenever $x\in\cac_1^\al$
with $\al>\frac12$ and $\vp\in C^1(\R)$, and it is easily shown, along the same lines as in the proof of Proposition \ref{prop:Lambda} that the integral
$\cj (dx \varphi(x) )$ defined by (\ref{defyg}) corresponds to the usual Young integral.

\subsubsection{Case of a $\al$-H\"older path with $\frac13<\al<\frac12$}


The construction (\ref{defyg}) doesn't work if $x\not\in\cac_1^{1/2+}$. However, if $x\in\cac_1^\al$ with $\al>\frac13$,
we can proceed further in the expansion of
equation (\ref{eq:dec-1}) by observing that, still in the smooth case,
we have, for $s,t\in\ott$,
$$
\int_s^t \lc d\vp(x) \rc_u
=\int_s^t  dx_u \, \vp'(x_u)
= [x_t-x_s] \vp'(x_s)+\int_s^t dx_u \int_s^u dx_v\, \vp''(x_v) ,
$$
or according to the notations of Section \ref{cpss},
\begin{equation}
  \label{eq:dec-3}
\der\vp(x)=
\cj \lp d\varphi(x)\rp  =  \cj \lp dx\, \varphi'(x) \rp
= \der x\, \varphi'(x)
+ \cj \lp dx\, d\varphi'(x) \rp.
\end{equation}
Injecting this equality in equation (\ref{eq:dec-1}), thanks to
(\ref{multintg}), we obtain
\begin{equation}
  \label{eq:dec-4}
\cj \lp dx\,\varphi(x) \rp = \der x \,\varphi(x)
+ \cj\lp dx \,dx \rp\,\varphi'(x)
+\cj\lp dx \, dx \, d\varphi'(x) \rp.
\end{equation}
Let us assume now that we are given a process $\cj(dx \,dx)\in\CC_2$,
usually (and somewhat improperly) called the Levy area of $x$, such that
\begin{equation}\label{dvpt}
\der\lp \cj(dx \,dx) \rp=\der x\, \der x
\quad\mbox{ and }\quad
\cj(dx \,dx)\in \CC_2^{2\al}.
\end{equation}
This assumption is of course not automatically satisfied, but it can be checked for instance in the Brownian and fractional Brownian cases. Then the right hand side of (\ref{dvpt}) is again well defined independently of the regularity of $x$, except for the last term. However, 
recast eq. (\ref{eq:dec-4}) into:
$$
-\cj\lp dx \, dx \, d\varphi'(x) \rp=
-\cj \lp dx\,\varphi(x) \rp + \der x \,\varphi(x)
+ \cj\lp dx \,dx \rp\,\varphi'(x),
$$
and 
apply again $\der$ to both sides of this last expression. Invoking Proposition \ref{dissec} 
and recalling that $\der (\cj(dx\, dx))= \der x\, \der x$, we obtain:
\begin{eqnarray}
  \label{eq:dec-5}
  - \der \cj\lp  dx\, dx d\varphi'(x)\rp
&=&  -\der x\, \der\vp(x) +\der x\, \der x \, \vp'(x)- \cj\lp dx \,dx \rp \,\der\vp'(x) \nonumber \\
&=& -\der x \lc \der\vp(x)- \der x \, \vp'(x)\rc - \cj\lp dx \,dx \rp \,\der\vp'(x). 
\end{eqnarray}
Everything in the r.h.s of equation (\ref{eq:dec-5}) is well defined at this
stage, and if we assume that all the terms belongs to $\OC_3^\mu$ with
$\mu>1$ (which can be justifyed  via Taylor's expansions whenever $x\in\cac_1^\al$
with $\al>\frac13$ and $\vp\in C^2(\R)$), we can conclude that
$$
\cj (dx\varphi(x)  )
= \der x \,\vp(x)
+\cj\lp dx \,dx \rp\, \vp'(x)  
-\laa\Big[ \cj\lp dx \,dx \rp \, \der\vp'(x)
- \der x\,\left( \der \varphi(x) - \der x\, \varphi'(x)  \right) \Big],
$$
or stated otherwise
$$
\cj (dx\varphi(x)  )
= (\id-\Lambda \delta) \left[ \der x \,\vp(x)
+\cj\lp dx \,dx \rp\, \vp'(x) \right],
$$
where we used the fact that $\delta \cj\lp dx \,dx \rp = \delta x
\delta x$ to put in evidence the fact that we are actually integrating
(in the sense of Corollary~\ref{cor:integration}) the 1-increment
$\der x \,\vp(x)
+\cj\lp dx \,dx \rp\, \vp'(x)$ which can be thought of  as a {\it corrected}
version of the more natural integrand $\der x \,\vp(x)$.
It is worth noticing at that point that this  integral has now to
be understood as an integral over the (step-2) rough path $(x, \cj(dx\,
dx))$ introduced in \cite{rough} and it coincides with the notion of  integral
over a rough path given by Lyons in~\cite{Lyons}.

\begin{remark}
  This algorithm has an obvious extension to higher orders if we
assume that a reasonable definition of the iterated integrals
$\cj(dx\, dx\cdots dx)$ can be given. To proceed further however we
need the notion of \emph{geometric} rough path (for more details
on this notion see~\cite{Lyons}) which must be  exploited
crucially to show that some terms are small enough and belong to the domain
of $\Lambda$. For a more general approach which do not rely on geometric rough-path see~\cite{g-ramif}.
\end{remark}



\section{Algebraic integration associated to a semigroup}
\label{sec:3}
The aim of this section is to set the basis for our future computations:
after recalling some basic facts about analytic semi-groups, we will
define a set of increments $\hca_*$ and a modified operator $\hd$
adapted to our evolution setting. Then we will give some basic
calculus rules for $(\hca_*,\hd)$ and eventually, we will fix the
notations for the main application we have chosen, that is the
stochastic heat equation.



\subsection{Analytical semigroups}\label{analsmg}
As in \cite{glt}, we will be able to develop our integration theory in the abstract setting of analytical semigroups on Banach spaces, whose basic features can be summarized as follows:
let $(\cb,|\cdot|)$ be a separable Banach space, and
$(A,\Dom(A))$ be a non-bounded linear operator
on $\cB$. We will assume in the sequel that
(see \cite[Sections~2.5 and 2.6]{paz}) $A$ is the generator of an analytical semigroup $\{S_t;t\ge 0  \}$, satisfying
$$
|S_t|\leq M e^{-\la t},
\quad\mbox{ for some constants } M,\la>0
\mbox{ and for all } t\ge 0,
$$
where $|\cdot|$ also stands for the operator norm on $\cb$.
Set now $\ap=-A$.
This allows in particular to define the fractional powers
$(\ap^\al,\Dom(\ap^\al))$ for any $\al\in\R$.

\vspace{0.2cm}

For $\alpha\geq 0$, let $\cb_{\al}$ be the space
$\Dom(\ap^\alpha)$ with the norm
$|x|_{\cb_{\al}}=|\ap^\alpha x|$.
Since $\ap^{-\alpha}$ is continuous,
it follows that the norm $|\cdot|_{\cb_{\al}}$
is equivalent to the graph norm of
$\ap^\alpha$.
If $\alpha=0$, then $\cb_{\al}=\cB$
and $\ap^0=\mbox{Id}$.
If $\alpha<0$, let $\cb_{\al}$ be
the completion of $\cB$ with respect
to $|x|_{\cb_{\al}}=|\ap^\alpha x|$,
which means in particular that $\cb_{\al}$ is a larger space than
$\cB$. We will also set $\cb_{-\infty}=\cup_{\al\in\R}\cb_\al$.

\vspace{0.2cm}

Among the important facts about these spaces,
note the following ones:
For any $\alpha\in\RR$ and any $\rho\geq 0$,
\begin{gather}
\ap^{-\rho}\text{ maps }\cb_{\al}\text{ onto }
\cB_{\alpha+\rho}\text{, for all }\alpha\in\RR,\ \rho\geq 0,\\
|x|_{\cb_{\al}}\leq C_{\alpha,\rho}|x|_{\cb_{\rho}}
\text{ for all }x\in\cb_{\al}
\text{ and all }\alpha\leq \rho.
\end{gather}
Moreover, for all $\alpha,\beta\in\RR$,
\begin{equation}
\label{eq-sg-5}
\ap^\alpha \ap^\beta=\ap^{\alpha+\beta}
\text{ on }\cB_\gamma
\end{equation}
with $\gamma=\max\{\alpha,\beta,\alpha+\beta\}$.
The semi-group $(S_t)_{t\geq 0}$ also satisfies
\begin{gather}
\label{eq-sg-1}
S_t\text{ may be extended to $\cb_{\al}$ for all $\alpha<0$
and all $t> 0$},\\
\label{eq-sg-2}
S_t\text{ maps }\cb_{\al}\text{ to }\cb_{\rho}
\text{ for all }\alpha\in\RR,\ \rho\geq 0,\ t>0,\\
\label{eq-sg-3}
\text{ for all $t>0$, $\alpha\geq0$, }
|\ap^\alpha S_t|
\leq M_\alpha t^{-\alpha}e^{-\lambda t},\\
\label{eq-sg-4}
\text{ for $0<\alpha\leq 1$, $x\in\cb_{\al}$, }
|S_tx-x|\leq C_\alpha t^\alpha |\ap^\alpha x|.
\end{gather}
We will denote with $\LL(\BB,\BB')$ the space of continuous
linear operators from the Banach space $\BB$ to the Banach space
$\BB'$. We let $\LL(\BB) = \LL(\BB,\BB)$. In order to be coherent with our previous notations, we also set $S_{t-s}=S_{ts}$ for a generic semigroup $S$, and $0\le s<t\le T$.

\subsection{Convolutional increments}
\label{sect:complex}

Let us turn now to the main concern of this section, that is the definition
of a complex $(\hca_*,\hd)$ which behaves nicely for
the definition of our evolution problem.

\vspace{0.2cm}

Notice that, due to the fact that the operator $S_{t_1  t_2}$ is
well defined only for $t_1>t_2$, our integration domains will
be of the form $\dn$, where $\dn$ stands for the n-simplex
$$
\cs_n = \{(t_1,\dots,t_n) :  T\ge t_1 \ge t_2 \ge \cdots \ge t_n  \ge 0\}.
$$
Let then $V$ be a separable Banach space.
The basic family of increments we will
work with is $\{ \hcac_n(V);n\ge 0 \}$, where $\hcac_n(V)$ denotes the space
of continuous functions from $\dn$ to $V$. Observe that an operator
$\der:\hcac_n(V)\to\hcac_{n+1}(V)$ can be
defined  just like in (\ref{eq:coboundary}). In particular, if $A\in\hcac_1(V)$
and $B\in\hcac_2(V)$, the relation (\ref{eq:simple_application}) is still
valid. However, let us see now why $\der$ is not adapted to the
resolution of equation  (\ref{eq:eq3}).

\vspace{0.2cm}

What made $\der$ an interesting operator in Section \ref{sec:one-dim} was the simple fact that,
if $F\in\hcac^\infty_1(\R)$, then, for $t,s\in\ott^2$, we have
\begin{equation}\label{deltintg}
[\der F]_{ts}=\int_s^t f_u\, du,
\quad\mbox{ with }\quad f=F'.
\end{equation}
However, if $S_t$ is the semigroup defined at Section \ref{analsmg},
and if we set
$$
\hat F_t=\int_0^t S_{t  u} f_u\, du,
\quad\mbox{ for }\quad
t\ge 0, f\in\hcac^\infty_0(\cb),
$$
then the same kind of relation doesn't hold true for $\hat F$. Indeed, for $s\le t$,
if we define the operator $a_{ts}:\cb\to\cb$ as
\begin{equation}\label{defats}
a_{ts}=S_{t  s}-\id,
\end{equation}
where $\id : \BB \to \BB$ is the identity operator, then
it is easily seen that
$$
[\der \hat F]_{ts}=\hat F_t-\hat F_s
=a_{ts}\hat F_s+\int_s^t S_{t  u} f_u\, du,
$$
and hence, in order to get a similar relation to (\ref{deltintg}) in this new context, one should consider an operator $\hd:\hcac_n(\cb)\to\hcac_{n+1}(\cb)$,
defined by
\begin{equation}\label{defhd}
[ \hd A ]_{t_1\dots t_{n+1}}
=
\lc \der A  \rc_{t_1\dots t_{n+1}}
- a_{t_1t_2}A_{t_2\dots t_{n+1}},
\quad\mbox{ for }\quad
A\in\hcac_n(\cb),\, (t_1\dots t_{n+1})\in\dnp.
\end{equation}
In the remainder of the paper, we will write
$\hd A=\der A-a A$, where we made use of the convention
(\ref{cvpdt}).
As in Section \ref{incr}, one can define, for $n\ge 1$,
$$
\cz\hca_n(\cb)=\hca_n(\cb) \cap \ker(\hd),
\quad\mbox{ and }\quad
\cb\hca_n(\cb)=\hca_n(\cb) \cap \mbox{Im}(\hd)
$$

Then the perturbed operator $\hd$ preserves some important properties
of the original coboundary $\delta$:

\begin{proposition}\label{acyc}
The couple $(\hcac_*,\hd)$ is an acyclic cochain complex: $\cZ \hcac_{n+1} =\cB \hcac_{n}$ for any $n \ge 0$.
\end{proposition}

\begin{proof}
Let us prove first that $\hd$ is a coboundary, i.e. $\hd\hd=0$. Indeed, if $F\in\hcac_n$
according to the fact that
$\der\der=0$ and thanks to the forthcoming Lemma \ref{daf}, we have
\begin{eqnarray*}
\hd\hd F&=&
\lp \der-a \rp \lc \lp \der-a \rp F  \rc
=\der\der F -\der(aF)-a\der F+ aa F\\
&=&-\der a F +a \der F -a\der F+ aa F
=aa F-\der a F.
\end{eqnarray*}
Furthermore, it is readily checked that
$$
(\delta a)_{tus} = a_{tu} a_{us}, \qquad (t,u,s) \in \cs_3,
$$
which gives $\hd\hd F=0$.

\medskip

The fact that
$\mbox{Im}\hd_{|\hcac_n}=\ker\hd_{|\hcac_{n+1}}$ can be proved along the
same lines as for the $(\CC_*,\delta)$ complex~\cite{rough}:
pick $A\in\hcac_{n+1}$ such that
$\hd A=0$, and set $B_{t_1\dots t_n}=A_{t_1\dots t_n s}$, with $s=0$. Then
\begin{eqnarray*}
[\hd B]_{t_1\dots t_{n+1}}&=&
\lc\der B\rc_{t_1\dots t_{n+1}s}
+(-1)^{n+1}A_{t_1\dots t_{n+1}}-a_{t_1t_2}A_{t_2\dots t_{n}s}\\
&=&[\hd A]_{t_1\dots t_{n+1}s}+(-1)^{n+1}A_{t_1\dots t_{n+1}}
=(-1)^{n+1}A_{t_1\dots t_{n+1}}.
\end{eqnarray*}
Thus, setting $C=(-1)^{n+1} B$, we get $\hd C=A$.

\end{proof}

The cochain complex $(\hca_*,\hd)$ will be the structure at the base of all the constructions in this paper.
Let us also mention at this point that, when the meaning is obvious, we will transpose the notations of Section \ref{sec:one-dim} to our
infinite dimensional setting. Furthermore, whenever this doesn't lead to an ambiguous situation, we will write $\hcac_n$ instead of $\hcac_n(\cb)$.



\vspace{0.2cm}

Let us give now a simple and useful extension of Proposition \ref{difrul},
which has already been used in the last proposition:
\begin{lemma}\label{daf}
Let $L\in\hcac_{n-1}(\cb)$ and $M\in\hcac_2(\cl(\cb))$. Then
$$
\der\lp ML \rp=\der M\, L- M\, \der L.
$$
\end{lemma}

\begin{proof}
Let $G_{t_1\cdots t_{n}}=M_{t_1t_2}L_{t_2\cdots t_{n}}$. Then
\begin{eqnarray*}
[\der G]_{t_1\cdots t_{n+1}}&=&
\sum_{i=1}^{n+1} (-1)^{i} G_{t_1\cdots \hat t_i \cdots t_{n+1}}\\
&=& -M_{t_2t_3}L_{t_3\cdots t_{n+1}}
+M_{t_1t_3}L_{t_3\cdots t_{n+1}}
+\sum_{i=3}^{n+1} (-1)^{i} M_{t_1t_2}L_{t_2\cdots \hat t_i \cdots t_{n+1}}\\
&=&[\der M]_{t_1t_2t_3}L_{t_3\cdots t_{n+1}}
+ M_{t_1t_2} \sum_{i=2}^{n+1}(-1)^{i}L_{t_2\cdots \hat t_i \cdots t_{n+1}},
\end{eqnarray*}
which yields our claim.

\end{proof}

%\vspace{0.2cm}

%\noindent

\subsection{Computations in $\hcac_*$}

Here again, like in Section \ref{sec:one-dim}, we will try to move
from a smooth setting to an irregular one. And we will start by giving
the equivalent, in our new setting, of Proposition
\ref{prop:Lambda}, which will require first the introduction of
some analytical structures on the spaces $\hca_n$:

\subsubsection{H\"older type spaces} \label{holtype}
First of all, we have to define some H\"older type
subspaces of $\hcac_{k}$, $k \le 3$ related to the spaces $\cb_\al$,
$\al\in\R$: for $\mu\ge 0$ and $g\in\cac_2(\cb_\al)$, we set
\begin{equation}\label{normba}
\norm{g}_{\mu,\al} \equiv
\sup_{t,s\in\cs_2}\frac{|g_{ts}|_{\cb_\al}}{|t-s|^\mu},
\quad\mbox{and}\quad
\hcac_2^{\mu,\al}=\lcl g\in\hcac_1(\cb_\al);\, \norm{g}_{\mu,\al}<\infty  \rcl,
\end{equation}
and the definition above also induces some semi-norms on $\cac_1$:
for $\gamma > 0, \al\in\R$, we say
that $f \in\hcac_1^{\gamma,\alpha}$ if
$$
\|f\|_{\gamma,\alpha} \equiv \|\hd
f\|_{\gamma,\alpha}< \infty.
$$
Another useful subspace of $\hcac_1$ will be
$\hcac_1^{0,\alpha}$, the space
of bounded paths in $\BB_\alpha$ with the supremum norm $\|f\|_{0,\alpha} =
\sup_{t \in [0,T]} |f_t|_{\alpha}$.

\vspace{0.2cm}

As far as $\hcac_3$ is concerned,
$\hcac_3^{\mu,\al}$ can be defined in the following way: set
\begin{eqnarray}
  \label{eq:normcc3}
  \norm{h}_{\gamma,\rho,\al} &=& \sup_{t,u,s\in\cs_3}
\frac{|h_{tus}|_{\cb_\al}}{|t-u|^\gamma |u-s|^\rho}\\
\|h\|_{\mu,\al} &\equiv &
\inf\left \{\sum_i \|h_i\|_{\rho_i,\mu-\rho_i,\al} ;\, h =
 \sum_i h_i,\, 0 < \rho_i < \mu \right\} ,\nonumber
\end{eqnarray}
where the last infimum is taken over all sequences $\{h_i \}_i$ such that $h
= \sum_i h_i$ and for all choices of the numbers $\rho_i \in (0,z)$.
Then  $\|\cdot\|_{\mu,\al}$ is again  easily seen to be a norm, and we set
$$
\hcac_3^{\mu,\al}=\lcl h\in\hcac_3(\cb_\al);\, \|h\|_{\mu,\al}<\infty \rcl.
$$

\vspace{0.2cm}

Eventually, we will need to introduce a slight extension of the spaces we
have just defined above: for $j=1,2$, let $\ce_j^{\mu,\al}$ be defined by
\begin{equation}\label{defce}
\ce_j^{\mu,\al}=\bigcap_{\ep\le\mu\wedge 1^-}\hcac_j^{\mu-\ep,\al+\ep},
\end{equation}
where $\ep\le\mu\wedge 1^-$ stands for the condition
$\ep\in [0,\mu]\cap [0,1)$, and where
the intersection is considered along any  arbitrary  family
$\{ 0\le \ep_1<\cdots<\ep_n \le \mu\wedge 1^-  \}$ for $n\ge 1$.
Obviously, some families of operators will play an important role
in the sequel, and this will lead us to the following specific
definitions for operator-valued increments:
\begin{definition}
For $\mu\ge 0$ and $\al,\beta\in\R$,
we will call $\hcac^{\mu}_2 \LL^{\beta,\al}$ the space
$\hcac_{1}^{\mu}(\cl(\cb_{\beta};\cb_{\al}))$, and will
denote by $\ce_2^{\mu} \LL^{\beta,\al}$ the space
$$
\ce^{\mu}_2 \LL^{\beta,\al}
=
\bigcap_{\ep\le  \mu\wedge 1^-}\hcac_2^{\mu-\ep} \LL^{\beta,\al+\ep},
$$
where the intersection is still considered along any arbitrary finite family
$\{ 0\le \ep_1<\cdots<\ep_n\le \mu\wedge 1^-  \}$ for $n\ge 1$.
The natural norm on $\hca_2^{\mu}\cl^{\beta,\al}$ will be defined by
\begin{equation}\label{defopnor}
\| A\|_{\mu,\beta,\al}
=
\sup_{t,s\in\cs_2}\frac{\| A_{ts} \nop}{|t-s|^{\mu}},
\end{equation}
and when we consider some Hilbert-Schmidt operators, the corresponding
spaces will be denoted by
$\hcac^{\mu}_2\LL^{\beta,\al}_{\text{HS}}$ and
$\ce^{\mu}_2\LL^{\beta,\al}_{\text{HS}}$.
\end{definition}

\subsubsection{The convolution sewing map and related properties}

Here is a first proposition showing how the analytical structures
introduced above interact with our previous algebraic notations:
\begin{proposition}
\label{prop:uniq}
If $\mu > 1$, then for any $\al\in\R$,
$\cZ \hcac_2^{\mu,\al} = \{ 0 \}$.
\end{proposition}
\begin{proof}
Take $h \in  \cZ \hcac_2^{\mu,\al}$. Then, according to Proposition
\ref{acyc}, there exists $f \in \hcac_1$ such that $h = \hd f$.
Consider the telescopic sum
\begin{equation*}
h_{ts} = (\hd f)_{ts} = \sum_{i=0}^n S_{t  t_{i+1}} (\hd f)_{t_{i+1} t_i},
\end{equation*}
with respect to the partition
$\Pi_{ts}^n = \{ t_{0 \le i \le n+1} : t_0 = s, t_{n+1}
= t \}$ of the interval $[s,t]$.
Since $\hd f \in \cZ \hcac_2^{\mu,\al}$ with $\mu > 1$, we have
\begin{equation*}
| (\hd f)_{ts}|_\alpha \le  \sum_{i=0}^n |(\hd f)_{t_{i+1} t_i}
|_\alpha \le \|\hd f\|_{\mu,\alpha} \sum_{i=0}^n |t_{i+1}-t_i|^\mu
\end{equation*}
which converges to zero as the size of the partition goes to zero. Since
$t,s$ are arbitrary we have
$\hd f = h=0$ in $\hcac_2^{\mu,\alpha}$.
\end{proof}

\vspace{0.2cm}

We can now state and prove the equivalent of Proposition~\ref{prop:Lambda}
in our evolution setting, which is the main aim of this section:
\begin{theorem}\label{th:conv-lambda}
Let $\mu>1$, $\al\in\R$. There exists a unique sewing map $\hat \Lambda :\cZ
\hcac_3^{\mu,\al} \to  \ce_2^{\mu,\al}$ such that $\hd \hat \Lambda =
\id_{\cZ \hcac_3}$.
Furthermore, for any $0\le \ep\le \mu\wedge 1^-$, there exists
a strictly positive constant $c_{\mu,\ep}$ such that
\begin{equation}\label{ineqhla}
\|\hat \Lambda h\|_{\mu-\eps,\alpha+\eps} \le c_{\mu,\ep} \|h\|_{\mu,\al},
\end{equation}
for any $h \in \cZ \hcac_3^{\mu,\al}$.
\end{theorem}

\begin{proof}
Like in the proof of Proposition~\ref{prop:Lambda}, we will divide our computations in two steps:

\vspace{0.2cm}

\noindent
{\it Step 1:}
The  uniqueness part of our theorem simply stems from the
fact that if we have $\hd a = h$ and $\hd a'= h$ with
$a,a' \in  \hcac_1^{\mu,\al}$, then $b  = a-a' \in \cZ
\hcac_1^{\mu,\al}$ and since $\mu > 1$, by Prop.~\ref{prop:uniq} we
must have $b = 0$.

\vspace{0.2cm}

\noindent
{\it Step 2:}
The existence part can be adapted from Proposition \ref{prop:Lambda},
and we will construct a process $M\in\ce_2^{\mu,\al}$ such that $\hd M=h$
starting from any $B\in\hcac_2(\cb_\al)$ satisfying $\hd B=h$
(this increment $B$ exists thanks to Lemma \ref{acyc}).
Now, similarly to (\ref{defmtsn}), we will set, for
a given $n\ge 1$, and $(t,s)\in\cs_2$,
$$
M_{ts}^{n}=B_{ts}-\sum_{i=0}^{2^n-1} S_{t  \rn_{i+1}} B_{\rn_{i+1},\rn_{i}},
$$
where $s,t$ and $\rn_i$ have been defined at (\ref{defpartin}). Then
$M_{ts}^{0}=0$ and
\begin{align*}
&M_{ts}^{n+1}-M_{ts}^{n}\\
&=
\sum_{i=0}^{2^n-1}
\lp S_{t  \rnp_{2i+2}} B_{\rnp_{2i+2},\rnp_{2i}}
-S_{t  \rnp_{2i+2}} B_{\rnp_{2i+2},\rnp_{2i+1}}
-S_{t  \rnp_{2i+1}} B_{\rnp_{2i+1},\rnp_{2i}}\rp\\
&=
\sum_{i=0}^{2^n-1}
S_{t  \rnp_{2i+2}}
\lp B_{\rnp_{2i+2},\rnp_{2i}}
- B_{\rnp_{2i+2},\rnp_{2i+1}}
- B_{\rnp_{2i+1},\rnp_{2i}}  \rp\\
&
\hspace{6cm}
- S_{t  \rnp_{2i+2}}
\lc S_{\rnp_{2i+2}  \rnp_{2i+1}}-\id \rc B_{\rnp_{2i+1},\rnp_{2i}}.
\end{align*}
Thus, according to the definition (\ref{defhd}) of $\hd$, we get
\begin{eqnarray*}
M_{ts}^{n+1}-M_{ts}^{n}
&=& \sum_{i=0}^{2^n-1} S_{t  \rnp_{2i+2}} \lc
\lp \der B \rp_{\rnp_{2i+2},\rnp_{2i+1},\rnp_{2i}}
- a_{\rnp_{2i+1},\rnp_{2i+1}} B_{\rnp_{2i+2},\rnp_{2i}}\rc\\
&=&\sum_{i=0}^{2^n-1} S_{t  \rnp_{2i+2}}
\lp \hd B \rp_{\rnp_{2i+2},\rnp_{2i+1},\rnp_{2i}}
=\sum_{i=0}^{2^n-1} S_{t  \rnp_{2i+2}}
h_{\rnp_{2i+2},\rnp_{2i+1},\rnp_{2i}}.
\end{eqnarray*}
Hence, for any $\ep<\mu$, we get, invoking (\ref{eq-sg-3}),
\begin{eqnarray*}
\lln A^{\al+\ep}\lp M_{ts}^{n+1}-M_{ts}^{n} \rp \rrn
&\le&
c_{\ep} \sum_{i=0}^{2^n-1} |t-\rn_{i+1}|^{-\ep} |h|_{\mu,\al} |t-s|^{\mu}\\
&\le& \frac{c_{\ep}|t-s|^{\mu-\ep}|h|_{\mu,\al}}{2^{n(\mu-1)}}
\int_0^1 u^{-\ep} du,
\end{eqnarray*}
which gives, like in Proposition \ref{prop:Lambda}, that
$M_{ts}\equiv\lim_{n\to\infty}M_{ts}^{n}$ exists, and
is an element of $\ce_2^{\mu,\al}$.
Now, the fact that $\hd M=h$ can be shown analogously to the case
of Proposition \ref{prop:Lambda}, and the proof of
(\ref{ineqhla}) is straightforward.

\end{proof}

\noindent
A direct consequence of the existence of the $\hat \Lambda$-map is a
result of convergence of finite sums:
\begin{corollary}
\label{cor:conv-integration}
Let $g \in \hcac_2$ such that $\hd g \in \hcac_3^{\mu,\alpha}$ for some $\mu > 1$. Then the 1-increment
$\hd f = (\id-\hat \Lambda \hd) g
\in \hcac_2^\al$ satisfies
$$
(\hd f)_{ts} = \lim_{|\Pi_{ts}| \to 0}
\sum_{i=0}^n S_{t  t_{i+1}} g_{t_{i+1}\,t_{i}},
$$
for all $(t,s)\in\cs_2$.
\end{corollary}
\begin{proof}
It follows the lines of the proof of Corollary~\ref{cor:integration}.
\end{proof}


\vspace{0.2cm}

We will now define an equivalent of the iterated integrals of Section
\ref{cpss} in our convolution context: consider some smooth functions
$g\in\hcac_1^\infty(\cl(\cba))$ and $f\in\hcac_1^\infty(\cba)$, for
some $\al\in\R$. Then $\cj(dg\, f)$ will be defined as an element of
$\hcac_3^\infty(\cba)$ by
$
\cj_{ts}\lp dg\, f \rp=\int_s^t dg_v\, f_v$,
 for  $(t,s) \in \cs_2$.
We will also need some integrals of processes weighted by the semi-group $S$, defined as follows, for $0\le s < t \le T$:
\begin{equation*}
\cj_{ts}\lp\hdd g\, f \rp=\int_s^t S_{t  v} dg_v\, f_v.
\end{equation*}
Once these elementary blocks have been defined, the iterated integrals
\begin{equation}\label{eq:not-iter-convol-intg}
\cj\lp d^{*_n}f_n\dots d^{*_1}f_1 \rp
\quad\mbox{ for }\quad
f_n,\dots,f_2\in  \hcac_1^\infty(\cl(\cba)),
f_1\in \hcac_1^\infty(\cba),
\end{equation}
where $d^{*_j}f_j$ stands for any of the increments of the form $df_j$ or
$\hdd f_j$, can be defined recursively
along the same lines as in Section \ref{cpss}.
In particular, the operator-valued increment $\cj(\hdd g S)$ is defined by
$$
\cj_{ts}(\hdd g S)=\ist S_{t  u} dg_u S_{u  s}.
$$
The relations between
$\hd$ and these integrals, which will be useful for our purposes, can
be summarized in the following:
\begin{proposition}\label{dishd}
Let $\al\in\R$, and $g\in\hcac_1^\infty(\cl(\cba))$,
$f\in\hcac_1^\infty(\cba)$. Then
$$
\hd\lp \cj(\hdd f) \rp=0,
\qquad
\hd\lp \cj(\hdd g\, f) \rp=0,
\qquad
\hd\lp \cj(\hdd g \hd f) \rp=
\cj(\hdd g S)\, \hd(f),
$$
and
$$
\hd\lp \cj(\hdd g \hdd f)\rp=
\cj(\hdd g S) \cj(\hdd f),
\qquad
\hd\lp \cj(\hdd g df)\rp=
\cj(\hdd g) \cj(df).
$$
\end{proposition}

\begin{proof}
The proof of these results is elementary. We will give some details about
the last relation for sake of completeness. For any $(t,u,s) \in \cs_3$, invoking the definition of $\hd$, we have
\begin{align*}
&\lc \hd\lp \cj(\hdd g df)\rp \rc_{tus}
=\lc \der\lp \cj(\hdd g df)\rp \rc_{tus}
- a_{tu} \cj_{us}(\hdd g df)\\
&=\int_s^t S_{t  v} dg_v\int_s^v df_w
-\int_u^t S_{t  v} dg_v\int_u^v df_w
-\int_s^u S_{u  v} dg_v\int_s^v df_w
%&\hspace{6cm}
-\lc S_{t  u}-\id  \rc \int_s^u S_{u  v} dg_v\int_s^v df_w\\
&=\cj_{tu}(\hdd g) \cj_{us}(df),
\end{align*}

which proves the claim.


\end{proof}

\subsection{Fractional heat equation setting}\label{sec:circle}

In this section, we will give the general setting under which
we  will try to define and solve the stochastic heat equation
driven by an infinite dimensional fractional Brownian motion:
as mentioned in the introduction, the main application we have
in  mind is the situation where
$A=\Delta-\id$, and $\Delta$ is the Laplace operator on the circle $S$,
assimilated to $\ou$.
This operator can be diagonalized in the trigonometric  basis of
$L^2([0,1];\C)$, namely $\{e_n;n\in\Z  \}$, where
$$
e_n(x)= e^{2\imath \pi n x},
\quad
x\in\ou.
$$
Associated to these eigenfunctions are the eigenvalues
$\la_n=-1- (2 \pi n)^2$. We have chosen to deal with
$A=\Delta-\id$ instead of $\Delta$ itself for computational convenience,
since this choice avoids the problem of a null eigenvalue for
constant functions.
Notice that in this case, $A$ is the generator of an analytical
semigroup, and all the constructions
of Section \ref{analsmg} goes through. Then
$\cb_\al$ can be identified with $H_\al$,
the usual Sobolev space
based on $L^2([0,1])$, for the definition of which we refer to
Adams \cite{Ad}, and
$\{ S_t;t\ge 0 \}$ stands for the heat semi-group, which admits
a kernel $G_t(\xi,\eta)$ for $t>0$ and $x,y\in\ou$.
In this context, set $G_t^{\al}(\xi,\eta)$
for the kernel of the operator $\ap^{\al}S_t$, and
$G^{\beta}(\xi,\eta)$ for the kernel of the operator $\ap^{-\beta}$.
Then, for $\al\in\R$ and $\beta>0$, $G_t^{\al}$ and $G^{\beta}$
admit the following spectral decomposition:
\begin{equation}\label{dcp:ker}
G_t^{\al}(\xi,\eta)
=\sum_{n\in\Z}\la_n^{\al} e^{-t\la_n} e_n(\xi) \bar e_n(\eta),
\quad\mbox{ and }\quad
G^{\beta}(\xi,\eta)
=\sum_{n\in\Z}\la_n^{-\beta}  e_n(\xi) \bar e_n(\eta).
\end{equation}



\vspace{0.2cm}

Let us specify now the noise $X$ we will consider: we will try to
stick to the existing literature on the topic, and
choose a fractional Brownian noise in time, defined on a certain
complete probability space $(\Omega,\cf,P)$,
which will be homogeneous in space,
with a spatial covariance function $Q$. Namely, $X$ will be a centered Gaussian field indexed by functions on $\ott\times \ou$, such that if
$\phi$ and $\psi$ are smooth enough, then
\begin{equation}\label{covstr}
E\lc X(\phi) X(\psi) \rc
=
c_H \int_{\ott^2}
\lp \int_{\ou^2}Q(\xi-\eta)\phi(u,\xi)\psi(v,\eta) d\xi d\eta\rp
|u-v|^{2H-2} du dv,
\end{equation}
with $c_H=H(2H-1)$, for $H>\frac12$. Notice that, in order to simplify our
statements, we will generally assume that $Q$ can be decomposed itself
on the basis $\{ e_n;n\in\Z \}$ in the following way:
\begin{equation}\label{dcp:cov}
Q(\xi)=\sum_{n\in\Z}q_n e_n(\xi),
\quad\mbox{ with }\quad
q_n=\la_n^{-\nu},
\mbox{ for }\nu\in[0,1),
\end{equation}
and notice that the case $\nu=0$ corresponds to the white noise in space,
while the case $\nu>1/2$ corresponds to a noise admitting a density
in space.
Some explicit construction of such kind of noise, as well as an account
on the related stochastic calculus,  can be found in
\cite{TTV}. The methodology we will develop in the rough case will
also enable us to handle the Brownian motion case, which means
a covariance structure given by
\begin{equation}\label{cov:bm}
E\lc X(\phi) X(\psi) \rc
=
\int_{\ott}
\lp \int_{\ou^2} \phi(u,\xi) Q(\xi-\eta) \psi(u,\eta) d\xi d\eta   \rp
du.
\end{equation}
% Recall also that, with these notations in mind, we try to solve
% the heat equation under the form (\ref{mulf}), which can be rewritten
% under the form (\ref{eq:eq3}) as well, as we mentioned in the introduction.

\vspace{0.2cm}


We give here a slight extension of a result result of~\cite{rough},
which will be used below to
prove existence of regular versions of some stochastic processes,
following the well-known approach of Garsia-Rodemich-Rumsey.
In order to state this extension, we shall introduce for the first time a 
variant of the operator $\hd$, called $\td$, acting on operator-valued increments which turns out to be useful in
the sequel, and which is defined by:
\beq\label{eq:def-tilde-delta}
\td Q = \hd Q - Q a= \der Q - a Q -Q a, \quad Q \in \cac^*(\LL(\BB)).
\eeq 
With this additional notation in hand, our regularity lemma is the following:
\begin{lemma}
\label{lemma:besov}
For any $\gamma > 0$,  $\alpha,\beta \in \RR$ and $p \ge 1$,  there exists a constant
$C$ such that for any $R \in \hcac_{2}(\cl^{\beta,\alpha})$, we have
\begin{equation}
\label{eq:generalboundxx}
\|R\|_{\gamma,\beta,\alpha} \le C (U_{\gamma+2/p,p,\beta,\alpha}(R)
+\|\td R \|_{\gamma,\beta,\alpha}),
\end{equation}
where
\begin{equation*}
 U_{\gamma,p,\beta,\alpha}(R) = \left[ \int_{\cs_2}
 \left(\frac{|R_{ts}|_{\beta,\alpha}}{|t-s|^\gamma}\right)^p dt ds \right]^{1/p}.
\end{equation*}
\end{lemma}

\begin{proof}
The proof of this result is similar to that appearing in~\cite{rough}. The crucial step in modifying that proof is to note that,
for any triple $(t,u,s) \in \cs_3$ we have, writing $S_{us}=S_{u  s}$ for $u\ge s$,
\begin{equation*}
  \begin{split}
R_{ts} & = R_{tu} + R_{us}  + (\der R)_{tus}
% \\ & = R_{tu} S_{us} + S_{tu} R_{us} + (\der R)_{tus} - R_{tu} a_{us} - a_{tu} R_{us}
 %\\ & 
 = R_{tu} S_{us} + S_{tu} R_{us} + (\td R)_{tus} 
\end{split}
\end{equation*}
which implies that
\begin{equation*}
\begin{split}
|R_{ts}|_{\beta,\alpha}& \le |R_{tu} S_{us}|_{\beta,\alpha} + |S_{tu} R_{us}|_{\beta,\alpha} + |(\td R)_{tus}|_{\beta,\alpha}
\\ & \le M |R_{tu}|_{\beta,\alpha} + M |R_{us}|_{\beta,\alpha} + |(\td R)_{tus}|_{\beta,\alpha}.
\end{split}
\end{equation*}
where we used the fact that $|S_\tau|_{\alpha,\alpha} \le M$ for any $\alpha$ and any $\tau \ge 0$.
\end{proof}


%%%%%%%%%%%%%%%%%%%%%%%%%%%%%%%%%%%%%%%
%%%%%%%%%%%%%%%%%%%%%%%%%%%%%%%%%%%%%%%

\section{Young theory}

We are now ready to analyze the Young integration in the evolution setting
along the same lines as in Section \ref{sec:dissection}: we will first
define the integral $\cj(\hdd x z)$ for two Young paths $x,z$ in an
abstract setting. Then we will solve Young SPDEs, and eventually,
check our main assumptions in the fractional heat equation setting of
Section~\ref{sec:circle}.


%%%%%%%%%%%%%%%%%%%%%%%%%%%%%%%%%%%%%%%%%
%%%%%%%%%%%%%%%%%%%%%%%%%%%%%%%%%%%%%%%%%

\subsection{Young integration}\label{young:integration}

The extension of the notion of integral weighted by an analytical
semi-group will be performed through the following algorithm,
which will be used in fact throughout the remainder of the paper:
\begin{enumerate}
\item
Assume first that $x$ is a regular operator-valued increment, and
$z$ a regular \hbox{$\cb$-valued} function and let
$
\cj_{ts}( \hdd x z ) \equiv
\ist S_{t  u} dx_u\, z_u$,
 for $
(t,s)\in\cs_2
$,
as an element of $\hca_2$.
\item
Through the application of $\hd$ and $\hl$, try to get an expression for
$\cj(\hd x z)$ which depends only on minimal regularity requirements for $x$ and $z$.
\item
Extend the notion of integral using the previous step, and see that
it induces the convergence of some well-chosen Riemann sums.
\end{enumerate}

\vspace{0.2cm}

Here is how this general strategy can be implemented here:
suppose for the moment that $x$ is a smooth operator-valued function and
$z$ a smooth function. Then it is easily checked that
\begin{equation}
  \label{eq:split-1}
\cj(\hdd x\,  z)
=
\cj(\hdd x S)  z+ \cj(\hdd x\, \hd z).
\end{equation}
Note that
in this last equality appears for the first time an incremental
operator which will play a fundamental role in the sequel, namely
the operator $X^1 \in \hcac_2(\LL(\BB))$ defined by
\begin{equation}\label{defnts}
X^1_{ts}=
\cj_{ts}(\hdd x S) = \int_s^t S_{t  u} dx_u S_{u  s}.
\end{equation}
And here is an important point of our strategy: the
noise $x$ does not appear by itself but always inside a convolution
of the form (\ref{defnts}),
so its action is {\it milded} by the regularizing properties of the
semigroup.

\smallskip

Applying  $\hd$ to the last term of eq.~(\ref{eq:split-1}) and
invoking Proposition \ref{dishd}, we get
\begin{equation*}
%\label{dishdx}
\hd [ \cj(\hdd x\, \hd z)]
%= - \hd[ \cj(\hdd x)  z ]
=
\cj(\hdd x S) \hd z = X^1 \hd z.
\end{equation*}
If the 2-increment $X^1 \hd z$ is small enough,
namely if
$X^1 \hd z \in \hcac_2^{\mu,\delta}$ for some $\delta$ and
some $\mu > 1$, then we can express $\cj(\hdd x\,  z)$ as
\begin{equation}\label{yalgobis}
\cj(\hdd x\,  z)= X^1   z+
\hl [X^1 \hd z ]  = (\id-\hat \Lambda
\hd) [X^1  z].
\end{equation}
The last equality is justified by noting that  when $x$ is a smooth incremental operator, we have $\hd X^1=X^1 a$ (i.e. $\td X^1 = 0$), and thus by Lemma \ref{daf} one obtains
that $\hd(X^1 z)=-X^1 \hd z$.

\vspace{0.2cm}

Let us turn now to the second point of our general strategy, which consists
in inverting the process which lead to (\ref{yalgobis}):
indeed, if we can define
properly the right hand side of (\ref{yalgobis}), then we will be able
to extend the notion of integral by a procedure which is coherent with the
basic properties required to any integral $\cj$.
Notice that this step only relies on the definition of
an operator $X^1$ associated to $x$, satisfying $\td X^1=0$, and such that
$X^1$ is regular enough. This will be formalized in the following Theorem (recall that the space $\cb_{-\infty}$ has been defined at
Section \ref{analsmg}.
):
\begin{theorem}
\label{th:young}
Let then $x$ be a path from $\ott$ to $\cb_{-\infty}$ such that the operator
$X^1$ associated to $x$ is well defined as an element of
$\hcac^{\kappa}_2 \LL^{\beta,\al}$,
where $\beta,\ka$ are positive constants, and $\al\in\R$.
We also assume that $X^1$  satisfies the algebraic relation $\hd X^1=X^1 a$.
Let  $z\in\hcac_{1}^{\eta,\beta}$,
with $\kappa+\eta>1$, and set
\begin{equation}
\label{eq:j-def-young}
\cj(\hdd x z)=X^1\,  z+\hl [X^1\, \hd z ]=( \id-\hl\hd  ) [X^1\, z].
\end{equation}
Then
\begin{enumerate}
\item
$\cj(\hdd x z)$ is well defined as an element of $\ce_2^{\kappa,\al}$.
\item
For a constant $c>0$,
we have
$$
\| \cj(\hdd x z) \|_{\kappa,\alpha}
\le
 \|X^1\|_{\kappa,\beta,\al} (\|z\|_{0,\beta} + c_\mu \|z\|_{\eta,\beta}),
$$
where
the norm $\|\cdot\|_{\kappa,\beta,\al}$ has been defined at relation
(\ref{defopnor}).
\item
It holds that, for any $0\le s<t\le T$
$$
\cj_{ts}(\hdd x z)
=\lim_{|\Pi_{ts}|\to 0}\sum_{i=0}^n S_{t-  t_{i+1}} X^1_{t_{i+1},t_{i}} z_{t_{i}},
$$
where the limit is over all partitions $\Pi_{ts} = \{t_0=t,\dots,
t_n=s\}$ of $[s,t]$ as the mesh of the partition goes to zero.
\end{enumerate}
\end{theorem}

\begin{proof}
Since $X^1 z$ is a well defined element of
$\hcac_2^{\kappa,\alpha}$, in order to show that the r.h.s. of eq.~(\ref{eq:j-def-young})
is well defined, it only remains to  check that $X^1 \hd z$ is in the
domain of $\hl$. However, since we have assumed that
$X^1\in\hcac^{\kappa}_2 \LL^{\beta,\al}$ and $\hd z \in \hca_2^{\eta,\beta}$,
we obviously get that $X^1 \hd z\in\hca_3^{\kappa+\eta,\al}$.
Thus, according to Theorem~\ref{th:conv-lambda},
$X^1 \hd z\in\mbox{Dom}(\hl)$, yielding the first assertion of our
theorem.

Moreover, thanks to the second part of Theorem~\ref{th:conv-lambda},
we have
\begin{equation}\label{estlai}
\| \hl[X^1\, \hd z] \|_{\mu,\al}
\le
c_{\mu} \|  X^1 \|_{\kappa,\beta,\al}
\| z \|_{\eta,\beta},
\end{equation}
and it is also readily checked that
\begin{equation}\label{estidx}
\| X^1\, z \|_{\kappa,\al}
\le
\|  X^1 \|_{\kappa,\beta,\al}
\| z \|_{0,\beta},
\end{equation}
which shows our second claim, by using eq.~(\ref{estlai}) and eq.~(\ref{estidx})
to estimate the r.h.s. of~(\ref{eq:j-def-young}).
Eventually, the third part of the theorem is a direct consequence of
Corollary~\ref{cor:conv-integration}.

\end{proof}

\begin{remark}
It is worth stressing at this point some elementary properties enjoyed
by the extension of the notion
of integral given by Theorem \ref{th:young}:
\begin{itemize}
\item
The third part of the theorem states that $\cj_{ts}(\hdd x z)$ is
associated to some natural Riemann sums involving the processes
$x$ (through $X^1$) and $z$.
\item
The arguments leading to relation (\ref{yalgobis}) also show that,
in case of some smooth processes $x$ and $z$, our integral $\cj_{ts}(\hdd x z)$
coincides with the usual one.
\end{itemize}
These first properties seem to imply that  our integral extension
is a reasonable one.
\end{remark}




%%%%%%%%%%%%%%%%%%%%%%%%%%%%%%%%%%%%%%%%%%%%%
%%%%%%%%%%%%%%%%%%%%%%%%%%%%%%%%%%%%%%%%%%%%%



\subsection{Young SPDEs}
\label{sec:young}

Recall that we wish to solve an equation of the form
\begin{equation}
\label{eq:evol-young}
dy_t = A y_t dt + dx_t f(y_t),
\qquad t\in\ott,
\end{equation}
with an initial condition $y_0=\psi\in\cb_{\kappa}$,
where $x$ is an operator-valued process which represents our
noise and $f : \BB \to \BB$ is a (possibly) non-linear regular
map. As mentioned in the introduction, we will consider equation
(\ref{eq:evol-young}) in the mild sense, i.e. we will say that
$y$ is a solution to (\ref{eq:evol-young}) if, for a given $\kappa > 0$ (specified below) we have $y\in\cac_1^{\kappa,\kappa}$ and if, for any $t\in\ott$, $y_t$ satisfies
\begin{equation}\label{miljto}
y_t = S_t \psi +
\int_0^t S_{t  u} dx_u f(y_u) = S_t y_0 + \cj_{t0} (\hat dx f(y)),
\end{equation}
where the integral $\cj_{t0} (\hat dx f(y))$ is understood in the sense
of Theorem \ref{th:young}. In fact, we will focus here on a slight extension
of the problem given by (\ref{miljto}): we will search for a (unique)
process $y\in\cac_1^{\kappa,\kappa}$ satisfying, for any $(t,s)\in\cs_2$,
\begin{equation}\label{extmiljto}
y_t =
S_{t  s} y_s + \cj_{ts} (\hat dx f(y)),
\qquad y_0=\psi,
\end{equation}
from which one recovers obviously (\ref{miljto}) by taking $s=0$.
Now, (\ref{extmiljto}) can be expressed in terms
of convolution increments, since it it equivalent to the following one:
\begin{equation}\label{algmil}
[\hd y]_{ts} = \cj_{ts}(\hdd x f(y))
=\lc(\id-\hat \Lambda \hd) [X^1 f(y)]\rc_{ts},
\quad\mbox{ for } (t,s)\in\cs_2,\quad
\mbox{ and }
\qquad y_0 = \psi,
\end{equation}
which sticks better to the algebraic formalism introduced in the previous
sections.

\vspace{0.2cm}

Let us specify also some of the assumptions under which our computations
will be performed: first of all, the incremental operator $X^1$ defined
by (\ref{defnts}) will be assumed to be in the following class:
\begin{hypothesis}\label{hyp:n}
Assume that
$
X^1 \in\hcac_2^{\tilde\gamma}\cl^{0,-\kappa} \cap \hcac_2^{\kappa_0}\cl^{\kappa,\kappa} 
$
for some $\tilde \gamma > \kappa_0 >\kappa > 1/4$ such that
$$
\tilde \gamma + \kappa > 1, \qquad \tilde \gamma - \kappa \ge \kappa_0, \qquad \kappa < 1/2.
$$
\end{hypothesis}
Notice that in the hypothesis above, the condition $\kappa<1/2$ is somehow redundant.
Indeed, if $\tilde \gamma \ge \kappa + \kappa_0\ge 2\kappa$, this forces the relation
$\kappa < 1/2$.

\smallskip

As far as the function $f$ is concerned,
we will also assume that the following holds true:
\begin{hypothesis}\label{hyp:f}
Let $\kappa$ be the strictly positive constant defined at
Hypothesis \ref{hyp:n}.
We assume that the function $f:\cb_\kappa\to\cb_\kappa$ is locally
Lipschitz, and satisfies $|f(x)|_{\cb_\kappa}\le c_f (1+|x|_{\cb_\kappa})$.
Furthermore, we suppose that $f$ can also be seen as a map from $\cb$
to $\cb$, and when considered as such, it holds that $f$ is
globally Lipschitz.
\end{hypothesis}
%\begin{remark}
%  We have imposed this set of assumptions, which is not
%  completely standard, so that it can be applied, in the case of the
%  heat equation in $\ou$, to a pointwise function
%  $[f(y)](\xi)=\si(y(\xi))$.
%\end{remark}
\vspace{0.2cm}

With these assumptions and notations in mind, we are now
able to solve our evolution equation in the Young sense:
\begin{theorem}
\label{th:young-spde}
Assume Hypothesis~\ref{hyp:n} and \ref{hyp:f} hold true, and that
$\psi\in\cb_{\ka}$.
Let $\hca_{1}^{*,\ka}$ be the subspace of $\hca_1$ defined by the norm
\begin{equation}\label{def:star-ka}
\| z\|_{*,\ka}
=
\|z\|_{0,\ka}+ \|\hd z\|_{\ka,\ka}.
\end{equation}
Then there exists a unique global solution to
(\ref{algmil}) in $\hca_1^{*,\kappa}$. Furthermore,
this solution enjoys the following properties:
\begin{itemize}
\item[(a)]
For any $t\in\ott$, $y_t$ can be written as
$y_t = S_t \psi + (\hd y)_{t0}$.
\item[(b)] Let us call $\Phi$ the map
$(\psi,X^1) \mapsto \Phi(\psi,X^1) = y$, where $y$ is the solution
to (\ref{algmil}). Then $\Phi$ is Lipschitz continuous from
$\BB_{\kappa_0} \times (\hca_2^{\tilde\gamma }\LL^{0,-\rho}
\cap \hca_2^{\ko}\cl^{\kappa,\kappa})$ to $\hca_1^{*,\kappa}$.
\end{itemize}
\end{theorem}

\begin{proof}

A classical fixed point argument will be sufficient to obtain the global solution.
Let us introduce the map
$\Gamma:\hca_1^{*,\kappa}\to\hca_1^{*,\kappa}$
defined in the following way: if $y\in\hca_1^{*,\kappa}$, we set
$\Gamma(y)=z$, where $z$ satisfies
\begin{equation}\label{def:hdz}
[\hd z]_{ts} = \cj_{ts}(\hdd x f(y))
=X^1 f(y) + \lc \hat \Lambda  [X^1 \hd f(y)]\rc_{ts},
\quad\mbox{ for } (t,s)\in\cs_2,\quad
\mbox{ and }
\quad z_0 = \psi.
\end{equation}
Let also $B$ be the ball defined by
\begin{equation}\label{def:ball}
B=
\lcl
y; y_0=\psi, \, \|y\|_{*,\ka}\le 2(1+ |\psi|_{\cb_\ka})
\rcl.
\end{equation}
Then the fixed point argument can be decomposed into two usual steps:
\begin{enumerate}
\item
Show that, on a small enough interval $\ott$, the ball $B$ is left
invariant by $\Gamma$.
\item
Prove that $\Gamma$, restricted to the ball $B$, is a contraction.
\end{enumerate}
We will mainly focus, in this proof, on the first of these steps,
since it contains most of the technical difficulties associated
to our method.



\vspace{0.2cm}

Take $y\in B$, and let us show that $z=\Gamma(y)\in B$
whenever $T$ is small enough
(recall that $\cs_2$ depends on the parameter $T$). To this purpose,
we will first bound the term $\hat \Lambda[X^1 \hd f(y)]$ in (\ref{def:hdz}).
Recall that
\begin{equation}\label{decfy}
\big[ \hd f(y) \big]_{ts}= \lc \der f(y) \rc_{ts} + a_{ts}f(y_s),
\end{equation}
and let us estimate the terms in the right hand side of (\ref{decfy})
separately: on one hand, recalling the notations of Section \ref{holtype},
and thanks to the fact that $f$ is Lipschitz on $\cb$, we have
\begin{eqnarray}\label{bndd:fy}
\lln\lc \der f(y) \rc_{ts}\rrn_{\cb}
&\le&
c_f \lln\lc \der y \rc_{ts}\rrn_{\cb}
\le
c_f \lp \big| [ \hd y ]_{ts}\big|_{\cb}
+  \lln a_{ts} y_s \rrn_{\cb} \rp \nonumber \\
&\le&
c_f \lc \|y\|_{\kappa,0}+ \|y\|_{0,\kappa} \rc |t-s|^{\kappa}
\nonumber \\
&\le&
c_f  \|y\|_{*,\kappa} |t-s|^{\kappa},
\end{eqnarray}
where $c_f$ is a positive constant which may change from line to line,
but which depends only on $f$. On the other hand, according to
Hypothesis~\ref{hyp:f}, it is readily checked that
$|f(y_s)|_{\cb_{\ka}}\le c_f (1+\|y\|_{0,\kappa})$. Thus, invoking
(\ref{eq-sg-4}), we obtain that
\begin{equation}\label{bnd:ay}
\lln a_{ts}f(y_s) \rrn_{\cb}
\le c_f \lp 1+\|y\|_{*,\kappa} \rp |t-s|^{\kappa},
\end{equation}
and plugging (\ref{bnd:ay}) and (\ref{bndd:fy}) into (\ref{decfy}),
we get
\begin{equation}\label{eq:bnd-fy2}
\big| [\hd f(y) ]_{ts} \big|_{\cb}
\le
c_f \lp 1+\|y\|_{*,\kappa} \rp |t-s|^{\kappa}.
\end{equation}
However, we know that $X^1\in\hca_2^{\tilde \ga}\cl^{0,-\kappa}$, and this fact,
together with the last estimate, yields
$$
\big| [X^1 \hd f(y) ]_{tus} \big|_{\cb_{-\kappa}}
\le
c_f  \| X^1 \|_{\tilde \ga,0,-\kappa} \lp 1+\|y\|_{*,\kappa} \rp
|t-u|^{\tilde \ga} |u-s|^{\ka}.
$$
Furthermore by Hypothesis \ref{hyp:n} we have
$\tilde \ga+\ka>1$. This means that Theorem \ref{th:conv-lambda} can
be applied here to obtain that
$\hl(X^1 \hd f(y))\in \ce_{2}^{\tilde \ga+\ka,-\ka}$. In particular,
invoking the definition (\ref{defce}) of the space $\ce_{2}^{\tilde \ga+\ka,-\kappa}$,
and since $2\ka<1$ and $\ko< \tilde\ga-\ka$, we get
$\hl(n\hd f)\in \hca_{2}^{\tilde \gamma-\ka,\ka}\subseteq \hca_{2}^{\ko,\ka}$. Moreover,
\begin{equation}\label{bndhl}
\| \hl(n\hd f)  \|_{\ko,\ka}
\le
c_{f,\tilde \ga,\ka}
\| n \|_{\tilde \ga,0,-\ka} \lp 1+\|y\|_{*,\kappa} \rp .
\end{equation}

\vspace{0.2cm}

A bound  similar to eq.~(\ref{bndhl}) can be found for the term $X^1 f(y)$
appearing in the definition of $\hd z$ in eq.~(\ref{def:hdz}). Indeed, owing to the fact that
$X^1\in\hca_2^{\ko}\cl^{\ka,\ka}$ and that $f$ has linear growth in
$\cb_\ka$, we get
\begin{equation}\label{bnd:ntsf}
 \lln X^1_{ts}f(y_s) \rrn_{\cb_\ka}
\le
c_f \| X^1 \|_{\ko,\ka,\ka} \lp 1+|y_s|_{\cb_\ka} \rp |t-s|^{\ko}
\le
c_f \| X^1 \|_{\ko,\ka,\ka} \lp 1+\|y\|_{*,\ka} \rp |t-s|^{\ko}.
\end{equation}
Hence, plugging (\ref{bnd:ntsf}) and (\ref{bndhl}) into (\ref{def:hdz}),
one obtains that
$
\|\hd z\|_{\ko,\ka}
\le
c_{f,X^1,\tilde \ga,\ka} \lp 1+\|y\|_{*,\ka} \rp
$.
Note here a crucial point: starting from $y \in \hcac_1^{\ka,\ka}$ we have constructed $z \in \hcac_1^{\ko,\ka}$ with $\ep = \ko - \ka> 0$. This little regularity gain can be used in order to write
\begin{equation}\label{bnd:hdz}
\|\hd z\|_{\ka,\ka}
\le
c_{f,X^1,\tilde \ga,\ka} \lp 1+\|y\|_{*,\ka} \rp T^{\ep},
\end{equation}
Now, the quantity $T^{\ep}$ can be made arbitrarily small as $T\to 0$.
Moreover recall that we still need a bound on $\| z\|_{*,\ka}$ defined
by (\ref{def:star-ka}), and thus an estimate on $\| z\|_{0,\ka}$
is needed at this point. However, it is easily checked that
\begin{equation}\label{bnd:zka}
|z_t|_{\cb_\ka}
\le
|S_t\psi|_{\cb_\ka} + |(\hd z)_{t0}|_{\cb_\ka}
\le
|\psi|_{\cb_\ka} + T^{\ko} \|\hd z\|_{\ko,\ka}.
\end{equation}
Putting together (\ref{bnd:hdz}) and (\ref{bnd:zka}), we finally get,
on $\ott$, that
$$
\| z\|_{*,\ka}
\le
|\psi|_{\cb_\ka} + c \lp 1+\|y\|_{*,\ka} \rp T^{\ep},
\quad\mbox{ with }\quad
c=c_{f,X^1,\tilde \ga,\ka},
$$
which yields that, whenever $c T^\eps \le 1/2$,  the ball $B$ defined by
(\ref{def:ball}) is left invariant by the map $\Gamma$.

\vspace{0.2cm}

Now that the invariance of $B$ has been shown, the contraction property
for $\Gamma$ in a small interval $\ott$ is a matter of standard arguments,
and is left to the reader for sake of conciseness. Let us just mention that $f$
is only supposed to be locally Lipschitz when considered as a function
from $\cb_\ka$ to $\cb_\ka$. However, we are able to establish the
contraction property here,
due to the fact that we are confined to the ball $B$.
This gives the existence and uniqueness result for equation (\ref{algmil})
in the small interval $\ott$ whose size does not depend on the initial condition $\psi$.
The construction
of a global unique solution from the solution in $\ott$ is also quite
standard, and its proof will be omitted here.


\end{proof}

%%%%%%%%%%%%%%%%%%%%%%%%%%%%%%%%%%%%%%%%%%%%%%%%%%%%%%%%%%%%%%%%%%%%%%%%%%
%%%%%%%%%%%%%%%%%%%%%%%%%%%%%%%%%%%%%%%%%%%%%%%%%%%%%%%%%%%%%%%%%%%%%%%%%%

\subsection{Application: the fractional heat equation}

Let us see now how the abstract results of Section \ref{sec:young}
can be applied in the case of the heat equation driven by a
fractional Brownian motion defined at Section \ref{sec:circle}.
Recall  that this means that we wish to solve equation (\ref{miljto})
in case $A=\Delta-\id$, where $\Delta$ is the Laplace operator on the circle,
$x$ is a fractional Brownian motion  defined by the
covariance function (\ref{covstr}), $\cb_{\ka}$ stands
for the usual Sobolev space on  $\ou$, and $f:\cb_\ka\to\cb_\ka$
is defined by $[f(y)](\xi)=\si(y(\xi))$ for $\xi\in\ou$ and a smooth
function $\si:\R\to\R$. In other words, we will try to solve the equation
\begin{equation}\label{mul:fbm}
y(t,\xi)=
\iou G_t(\xi,\eta) \psi(\eta) d\eta
+\iot\iou G_{t-s}(\xi,\eta) X(ds,d\eta)\, \si(y_s(\eta)),
\end{equation}
where the last integral has to be understood in the sense of
Theorem \ref{th:young}. Notice that we have chosen here a multiparametric
formulation for our equation, for computational purposes. However,
as mentioned in the introduction, this setting can be translated easily
into the infinite-dimensional one. Now, the application of
Theorem \ref{th:young-spde} in this context amounts to define
an incremental operator $X^1$  related to our problem, and then
to show that Hypothesis \ref{hyp:n} and \ref{hyp:f} are fulfilled.

\vspace{0.2cm}

Let us give then a natural definition of the operator $X^1$ associated
to our equation: we will set, for $\psi\in\cb$ and $(t,s)\in\cs_2$,
\begin{eqnarray}\label{def:nts-psi}
[X^1_{ts}\psi](\xi)
&=&
[\cj_{ts}(\hdd X S)]\psi(\xi)\\
&=&
\ist\iou G_{t-u}(\xi,\eta_1) X(du,d\eta_1)
\lp \iou G_{v-s}(\eta_1,\eta_2)\psi(\eta_2) d\eta_2 \rp,\nonumber
\end{eqnarray}
which has to be understood now in the Wiener sense, as a centered Gaussian
random variable whose variance is given by (\ref{covstr}). In this
context, the regularity result we obtain on $X^1$ is the following:

%\note{We will need also the case $H=1/2$.}
\begin{proposition}\label{prop:reg-n}
Let $X$ be an infinite dimensional  fractional Brownian motion defined by
the covariance function (\ref{covstr}) for a given $H>1/2$, with
$Q$ given by (\ref{dcp:cov}) for $\nu\in[0,1)$. Suppose that $H+\bar\nu/2>3/4$,
with the convention $\overline\nu = \nu \wedge (1/2)$.
Let $X^1$ be the incremental operator defined by
(\ref{def:nts-psi}). Then for any $\tilde \gamma<H-1/4+\bar\nu/2$, 
$\ka \in (1/4,1/2)$,  $\kappa_0 =\tilde \gamma - \kappa$ and $\gamma < H$  we have 
$$
X^1\in\hcac_2^{\tilde \ga}\cl_{\mbox{\tiny HS}}^{0,-\ka} \cap \hcac_2^{\ko}\cl_{\mbox{\tiny HS}}^{\ka,\ka} \cap \hcac_2^{\ga}\cl_{\mbox{\tiny HS}}^{\ka,-\ka}
$$
almost surely.
\end{proposition}

\begin{remark}
The reader will probably notice that the assumption $\ka>1/4$ is not necessary in order
to prove the proposition above. However, we include it already at this stage, since this restriction
is crucial for Proposition \ref{hyp:si} to hold true.
\end{remark}

The proof of Proposition \ref{prop:reg-n} relies on the following
elementary lemmas, that we label for further use:

\begin{lemma}
\label{lemma:asym-sum}
For any $\alpha < \beta$ such that $\alpha+\beta>1/2$
there exists a constant $C$ such that
$$
\sum_{i,j: i+j=k } \lambda_i^{-\alpha} \lambda_j^{-\beta} \le C \lambda_k^{-\alpha - \overline{ \beta}+1/2},
$$  
where $\overline{\beta} = \min(\beta;1/2)$.
\end{lemma}

\begin{lemma}\label{intg:fr}
Let $a$ and $b$ be two positive constants, and $H>1/2$.
Then the integral
$$
\iou\iou |u-v|^{2H-2} |2-u-v|^{-a} |u+v|^{-b} du dv
$$
is finite whenever $2H-a>0$ and $2H-b>0$.
\end{lemma}

We leave the easy proof of these results to the reader.

\begin{proof}[Proof of Proposition \ref{prop:reg-n}]
We need to prove that the r.v. $X^1$ has a version with the claimed regularity. For random operators,  
up to our knowledge, no standard method is available to prove regularity properties. 
So we have chosen the following simple (though arguably non-optimal) strategy in
order to obtain a regular version: first we determine the kernel associated to the operator $X^1$,
then using the kernel  we estimate its  Hilbert-Schmidt norm in some $L^2$ space. 
This will be enough to apply the modified Garsia-Rodemich-Rumsey lemma ~\ref{lemma:besov} 
and conclude the proof. We will develop now this strategy into several steps, discussing in detail the proof of $
X^1\in\hcac_2^{\tilde \ga}\cl_{\mbox{\tiny HS}}^{0,-\ka}$.
The other path-wise statements can be proven similarly. 



\vspace{0.2cm}

\noindent
{\it Step 1: Definition of a random kernel.}\\
For $(t,s)\in\cs_2$, $X^1_{ts}$ is considered as an operator from
$\cb=L^2(\ou)$ to $\cb_{-\ka}$, and thus
$\|X^1_{ts}\|_{\mbox{\tiny HS},\,\cb\to\cb_{-\ka}}=
\|\ap^{-\ka}X^1_{ts}\|_{\mbox{\tiny HS},\,\cb\to\cb}$, which is the expression
we are going to evaluate. Pick $\psi\in\cb$ smooth enough.
Applying Fubini's theorem for the fractional Brownian motion, we get
\begin{eqnarray*}
\lc \ap^{-\ka}X^1_{ts}\psi\rc(\xi)&=&
\ap^{-\ka} \ist\iou G_{t-u}(\xi,\eta_1)X(du,d\eta_1)
\lp \iou G_{u-s}(\eta_1,\eta_2)\psi(\eta_2) d\eta_2  \rp\\
&=&
\ist\iou G_{t-u}^{-\ka}(\xi,\eta_1)X(du,d\eta_1)
\lp \iou G_{u-s}(\eta_1,\eta_2)\psi(\eta_2) d\eta_2  \rp\\
&=&
\iou K_{ts}(\xi,\eta_2)\psi(\eta_2) d\eta_2,
\end{eqnarray*}
where the kernel $G_{t-u}^{-\ka}$ has been defined at
Section \ref{sec:circle}, and
where $K_{ts}(\xi,\eta)$ is the random kernel on $\ou^2$ defined
by the Wiener integral
$$
K_{ts}(\xi,\eta)
=
\ist\iou G_{t-u}^{-\ka}(\xi,\eta_1) G_{u-s}(\eta_1,\eta) X(du,d\eta_1).
$$
Hence, the Hilbert-Schmidt norm of $X^1_{ts}$, seen as an operator
from $\cb$ to $\cb_{-\ka}$, will be given by
\begin{equation}\label{norm:nts}
\|X^1_{ts}\nhs^{2}
=
\iou  \iou  [K_{ts}(\xi,\eta)]^2\, d\xi d\eta.
\end{equation}
Our next aim will then be to evaluate this last quantity.


\vspace{0.2cm}

\noindent
{\it Step 2: $L^2$ computations.}\\
A direct application of (\ref{covstr}) gives
\begin{multline*}
E[K_{ts}^{2}(\xi,\eta)]
=
c_H \int_{s}^{t}\int_{s}^{t}
\Bigg( \int_{\ou^2}G_{t-u}^{-\ka}(\xi,z) G_{u-s}(\eta,z)
Q(z-\hz)
 \\
\times G_{t-v}^{-\ka}(\xi,\hz) G_{v-s}(\eta,\hz) dz d\hz \Bigg)
|u-v|^{2H-2} du dv.
\end{multline*}
Furthermore, for $z,\hz\in \ou$, it holds that
\begin{eqnarray*}
\iou  G_{t-u}^{-\ka}(\xi,z) G_{t-v}^{-\ka}(\xi,\hz) d\xi
&=&
G_{2t-u-v}^{-2\ka}(z,\hz)\\
\iou  G_{u-s}(\eta,z) G_{v-s}(\eta,\hz) d\eta
&=&
G_{u+v-2s}(z,\hz).
\end{eqnarray*}
Thus, going back to relation (\ref{norm:nts}), we obtain
\begin{eqnarray}\label{ats}
A_{ts}&\equiv &E\lc \|X^1_{ts}\nhs^{2}  \rc \\
&=&
c_H \int_{s}^{t}\int_{s}^{t}
\lp \int_{\ou^2}
Q(z-\hz)
G_{u+v-2s}(z,\hz) G_{2t-u-v}^{-2\ka}(z,\hz) dz d\hz\rp
|u-v|^{2H-2} du dv.\nonumber\\
&=&
c_H \int_0^{\ep}\int_0^{\ep} F(u,v) |u-v|^{2H-2} du dv,\nonumber
\end{eqnarray}
where we have set $\ep=t-s$, and with $F:[0,\ep]^2\to\R_+$ defined by
\begin{equation}\label{def:fuv}
F(u,v)
=
\int_{\ou^2}
Q(z-\hz)
G_{u+v}(z,\hz) G_{2\ep-u-v}^{-2\ka}(z,\hz) dz d\hz.
\end{equation}
Furthermore, plugging the definitions (\ref{dcp:ker}) and (\ref{dcp:cov})
into (\ref{def:fuv}), and invoking the fact that $\{e_n;n\in\Z \}$
is an orthonormal basis of $L^2(\ou)$, we get
$$
F(u,v)
=
\sum_{m,n,l\in D}
\la_n^{-\nu} \la_l^{-2\ka}
e^{-\la_m (u+v)}e^{-\la_l (2\ep-u-v)},
$$
where $D =\{m,n,l \in \mathbb{Z}^3 : m+n+l = 0\}$. Then
$$
A_{ts}
=
c_H \sum_{m,n,l\in D}
\la_n^{-\nu} \la_l^{-2\ka}
\int_0^\ep\int_0^\ep
\frac{e^{-\la_m (u+v)}e^{-\la_l (2\ep-u-v)}}
{|u-v|^{2-2H}} du dv.
$$
Owing now to the fact that $x\mapsto x^a e^{-x}$ is a bounded function
on $\R_+$ for any $a>0$, we obtain, for a constant $c$ which may
change from line to line,
\begin{eqnarray}\label{eq:bnd:A-ts}
A_{ts}&\le& c
\sum_{m,n,l\in D} \la_n^{-\nu} \la_l^{-2\ka} \la_m^{-a}
\int_0^\ep\int_0^\ep
\frac{du dv}{|u-v|^{2-2H} (u+v)^{a}} \nonumber \\
&\le& c \ep^{2H-a}
\sum_{m+n+l=0} \la_n^{-\nu} \la_l^{-2\ka} \la_m^{-a},
\end{eqnarray}
where we have used  Lemma \ref{intg:fr}
under the condition $a<2H$. Let us now analyze the sum. Of course we can write
$$
\sum_{m+n+l=0} \la_n^{-\nu} \la_l^{-2\ka} \la_m^{-a}
\le \sum_{l,k: l+k=0}  \la_l^{-2\ka} \sum_{m,n: m+n=k}\la_m^{-a} \la_n^{-\overline\nu}
$$
Moreover, taking $a=1/2-\overline \nu+\eta$ for some small $\eta > 0$ and using Lemma~\ref{lemma:asym-sum} we have
$$
\sum_{m+n+l=0} \la_n^{-\nu} \la_l^{-2\ka} \la_m^{-a}
\le c \sum_{l,k: l+k=0}  \la_l^{-2\ka} \la_k^{-a-\overline \nu +1/2}
=c \sum_{l,k: l+k=0}  \la_l^{-2\ka} \la_k^{-\eta},
$$
and this sum is always finite under the condition $\ka>1/4$.
Then, going back to (\ref{eq:bnd:A-ts}),  we have found that $A_{ts}\le c \ep^{2\tilde \ga'}$,
for any $\tilde \ga'=H-a/2 < H -1/4 +\overline \nu /2$, where we recall that
$\overline \nu=\inf(\nu;1/2)$.

\vspace{0.2cm}

\noindent
{\it Step 4: $L^p$ estimates.}\\
We will  prove now that, for any $p\ge 1$, we have
\begin{equation}\label{nhsp}
E\lc \|X^1_{ts}\nhs^{2p}  \rc \le c_p (t-s)^{2\tilde \gamma' p},
\quad\mbox{ for }\quad
0\le s<t\le T.
\end{equation}
Indeed, a simple application of H\"older's inequality yields
\begin{eqnarray*}
E\lc \|X^1_{ts}\nhs^{2p}  \rc&=&
\int_{\ou^{2p}}
E\lc \prod_{i=1}^{p} K_{ts}^2(\xi_i,\eta_i) \rc d\xi_1 d\eta_1\ldots
d\xi_p d\eta_p\\
&\le&
\int_{\ou^{2p}}
\prod_{i=1}^{p} E^{1/p}\lc K_{ts}^{2p}(\xi_i,\eta_i) \rc
d\xi_1 d\eta_1\ldots d\xi_p d\eta_p,
\end{eqnarray*}
and since $K_{ts}(\eta_i,\eta_i)$ is a Gaussian variable, we get
\begin{eqnarray*}
E\lc \|X^1_{ts}\nhs^{2p}  \rc&\le&
c_p \int_{\ou^{2p}}
\prod_{i=1}^{p} E\lc K_{ts}^{2}(\xi_i,\eta_i) \rc
d\xi_1 d\eta_1\ldots d\xi_p d\eta_p\\
&=&
c_p \lp \int_{\ou} E\lc K_{ts}^{2}(\xi,\eta) \rc d\xi d\eta\rp^{p}\\
&=& c_p
E^{p}\lc \|X^1_{ts}\nhs^{2}  \rc,
\end{eqnarray*}
which easily yields (\ref{nhsp}).



\vspace{0.2cm}

\noindent
{\it Step 5: Conclusion.}\\
Recall that $X^1$ is considered as an element
of $\hca_2(\cl_{\mbox{\tiny HS}}^{0,-\ka})$.
We can use now inequality~(\ref{eq:generalboundxx}), which can be
read here as:
\begin{equation}\label{n:ga}
\|X^1\|_{\tilde \gamma,0,-\ka} \le C \lc U_{\tilde \gamma+2/p,p,0,-\ka}(n)+\|\td X^1\|_{\tilde \gamma,0,-\ka}\rc,
\end{equation}
in order to bound $\| X^1 \|_{\tilde \ga,0,-\ka}$ for any
$\tilde \ga<\tilde \ga'<H-1/4+\overline \nu/2$. Indeed, if $p$ is large enough, we have
that $\tilde \gamma+2/p<\tilde \ga'$, and the term $U_{\tilde \gamma+2/p,p,0,-\ka}(X^1)$ is easily
handled thanks to (\ref{nhsp}). This yields
\begin{equation}\label{n:uga}
E\lc U_{\tilde \ga',p,0,-\ka}(X^1) \rc<\infty.
\end{equation}
We are now left with the estimation of $\|\td X^1\|_{\tilde \ga}$. However,
remember that $\td X^1=0$ in case of a regular signal $x$, and it is readily checked that this
relation is still valid in the current fractional Brownian setting, so this term is identically zero.
 Thus, we have obtained that
$$
E\lc \|X^1\|_{\tilde \ga,0,-\ka} \rc \le
c E\lc U_{\tilde \ga',p,0,-\ka}(X^1) \rc<\infty,
$$
which implies that $\|X^1\|_{\tilde \gamma,0,-\ka}<\infty$ almost surely, concluding the proof.

\smallskip

Along the same lines as in the preceding steps, some $L^2$ bounds state that 
\begin{equation}\label{nhsp-other-2}
E\lc \|X^1_{ts}\|_{\text{HS},\LL(\BB_\kappa,\BB_\kappa)}^{2}  \rc \le c (t-s)^{2 \kappa'_0 },
\quad\mbox{ for }\quad
0\le s<t\le T
\end{equation}
and
\begin{equation}\label{nhsp-other-3}
E\lc \|X^1_{ts}\|_{\text{HS},\LL(\BB_{-\kappa},\BB_\kappa)}^{2}  \rc \le c (t-s)^{2 \gamma' },
\quad\mbox{ for }\quad
0\le s<t\le T,
\end{equation}
for any $\kappa'_0 < H-1/4-\ka+\overline \nu/2$ and $\gamma'<H$ respectively. Following the same strategy as before, these bounds are enough to prove the remaining assertions of the
proposition.

\end{proof}

Let us see now how this results can be related to our Hypothesis
\ref{hyp:n}. 
Recall that the restriction $\ka > 1/4$ is dictated by the fact that we need to work in
a space $\cb_\ka$ embedded in the space $C(\ou)$ of continuous functions
on $\ou$ in order to prove Proposition \ref{hyp:si} below.

\begin{corollary}
Suppose $X$ is an infinite dimensional  fractional Brownian motion defined by
the covariance function (\ref{covstr}) for a given $H>1/2$, with
$Q$ given by (\ref{dcp:cov}) for $\nu\ge 0$.
Assume moreover that $H>7/8-\overline \nu/2$. Then the incremental operator $X^1$
satisfies Hypothesis \ref{hyp:n} for some
$$
\kappa \in (1/4,1/2), \quad \kappa_0 < H-1/4-\ka+\overline \nu/2, \quad \tilde \gamma = \ko + \ka
$$
\end{corollary}

\begin{proof}
By the previous result we have that $X^1$ has the required regularity for any $1/4<\ka<1/2$,
$
\kappa_0 < H-1/4-\ka+\overline \nu/2
$ and $\tilde \gamma = \ko+\ka<H-1/4+\overline \nu/2$. 
In order to check Hypothesis~\ref{hyp:n},
we now need to require that $\tilde \gamma + \kappa > 1$. In fact,  there exists $1/4<\kappa < \kappa_0$ satisfying this inequality if and only if
$
\tilde \gamma + \kappa_0 > 1
$, i.e. $2H-1/2-\ka+\overline \nu > 1$. This is equivalent to assume
$$
H > 3/4 + \ka/2 - \overline \nu/2 > 7/8 - \overline \nu/2.
$$
In this latter case, it is easily seen that 
there exist $\tilde\gamma,\ka,\ko$ satisfying our requirements.
\end{proof}


\begin{remark}
If we are only interested in obtaining a local solution for our Young PDE, then
the estimate (\ref{eq:bnd-fy2}) can be replaced by a bound in $\cb_\ka$, which will
be quadratic in $y$. Hence, using the fact that
$$
X^1\in\hcac_2^{\gamma}\cl^{\ka,-\kappa} \cap \hcac_2^{\kappa_0}\cl^{\kappa,\kappa} 
$$
for any $\gamma < H$ and $1/4 < \kappa < \kappa_0 <
H-1/4-\ka+\overline \nu/2$, the condition for the construction of the (local)
fix-point map $\Gamma$ becomes $\gamma + \kappa > 1$. To fulfill this
requirement with our fractional Brownian noise we only have to impose $H >
3/4-\overline \nu/4$. This condition is comparable (but a bit worser) with the results
of~\cite{glt}, where the Hilbert spaces
$W^{2\kappa,2}$ were considered, and where we found $H > 3/4-\overline \nu/2$. 
One of the drawback of the approach presented in this
paper is that the esimation of the random operators like $X^1$ in Banach
spaces $W^{2\alpha,p}$ for $p > 2$ seems very difficult. Moreover it
seems that the estimation in the Hilbert-Schmidt norm causes another
small loss of regularity, which means that even in the case
of a ``regular'' noise $\nu =1/2$, our bound on $H$ is $H > 5/8$ and not
$H > 1/2$ as should be natural to expect and found in~\cite{glt}. On
the other hand, as
we will see later, the operator approach seems better suited than the
approach of~\cite{glt} for a true rough-path expansion of SPDEs.
\end{remark}

\vspace{0.2cm}





Now that we have checked the assumptions on $X^1$, let us turn to the hypothesis on the non-linear coefficient $\si$ in equation (\ref{mul:fbm}).
In order to deal with the Sobolev norms, it is worth mentioning that,
instead of working with the spaces $\cb_\ka=H_\ka$ we have used
so far, characterized by their Fourier decomposition, we will consider
the Sobolev spaces $W^{2\ka,2}$, induced by the norms
\begin{equation}\label{norm:wka}
\lc W_\ka(\psi) \rc^2
\equiv
\| \psi\|_{L^2(\ou)}^{2}
+
\int_{\ou^2} \frac{|\psi(\xi)-\psi(\eta)|^2}{|\xi-\eta|^{1+4\ka}} d\xi d\eta.
\end{equation}
These spaces are obviously more convenient than the spaces $\cb_\ka$
for the computations  on $f$, and they are closely related to
these latter spaces, since the following classical relation holds true (see \cite{Ad}):
$$
\cb_{\ka+\ep} \subset W^{2\ka,2} \subset \cb_{\ka-\ep},
\quad\mbox{ for any }\quad
\ep>0.
$$
Using these embeddings we can consider the operator $X^1_{ts}$ going from a space $W^{2\ka,2}$
to a space $\cb_\ka$ by just loosing a little  regularity in $t,s$. Then we can verify that $f$ satisfy a slight modification of Hyp.~\ref{hyp:f}:


\begin{proposition}\label{hyp:si}
Let $\si\in C_b^2(\R)$ be a real-valued function. Then, for any $\ka>1/4$,
the function
$f:W^{2\ka,2}\to W^{2\ka,2}$ defined by $[f(y)](\xi)=\si(y(\xi))$ 
 is locally
Lipschitz,  satisfies $|f(x)|_{W^{2\ka,2}}\le c_f (1+|x|_{W^{2\ka,2}})$ and is globally Lipschitz as a map
$f : \cb \to \cb$.
\end{proposition}

\begin{proof}
Recall that, for our particular situation, $\cb=L^2(\ou)$, and it is
easily checked that, whenever $\si\in C_b^2(\R)$, the function
$f:\cb\to\cb$ is bounded and globally Lipschitz.

\vspace{0.2cm}


%\vspace{0.2cm}

With these considerations in mind, it is readily seen that
$f:W^{2\ka,2}\to W^{2\ka,2}$ has linear growth. In order to check that
$f$ is also locally Lipschitz, note that its gradient can be computed
as follows for $y,h\in W^{2\ka,2}$:
$$
\nabla f(y):W^{2\ka,2}\to W^{2\ka,2},
\qquad
\lc \nabla f(y)\cdot h  \rc (\xi)
=
\si'(y(\xi)) h(\xi).
$$
Let us estimate now the norm (\ref{norm:wka}) of $\nabla f(y)\cdot h$:
first, if $\si'$ is a bounded function, then
\begin{equation}\label{bnd:nld}
\| \nabla f(y)\cdot h \|_{L^2(\ou)}
\le
\| \si'\|_{\infty} \| h\|_{L^2(\ou)}
\le
\| \si'\|_{\infty} \| h\|_{W^{2\ka,2}}.
\end{equation}
As far as the variational term of (\ref{norm:wka}) is concerned, notice
that we have assumed $\ka>1/4$, which means that
$W^{2\ka,2}\subset C(\ou)$, and for any $h\in W^{2\ka,2}$,
$\| h \|_{\infty}\le c \| h\|_{W^{2\ka,2}}$. Thus,
\begin{align}\label{bnd:nvar}
&\int_{\ou^2}
\frac{|\lc \nabla f(y)\cdot h  \rc (\xi)-\lc \nabla f(y)\cdot h  \rc (\eta)|^2}
{|\xi-\eta|^{1+4\ka}} d\xi d\eta \nonumber \\
&\le \| \si'\|_{\infty}^2
\int_{\ou^2}  \frac{|h(\xi)-h(\eta)|^2}{|\xi-\eta|^{1+4\ka}} d\xi d\eta
+
\| h\|_{W^{2\ka,2}}^2 \int_{\ou^2}
\frac{|\si'(y(\xi))-\si'(y(\eta))|^2}{|\xi-\eta|^{1+4\ka}} d\xi d\eta
\nonumber \\
&\le c_{\si} \| h\|_{W^{2\ka,2}}^2
\lc 1+ \| \si''\|_{\infty}^2 \| y\|_{W^{2\ka,2}}^2  \rc.
\end{align}
Putting together (\ref{bnd:nld}) and (\ref{bnd:nvar}), we have thus shown
that
$$
\| \nabla f(y)\|_{\cl(W^{2\ka,2})}
\le
c_{\si} \lp 1+ \| y\|_{W^{2\ka,2}}  \rp,
$$
which easily yields that $f:W^{2\ka,2}\to W^{2\ka,2}$ is locally Lipschitz.


\end{proof}


\begin{remark}
Notice that, in spite of the fact that $\si$ is assumed to be a nicely
behaved coefficient, its interpretation as an application from
$W^{2\ka,2}$ to $W^{2\ka,2}$ does not enjoy the usual assumptions of boundedness made on
coefficients in rough path theory (see e.g. \cite{Lyons,LyonsBook,Lej}).
This is one of the major sources
of problems in our computations, and in general in the extension of rough path theory to SPDEs.
\end{remark}

\vspace{0.2cm}

Let us summarize now the considerations of the current section into the following
theorem:
\begin{theorem}\label{exuni:heat-fbm}
Let $X$ be an infinite-dimensional fractional Brownian motion
on $\ott\times\ou$, defined by the covariance function
(\ref{covstr}) and (\ref{dcp:cov}), with $H>1/2$ and $\nu\in[0,1)$
such that $H>7/8-\overline \nu/2$ and let  $\si\in C_b^2(\R)$. Then, there exists $\ka \in (
1/4,2H-3/2+\overline \nu)$ such that for any initial condition $\psi \in \cb_\ka$, the equation
\begin{equation}\label{eq:calor}
Y(0,\xi) = \psi(\xi), \qquad \partial_t Y(t,\xi)
= \Delta Y(t,\xi) dt + \si(Y(t,\xi)) X(dt,d\xi),
\quad
t\in\ott, \xi\in\ou,
\end{equation}
with periodic boundary conditions, understood in the mild sense given by
(\ref{algmil}), has a unique global solution in $\hca_1^{\ka,\ka}$.
\end{theorem}

\begin{proof}
By Prop.~\ref{hyp:si} the map $f$ is Lipshitz and with linear growth from $\BB_{\ka}$ to $\BB_{\ka-\eps}$ for arbitrarily small $\eps$. As already noted this little mismatch of regularity can be compensated by the time-regularity of $X^1$. Then by a small modification of the arguments of
Theorem \ref{th:young-spde} and by Proposition \ref{prop:reg-n} we can  directly to solve the equation
$$
\partial_t Y(t,\xi)
= (\Delta-\id) Y(t,\xi) dt + \si(Y(t,\xi)) X(dt,d\xi),
\quad
t\in\ott, \xi\in\ou
$$
as a rough evolution equation in $\hca_1^{\ka,\ka}$. Now, if one wants to solve (\ref{eq:calor}),
it is sufficient to get an existence and uniqueness result for
the equation
$$
\partial_t Y(t,\xi)
= (\Delta-\id) Y(t,\xi) dt +
Y(t,\xi) dt +
\si(Y(t,\xi)) X(dt,d\xi),
\quad
t\in\ott, \xi\in\ou,
$$
which can be done along the same lines as for Theorem \ref{th:young-spde},
by taking care of the additional drift term $Y dt$. This step is left to
the reader.
\end{proof}






%%%%%%%%%%%%%%%%%%%%%%%%%%%%%%%%%%%%%%%
%%%%%%%%%%%%%%%%%%%%%%%%%%%%%%%%%%%%%%%

\section{Rough evolution equations: the linear case}
\label{sec:rough}
We pass now to the development of an expansion which allows to consider equation (\ref{eq:evol-young}) in a case which goes beyond the Young theory, in terms of the H\"older regularity of the
driving noise $x$. We start with a simple linear case, i.e. $f\equiv\id$,
which will hopefully lead to a better understanding of our method.


\subsection{Strategy}

Recall that we wish to get some existence and uniqueness results
for the equation
\begin{equation}\label{eq:rough-ev}
\hd y_{ts}=
\ist S_{t  u} dx_u y_u=\cj_{ts}(\hdd x \, y),
\quad\mbox{ and }\quad
y_0=\psi.
\end{equation}
Just like in the case of the Young integral, sketched at the
beginning of Section 4.1,
%\ref{young:integration}
we will proceed as follows:
\begin{enumerate}
\item
Expand (\ref{eq:rough-ev}) as if $x$ were a regular process, until
we get some terms which can be analyzed through the application
of the operators $\hd$ and $\hl$.
\item
Define a natural extension of the notion of integral thanks to the
first step, and show that this allows to integrate a reasonably wide
class of functions.
\item
Solve the equation in the sense given by this notion of integral.
\end{enumerate}
In the current section, we will mostly address the first of these three steps.

\vspace{0.2cm}

If $x$ is a regular process, equation (\ref{eq:rough-ev}) can
be solved by means of the classical evolution theory. Furthermore,
if $y$ designates the unique solution to (\ref{eq:rough-ev}),
then according to our expansion strategy, $y$ also satisfies,
for $t,s\in\cs_2$:
\begin{eqnarray*}
\hd y_{ts}&=&
\ist S_{t  u} dx_u y_u
%\\
%&
=
%&
 \ist S_{t  u} dx_u S_{u  s} y_s
+ \ist S_{t  u} dx_u \hd y_{us}.
\end{eqnarray*}
However, the last term of this equation cannot be defined by applying
the map $\hl$ when $x$ has low time regularity. In order
to cope with this difficulty, let us expand again $\hd y$ by plugging
relation (\ref{eq:rough-ev}) into the previous equation. Doing this twice, we get
\begin{eqnarray*}\label{eq:eq6}
\hd y_{ts}&=&
\ist S_{t  u} dx_u S_{u  s} y_s
+ \ist S_{t  u} dx_u \int_s^u S_{u  v} dx_v y_v\\
&=&
\ist S_{t  u} dx_u S_{u  s} y_s
+ \ist S_{t  u} dx_u \int_s^u S_{u  v} dx_v S_{v  s} y_s
\\ &&+ \ist S_{t  u} dx_u \int_s^u S_{u  v} dx_v \int_s^v S_{v  w} dx_w S_{w  s} y_{s}
\\ &&+ \ist S_{t  u} dx_u \int_s^u S_{u  v} dx_v \int_s^v S_{v  w} dx_w \hd y_{ws}.
\end{eqnarray*}
Thus, going back to our notations on iterated integrals (\ref{eq:not-iter-convol-intg}), we can recast (\ref{eq:rough-ev}) into
\begin{equation}\label{eq:rough-integral}
\hat \delta y  = X^1 y + X^2 y + X^3 y + \cj(\hdd x\, \hdd x\, \hdd x\, \hdd x\, y),
\end{equation}
where, for $t,s\in\cs_2$ and $\phi\in\cb$, the operators $(X^i)_{i=1,2,3}$ are defined by:
\begin{equation}\label{def:levy}
X^{i}_{ts}\phi :=
\cj_{ts}\lp \hdd x X^{i-1} \rp \phi
=\int_s^t S_{t  u} dx_u X^{i-1}_{us}\phi.
\end{equation}
with $X^0_{ts}=S_{t  s}$.
These operators are the new building block we will need in order to
solve equation (\ref{eq:rough-ev}),  and they play the role of the
iterated integrals of rough path theory in our bilinear evolution
context.  
Notice that the last term in equation~(\ref{eq:rough-integral})
is considered as a remainder: suitable assumptions should be made to ensure that it will be
 {\it small} enough. Notice also that we stopped our expansion at the third order. We will
 see that this is the minimum order which allows to handle the Brownian case.

\smallskip

Let us say a few words now about the algebraic properties of the operators $X^i$:
when $x$ is a smooth process, we have for example
\begin{eqnarray*}
\hd X^2_{tus}&=&
\int_u^t S_{t  v} dx_v \int_s^v S_{v  w} dx_w S_{w  s} 
%\\ &&\hspace{3cm}
- \int_u^t S_{t  v} dx_v \int_u^v S_{v  w} dx_w S_{w  u},
\end{eqnarray*}
and using some elementary algebra, we end up with
\begin{eqnarray*}
\hd X^2_{tus}&=&
\int_u^t S_{t  v} dx_v S_{v  u}\int_s^u S_{u  w} dx_w S_{w  s}
%\\&&
+ \int_u^t S_{t  v} dx_v \int_u^v S_{v  w} dx_w S_{w  u}[S_{u  s}-\id]
\\ &=& X^1_{tu} X^1_{us} + X^2_{tu} a_{us}.
\end{eqnarray*}
Thus, taking into account our algebraic
convention (\ref{cvpdt}) and the definition of $\td$ given at (\ref{eq:def-tilde-delta}), 
we have obtaind the relation
$\td X^2 = X^1 X^1$. In a more general way, it is not difficult to show by induction that
\begin{equation*}
\td X^n = \sum_{i=1}^{n-1} X^i X^{n-i},
\end{equation*}
which are exactly the Chen relations in this setting. 

\smallskip



We can now specialize our previous program into the following:
\begin{enumerate}
\item[2a.] Assume that the operator-valued 1-increments $X^1,X^2,X^3$
are defined by some kind of
operation which preserves the usual algebraic relations between integrals
(e.g. use stochastic calculus with respect to an
Hilbert space valued fractional Brownian motion or some other limiting
procedure on discrete sums). They
will be our (step-3) \emph{rough path}.

\item[2b.] Using $(X^1,X^2,X^3)$ define an integration theory for a sufficiently
  large class of functions $\QQ$ so that it will be possible to give a
  meaning to integrals of the form
$
z_t = \int_0^t S_{t  u} dx_u y_u
$
for any $y \in \QQ$.  We will call $\QQ$ the space of paths
\emph{controlled by} $X$.

\item[3'.] Study the map $\Gamma: \QQ \to \QQ$ defined by
$
\Gamma(y)_t = \int_0^t S_{t  u} dx_u y_u,
$
and prove that it has a fixed point $y = \Gamma(y)$ which will be then
a  solution  of the evolution problem (\ref{eq:rough-ev}).
\end{enumerate}


\subsection{Integration of weakly controlled paths}

We start by postulating some reasonable properties for $X^n$.
\begin{hypothesis}\label{hyp:nm}
We will assume that the process $X$ allows to define some
operator-valued increments $X^1,X^2,X^3$, representing morally
(\ref{defnts}) and (\ref{def:levy}) respectively. This amounts for us to suppose
that the $X^i$'s  satisfy the algebraic relations
$$
\td X^1=0,\quad \td X^2=X^1\, X^1, \quad \td X^3= X^1 \, X^2 + X^2\, X^1,
$$
and that the following H\"older-regularity properties holds true:
$$
X^i \in \hcac^{\ga+(n-1)\ko}_2 \LL^{\eta,-\rho}\cap \hcac^{n\ko}_2 \LL^{\eta,\eta},
\quad i=1,2,3,
$$
for some $\eta,\rho \ge 0$ and $\ga,\ko$ such that $\ga = \ko+\eta+\rho$ and 
$\ga + 3 \ko > 1$.
\end{hypothesis}



\vspace{0.2cm}

We will define now the class $\cq$ of processes we wish to be able to
integrate against $x$: in the current situation, it will include
any process which can be decomposed into a part depending on $X^1,X^2$,
plus a remainder term which is assumed to be {\it small} enough.
For the sake of a contraction argument needed below (compare to the Young case)
 we fix a given time regularity
$\ka$ such that $0<\ka< \ko$.
\begin{definition}[Weakly controlled paths]\label{def:weak-controlled}
Let $\psi \in \cb_{\eta}$ be a given initial condition.
A path $y \in \hcac_1^{\ka,\eta}$ is said to be \emph{weakly controlled}
by $X^1,X^2$ if $y_0=\psi$ and $\hd y$ can be decomposed into
\begin{equation}\label{contol:dcp}
\hd y  = X^1 y^1 + X^2 y^2 + y^r, \qquad \hd y^1  = X^1 y^2 +  y^{1,r}
\end{equation}
with $y^i \in \hcac_1^{\ka,\eta}$ $i=1,2$, and a regular part
$y^{r}\in\hcac_2^{\ga+2\ka,\eta}$, $y^{1,r}\in\hcac_2^{\ga+\ka,\eta}$ with 
$\ka<\ko\wedge \eta$.
Furthermore, we asssume that the regularity of $y^1,y^2$
and $y^r,y^{1,r}$ can be related to those of $X$  by the following
relation:
$
\ga+3\ka>1$,
a condition that can be always fullfilled by a suitable choice of
$\ka$ whenever $\ga+3\ko>1$.
Denote this space of controlled paths
by $\cq_{\ka,\eta,\psi}$, or when this does not lead to an ambiguous situation,
simply by $\cq_{\ka,\eta}$ or $\cq$. Moreover, one can define a
semi-norm $\cn$ on $\cq_{\ka,\eta}$ in the following way:
$$
\Norm{y}{\QQ_{\ka}}
= \Norm{y}{\hcac_1^{\ka,\eta}} +
\sum_{i=1,2}\Norm{y^i}{\hcac_1^{\infty,\eta}} +
\sum_{i=1,2} \Norm{y^i}{\hcac_1^{\ka,\eta}}
+ \Norm{y^r}{\hcac_2^{\ga+2\ka,\eta}}
+ \Norm{y^{1,r}}{\hcac_2^{\ga+\ka,\eta}}.
$$
\end{definition}

\begin{remark}
Even if a weakly controlled path is, strictly speaking, given by
      a tuple $(y,y^1,y^2,y^r,$ $y^{1,r})$ we will, with a slight abuse of notations,
      denote it with its first component, i.e. simply $y$.
\end{remark}

\begin{remark}
 The notion of weakly controlled path appeared first in~\cite{rough} in the
finite dimensional context as a way to linearize the space of
rough paths around the driving control. Even if this linearization
does not preserve the whole structure of the space of rough paths, it is
enough to find solutions of rough differential equations. 
%In some intuitive sense controlled paths are close to the driving rough path,
%in a kind of Radon-Nykodim sense.
\end{remark}

\vspace{0.2cm}

With these notations at hand, we will try to implement now the strategy
designed at the beginning of Section 4.1 in order to integrate a weakly
controlled process $y$: let us first assume $x$ is a smooth process,
and $y\in\cq$. Then $\cj(\hdd x y)$ is well defined, and thanks
to equations (\ref{eq:split-1}) and (\ref{contol:dcp}), we have
$$
\cj(\hdd x y)=
\cj(\hdd x S) y + \cj(\hdd x \hd y)
= \cj(\hdd x S) y + \cj(\hdd x \, X^1 y^1) + \cj(\hdd x \, X^2 y^2) + \cj(\hdd x \, y^r).
$$
Furthermore, for $s<t$, the term $\cj_{ts}(\hdd x \, X^1 y^1)$ above only involves
$y^1_s$, and hence the increment $\cj_{ts}(\hdd x \, X^1 y^1)$ is equal to
$\cj_{ts}(\hdd x \, X^1) y^1_s$, that is $X^2_{ts}y^1_s$. This yields:
\beq\label{dcp:weak}
\cj(\hdd x y)=X^1 y + X^2 y^1 + X^3 y^2 + \cj(\hdd x\,  y^r).
\eeq
Note that, in this last expression, the terms $X^1 y$, $X^2 y^1$ and
$X^3 y^2$ are well-defined under
Hypothesis \ref{hyp:nm}. In order to have a well-defined expression for
$\cj(\hdd x y)$ in the rough case, it remains to handle the term
$\cj(\hdd x  y^r)$. Then, let us write
$$
\cj(\hdd x  y^r)= \cj(\hdd x y) - X^1 y - X^2 y^1 - X^3 y^2,
$$
and let us analyze this relation by applying $\hd$ to both sides. This
gives
\begin{equation}\label{first:hdyr}
\hd \lc \cj(\hdd x  y^r) \rc
=
-\hd \lc  X^1 y \rc - \hd \lc X^2 y^1 \rc - \hd \lc X^3 y^2 \rc,
\end{equation}
and notice that in the last expression, $\hd[\cj(\hdd x  y^r)]\neq 0$, since
$y^r$ belongs to $\hcac_2$ instead of $\hcac_1$.
Moreover, a slight extension of Lemma \ref{daf} shows that, for
$M\in\hca_2(\cl(V))$ and $L\in\hca_1(V)$, we have
\begin{equation*}
\hd(M\, L)= \hd M\, L- M\, \der L= \td M\, L- M\, \hd L.
\end{equation*}
Applying this elementary relation to (\ref{first:hdyr}), we
end up with
\begin{eqnarray}\label{sec:hdyr}
\hd \lc \cj(\hdd x  y^r) \rc&=&
-\td X^1\, y + X^1\, \hd y - \td X^2 \, y^1 + X^2\, \hd y^1- \td X^3 \, y^2 + X^3\, \hd y^2\nonumber\\
&=& X^1 (\hd y - X^1 y^1 - X^2 y^2) + X^2 (\hd y^1 - X^1 y^2) + X^3
\hd y^2
\nonumber\\
&=& X^1 y^r + X^2 y^{1,r} + X^3 \hd y^2
\end{eqnarray}
under our hypothesis on $y$ and $X$ we have the following
regularities: 
$$
 X^1 y^r \in \hcac_3^{\ga+3\ka,-\rho}, \qquad X^2 y^{1,r} \in
 \hcac_3^{\gamma+\ko+2\ka,-\rho}, \qquad X^3 \hd y^2  \in
 \hcac_3^{\gamma+2\ko+\ka,-\rho},
$$
so if $\ga+3\ka >1$ we can apply the operator $\hl$ and express $\cj(\hdd x  y)$
in terms of $\hd$ and $\hl$ only. Plugging (\ref{sec:hdyr})
into (\ref{dcp:weak}) we get
\begin{equation}\label{def:roughintg}
\cj(\hdd x  y)=
 X^1 y + X^2 y^1 + X^3 y^2 + \hl\lp X^1 y^r + X^2 y^{1,r} + X^3 \hd
 y^2  \rp.
\end{equation}
Similarly to what we did in the Young case, we are now able to invert
the procedure which lead to relation (\ref{def:roughintg}), by
just invoking the assumptions made on $X^i$ and $y$:
\begin{theorem}\label{th:int-rough-lin}
Let $x$ be a path such that $X^i, i=1,2,3$ are well-defined, and such that
Hypothesis \ref{hyp:nm} holds true.
 Let also $y\in\cq_{\ka,\eta,\psi}$ for
$0<\ka<\ko<\ga-\ka$ and $\ka \le \eta$.
Define $z\in\hcac_1(\cb_\eta)$
such that $z_0=\psi$ and $\hd z$ satisfies:
$$
\hd z\equiv
\cj(\hdd x  y)=
 X^1 y + X^2 y^1 + X^3 y^2 + \hl\lp X^1 y^r + X^2 y^{1,r} + X^3 \hd
 y^2  \rp.
$$
and let $z^1 = y$, $z^2 = y^1$, $z^{1,r}=X^2 y^2 + y^{r}$ so that
$
\hd z^1 = X^1 z^2 + z^{1,r}
$.
Then
\begin{enumerate}
\item
$z$ is well-defined as an element of $\cq_{\ka,\eta}$, and coincides with the usual Riemann convolution of $y$ by $x$ in case $x$ and $y$ are smooth processes.
\item
The semi-norm of $z$ in $\cq_{\ka,\eta}$ can be estimated as:
\begin{equation}\label{bn:hdz-rough}
\cn[z;\cq_{\ka,\eta}]\le
c_{X} T^{\ko-\ka}
\lp \|\psi\|_{\cb_{\ka}} + \cn[y;\cq_{\ka,\eta}]\rp,
\end{equation}
for a positive constant $c_{X}$ depending only on $X^i$, $i=1,2,3$.
\item
It holds that, for any $0\le s<t\le T$
\begin{equation}\label{cvgce:rsums}
\cj_{ts}(\hdd x y)
=\lim_{|\Pi_{ts}|\to 0}\sum_{i=0}^n
S_{t  t_{i+1}}
\lc X^1_{t_{i+1},t_{i}} y_{t_{i}}+X^2_{t_{i+1},t_{i}}
y^1_{t_{i}}+X^3_{t_{i+1},t_{i}} y^2_{t_{i}} \rc,
\end{equation}
where the limit is over all partitions
$\Pi_{ts} = \{t_0=t,\dots,t_n=s\}$
of $[s,t]$ as the mesh of the partition goes to zero.
\end{enumerate}
\end{theorem}

\begin{proof}
We will divide again this proof in several steps.

\vspace{0.2cm}

\noindent
{\it Step 1:}
Let us start by evaluating the regularity of the terms in the right
hand side of (\ref{def:roughintg}), i.e.
$$
A=X^1 y,\qquad B=X^2 y^1, \qquad C=X^3 y^2 ,\qquad D=\hl\lp  X^1 y^r + X^2 y^{1,r} + X^3 \hd
 y^2  \rp,
$$
under our standing assumptions.

\smallskip

In order to bound $A$, we will first
estimate $\|y_s\|_{\cb}$ itself for $s\le T$: if $y\in\cq_{\eta}$, we
have $y_s=S_s\psi+\hd y_{s0}$, and hence
\begin{equation}\label{bn:y}
\| y\|_{\infty,\cb_\eta}\le
\|\psi\|_{\cb_{\eta}} + T^{\ka} \cn[y;\hca_1^{\ka,\eta}].
\end{equation}
In particular, $y$ is bounded in $\cb_\eta$ on $\ott$. Thus, if
$X^1 \in\hcac_2^{\ga}\cl^{\eta,-\rho}\cap\hcac_2^{\ko}\cl^{\eta,\eta}$, we
have $X^1 y\in\hcac_2^{\ga,-\rho}$,
and also $X^1 y\in\hcac_2^{\ko,\eta}$. Moreover,
$$
\| X^1_{ts} y_s\|_{\cb^\eta}
\le
\| X^1\|_{\ko,\eta,\eta} (t-s)^{\ko} 
\lp \|\psi\|_{\cb_{\eta}} + \cn[y;\hcac_1^{\ka,\eta}]  \rp,
$$
and thus
\begin{equation}\label{bn:ny}
\cn[X^1 y;\hcac_2^{\ka,\eta}]
\le
\| X^1 \|_{\ko,\eta,\eta}
\lp \|\psi\|_{\cb_{\eta}} + \cn[y;\hcac_1^{\ka,\eta}]  \rp T^{\ka-\ko}.
\end{equation}

\smallskip

Let us estimate now the term $B$, that is $\cn[X^2 y^1;\hcac_2^{\ga+\ka,\eta}]$:
since $y^1\in\hcac_1^{\infty,\eta}$ and $X^2 \in\hcac_2^{2\ko}\cl^{\eta,\eta}$,
we obtain again that $X^2 y^1\in\hcac_2^{2\ko+\ka,\eta}$, and we have
\begin{equation}\label{bn:myn}
\cn[X^2 y^1;\hcac_2^{2\ka,\eta}]
\le
\| X^2 \|_{2\ko,\eta,\eta} \,
\cn[y^1;\hcac_1^{\infty,\eta}]   T^{2(\ko-\ka)}.
\end{equation}

\smallskip

The term $C$ can now be bounded along the same lines as for $A$ and
$B$. Moreover for the term $D$, as we already observed above,
 $X^1 y^r\in\hcac_3^{\ga+3\ka,-\rho}$,  $X^2 y^{1,r}\in\hcac_3^{\ga+\ko+2\ka,-\rho}$
and $X^3\hd y^2\in\hcac_3^{\ga+2\ko+\ka,-\rho}$,
and observe that we have assumed that $\ga+3\ka > 1$.
Thus
the operator $\hl$ can be applied to $ X^1 y^r + X^2 y^{1,r} + X^3 \hd
 y^2 $,
and invoking inequality (\ref{ineqhla}),
we get that
\begin{align}\label{bn:hlnyr}
&\| \hl\lp  X^1 y^r + X^2 y^{1,r} + X^3 \hd
 y^2 \rp \|_{\ga+3\ka,-\rho}
\le
c \|  X^1 y^r + X^2 y^{1,r} + X^3 \hd
 y^2  \|_{\ga+3\ka,-\rho}\nonumber\\
&\le c
 \lp \| X^1\|_{\ga,\eta,-\rho}\, \cn[y^r;\hcac_2^{3\ka,\eta}]
+ \| X^2\|_{\ga+\ko,\eta,-\rho}\, \cn[y^{1,r};\hcac_2^{2\ka,\eta}]
+ \| X^3\|_{\ga+2\ko,\eta,-\rho}\, \cn[y^{2};\hcac_2^{\ka,\eta}]
\rp.
\end{align}
Summarizing inequalities (\ref{bn:y})-(\ref{bn:hlnyr}), we have obtained
that $z$ is a well-defined element of $\hcac_1^{\ka,\eta}$, and that it
satisfies
$$
\|\hd z \|_{\ka,\eta}\le
c_{n,m} T^{\ko-\ka}
\lp \|\psi\|_{\cb_{\eta}} + \cn[y;\cq_{\ka,\eta}]\rp.
$$


\vspace{0.2cm}

\noindent
{\it Step 2:}
Let us estimate now $z$ as an element of $\cq_{\ka,\eta}$. The natural
decomposition of $\hd z$ is obviously $\hd z = X  z^1 + X^2 z^2 + z^r$, with
$$
z^1 =y, \quad z^2=y^1
\quad\mbox{ and }\quad
z^r= X^3 y^2  + \hl\lp   X^1 y^r + X^2 y^{1,r} + X^3 \hd
 y^2 \rp.
$$
It is now easily checked, along the same lines as for Step 1, that
$z$ satisfies relation (\ref{bn:hdz-rough}).

\vspace{0.2cm}

\noindent
{\it Step 3:}
In order to see how to get the convergence of the Riemann sums
to $\cj(\hdd x y)$ it is enough to remark that
 $\hd z$ can be written as
$
\hd z 
= (\id-\hl\hd)\lc  X^1 y+X^2 y^1 +X^3 y^2  \rc
$.
Applying Corollary \ref{cor:conv-integration}, we now get
relation (\ref{cvgce:rsums}).

\end{proof}




\begin{remark}
The space of weakly controlled paths is a vector space with respect to
the action of $\R$  but not with respect to other interesting linear endomorphisms
of $\BB$. The problem lies in the fact that for general linear $L :
\BB \to \BB$ we can have $\hd L y \neq L \hd y$ since $L$ does not
necessarily commute with the semigroup
(which appears in the definition of $\hd =\der - a$).
%This will be the cause of some troubles when dealing with
%integrals of nonlinear functions of weakly controlled paths.
\end{remark}

\subsection{Linear evolution problem}\label{sec:lin-evolution}

Let us turn now to the main aim of this section, which is to
get an existence and uniqueness result for equation (\ref{eq:rough-ev}):
\begin{theorem}\label{thm:lin-rough}
Assume that Hypothesis~\ref{hyp:nm} holds for the triple of incremental operators $X^1,X^2,X^3$
with
$\ga,\ko,\ka,\eta,\rho$ such that $\ga = \ko+\eta+\rho$, 
$\ga + 3 \ko > 1$ and $\ka<\ko$. Then
\begin{enumerate}
\item
Equation (\ref{eq:rough-ev}) admits a unique solution
$y\in\cq_{\eta}$.
\item
The map $(\psi,X^1,X^2,X^3)\mapsto y$ is continuous.
\item
For $(t,s)\in\cs_2$, the map
$\Phi_{ts}:\cb_{\eta}\to \cb_{\eta}$, such that $\Phi_{ts}\psi=y_t$ when
$y_s=\psi$ and $\hd y_{ts}=\cj_{ts}(\hdd x y)$
is a bounded linear endomorphism of $\cb_{\eta}$, and it satisfies
the cocycle property $\Phi_{tu} \Phi_{us} = \Phi_{ts}$.
\end{enumerate}
\end{theorem}


\begin{proof}
Like in the Young case, the solution $y$ will be identified as the fixed point
of the map $\Gamma:\cq_{\ka,\eta}\to\cq_{\ka,\eta}$ defined by $z=\Gamma(y)$,
with $z_0=\psi$ and
$
\hd z_{ts}=\cj_{ts}( \hdd x \, y )
$.
And here again, we will concentrate on the fact that, on a small
interval $\ott$, the ball
$$
B=\lcl
y;\, y_0=\psi, \cn[y;\cq_{\ka,\eta}]\le \|\psi\|_{\cb_{\eta}}
\rcl
$$
is left invariant by the map $\Gamma$.

\vspace{0.2cm}

Indeed, whenever $y\in B$, then Theorem \ref{th:int-rough-lin}
asserts that for $z=\Gamma(y)$,
the following estimate holds true:
$$
\cn[z;\cq_{\ka,\eta}]\le
c_{X} T^{\ko-\ka}
\lp \|\psi\|_{\cb_{\eta}} + \cn[y;\cq_{\ka,\eta}]\rp.
$$
Hence, if one chooses a small enough $T$, so that
$c_{n,m} T^{\ko-\ka}<1/2$, it is readily checked that
$\cn[z;\cq_{\ka,\eta}]\le \|\psi\|_{\cb_{\eta}}$, which proves that $z\in B$.
The contraction property is now a matter of standard arguments,
and the remainder of the theorem follows easily.


\end{proof}

\subsection{Application: stochastic heat equation}

In the sequel of the paper, for sake of simplicity, the generic
situation of a process $X$ with $\ga$-H\"older continuity in time
with $\ga\le 1/2$ will be the case of an infinite dimensional
Brownian motion, given by the covariance function
(\ref{cov:bm}). For this special process, we will try to
construct a pathwise solution to the linear stochastic heat
equation on $\ou$.
At the end of the section, we will give some hints
about the way the fractional Brownian case should be treated.

\subsubsection{The Brownian case}

Like in the Young case, the key step in order to apply Theorem
\ref{thm:lin-rough} to the Brownian setting is to define $(X^i)_{i=1,2,3}$
in a reasonable way, and then to check
Hypothesis \ref{hyp:nm}. We have chosen here to deal with an It\^o
type definition for $X^n$, and we get the following result:
\begin{proposition}\label{reg:m-brownian}
Let $X$ be an infinite dimensional Brownian motion defined by the
covariance structure (\ref{cov:bm}), with $Q$ given by (\ref{dcp:cov})
for $\nu\in\ou$. For $n=1,2,3$, let $X^n$  be the incremental operators
given by (\ref{def:nts-psi}) and (\ref{def:levy}) respectively,
where the stochastic integral has to be understood in the It\^o sense
(see e.g. \cite{Da,Wa} for a complete definition). Then, almost surely,
$$
X^n\in\hcac_2^{\ga+(n-1)\ko}\cl_{\mbox{\tiny HS}}^{\eta,-\rho} \cap \hcac_2^{n\ko}\cl_{\mbox{\tiny HS}}^{\eta,\eta}
$$ 
for any $\eta > 1/4$, $\ga > \ko > \ka$ satisfying
$$
\ko < 1/4-\eta+\overline \nu/2, \quad\mbox{and}\quad \ga < 1/2,
$$
with $\bar\nu=\inf(\nu;1/2)$.
\end{proposition}



\begin{proof}%[Proof of Proposition \ref{reg:m-brownian}]
We have already proved the regularity of $X^1$ in the fractional Brownian case.  The proof
 in the current case would be similar, and
we omit it. It will be enough to take $\eta = 1/4+\ep$, $\ko = \overline \nu/2 -2\ep$ and $\ga = \ko + \rho+\ep$, $\rho =1/4-\overline \nu/2$ for a given small $\ep >0$.

 Let us concentrate then on
the regularity properties of $X^2$: we will prove in fact first that
$X^2\in\hca_2^{\ga+\ko}\cl_{\mbox{\tiny HS}}^{\eta,-\rho}$, and for this,
we will proceed along the same lines as for the proof of
Proposition \ref{prop:reg-n}. Let us sketch the main steps which have to
be followed.


\vspace{0.2cm}
\noindent
{\it Step 1:}
First of all, we have to estimate
$\|\ap^{-\rho}X^2_{ts}A^{-\eta}\|_{\mbox{\tiny HS};\cb\to\cb}$, and it is readily
checked that $\ap^{-\rho}X^2_{ts}A^{-\eta}$ is represented by the kernel
$$
\tk_{ts}(\xi,\eta)
=
\ist\iou G_{t-u}^{-\rho}(\xi,\eta_1) X(du, d\eta_1)
\int_s^u \iou G_{u-v}(\eta_1,\eta_2) X(dv, d\eta_2)
G^{-\eta}_{v-s}(\eta_2,\eta).
$$
Thus, when considered as an operator from $\cb^\eta$ to $\cb_{-\rho}$,
we obtain that
\begin{equation}\label{exp:hs-mts}
E\lc \| X^2_{ts}\nhs^2 \rc
=
\int_{\ou^2} E\lc ( \tk_{ts}(\xi,\eta) )^2 \rc d\xi d\eta.
\end{equation}
Moreover, some standard considerations about iterated integrals
for Brownian noises (see e.g. \cite{Wa,Da}) yield
$$
E\lc ( \tk_{ts}(\xi,\eta) )^2 \rc
=
\ist du\int_{\ou^2} G_{t-u}^{-\rho}(\xi,\eta_1)
G_{t-u}^{-\rho}(\xi,\hat\eta_1) Q(\eta_1-\hat\eta_1)
H_{us}(\eta,\eta_1,\hat\eta_1)d\eta_1 d\hat\eta_1,
$$
with
\begin{align*}
&H_{us}(\eta,\eta_1,\hat\eta_1)\\
&=
\int_s^u dv \int_{\ou^2}
G_{u-v}(\eta_1,\eta_2) G^{-\eta}_{v-s}(\eta_2,\eta) Q(\eta_2-\hat\eta_2)
G_{u-v}(\hat\eta_1,\hat\eta_2) G^{-\eta}_{v-s}(\hat\eta_2,\eta)
d\eta_2 d\hat\eta_2.
\end{align*}
Plugging this equality into (\ref{exp:hs-mts}), we end up with
\begin{equation}\label{mts:first}
E\lc \| X^2_{ts}\nhs^2 \rc
=
\int_0^\ep du \int_{\ou^2}
G_{2(\ep-u)}^{-2\rho}(\eta_1,\hat\eta_1) \Psi_u(\eta_1,\hat\eta_1)
Q(\eta_1-\hat\eta_1) d\eta_1 d\hat\eta_1,
\end{equation}
where we have set $\ep=t-s$ and
$$
\Psi_u(\eta_1,\hat\eta_1)
= \int_0^u dv \int_{\ou^2} G_{u-v}(\eta_1,\eta_2)
G_{u-v}(\hat\eta_1,\hat\eta_2) G_{2v}^{-2\eta}(\eta_2,\hat\eta_2)
Q(\eta_2-\hat\eta_2) d\eta_2 d\hat\eta_2.
$$
Furthermore, using the spectral decomposition of $G_t$ and $Q$,
introduced respectively by (\ref{dcp:ker}) and (\ref{dcp:cov}),
we obtain
\begin{align*}
&\int_{\ou^2} G_{u-v}(\eta_1,\eta_2)
G_{u-v}(\hat\eta_1,\hat\eta_2) G^{-2\eta}_{2v}(\eta_2,\hat\eta_2)
Q(\eta_2-\hat\eta_2) d\eta_2 d\hat\eta_2\\
&=
\sum_{i,j,k,l\in\Z}
e_i(\eta_1) e_j(\hat\eta_1)
\frac{e^{-(\la_i+\la_j)[u-v]}e^{-2 \la_k v}}{\la_l^{\nu} \la_k^{2\eta}}
\1_{\{ i-k-l=0 \}} \1_{\{ j+k+l=0 \}}.
\end{align*}
Injecting again this value into (\ref{mts:first}) and using the fact that $\lambda_{-i} = \lambda_i$, we have that
$$
E\lc \| X^2_{ts}\nhs^2 \rc
=
\sum_{i,j,k,l,m,n\in E}
\frac{1}{\la_l^\nu\la_m^{2\rho}\la_n^\nu \la_l^\nu \la_k^{2\eta}}
\int_0^\ep du \, e^{-2 \la_m (\ep-u)}
\int_0^u dv \, e^{-2\la_j[u-v]}e^{-2 \la_k v},
$$
with
\beq\label{def:set-e}
E=\lcl
j,k,l,m,n\in\Z;\, m+n=j, \, k+l=-j
\rcl.
\eeq
Thus, we get
\begin{eqnarray*}
E\lc \| X^2_{ts}\nhs^2 \rc&=&
\sum_{j,k,l,m,n\in E}
\frac{1}{\la_m^{2\rho}\la_n^\nu  \la_l^\nu \la_k^{2\eta}}
\int_0^\ep du \int_0^u dv \, e^{-2\la_m (\ep-u)-2\la_j[u-v]-2\la_k v}\\
&\le& c
\sum_{j,k,l,m,n\in E}
\frac{1}{\la_j^b  \la_m^{2\rho+a}\la_n^\nu  \la_k^{2\eta} \la_l^\nu }
\int_0^\ep du \int_0^u dv \,
 du
\frac{dv}{(\ep-u)^a (u-v)^b}\\
&\le& c  \ep^{2-a-b}
\sum_{j,k,l,m,n\in E}
\frac{1}{\la_j^b  \la_m^{2\rho+a}\la_n^\nu  \la_k^{2\eta} \la_l^\nu }
\int_0^1 du \int_0^u dv \,
 du
\frac{dv}{(u-v)^b }.
%\\
%&\le& c  \ep^{2-(a+b+c)},
\end{eqnarray*}
The double integral above is finite whenever $a,b \in (0,1)$, while the sum can be handled along the following lines: first rewrite
$$
S = 
\sum_{m,n,j:m+n-j=0}
\frac{1}{\la_j^b  \la_m^{2\rho+a}\la_n^\nu }\sum_{k,l:k+l=-j}\frac{1}{ \la_k^{2\eta} \la_l^\nu },
$$
and observe that,
thanks to Lemma~\ref{lemma:asym-sum} and according to our hypothesis $\eta>1/4$, we have
$$
\sum_{k,l:k+l=-j}\frac{1}{ \la_k^{2\eta} \la_l^\nu } \le \sum_{k,l:k+l=-j}\frac{1}{ \la_k^{2\eta} \la_l^{\overline\nu} } \le C \lambda_j^{-\overline \nu},
$$
where $C$ stands again for a positive constant which can change from line to line.
Then
$$
S \le C
\sum_{m,n,j:m+n-j=0}
\frac{1}{\la_j^{b+\overline \nu}  \la_m^{2\rho+a}\la_n^\nu }
\le C \sum_m 
\frac{1}{\la_m^{2\rho+a} } \sum_{n,j:n-j=-m}
\frac{1}{\la_j^{b+\overline \nu}  \la_n^{\overline \nu} }
$$
and choose $b > 1/2 -2\overline \nu$,
so that another application of Lemma~\ref{lemma:asym-sum} gives $$
S \le C \sum_m
1/{\la_m^{2\rho+a+b+2\overline \nu -1/2} }.
$$
This latter sum is  finite when $a+b > 1-2\rho-2\overline \nu$. Then for any
$\theta<\theta^*$ such that  $2\theta^* < 1+2\rho+2\overline{\nu} $ we have found that
 $E\lc \| X^2_{ts}\nhs^2 \rc \le c \eps^{2-a-b} \le c \eps^{2\theta}.$


\vspace{0.2cm}

\noindent
{\it Step 2:}
One can go from $L^2$ to $L^p$ estimates for $m$ just like in Proposition
\ref{prop:reg-n} Step 4: indeed, we have
$$
E\lc \|X^2_{ts}\nhs^{2p}  \rc
\le
c_p \int_{\ou^{2p}}
\prod_{i=1}^{p} E^{1/p}\lc \tk_{ts}^{2p}(\xi_i,\eta_i) \rc
d\xi_1 d\eta_1\ldots d\xi_p d\eta_p.
$$
Moreover, $\tk_{ts}(\xi_i,\eta_i)$ is a variable of the second chaos
$\ch_2$ with respect to the Gaussian field $X$, and invoking
\cite[Relation (1.61)]{Nbk}, the $L^2$ and $L^{2p}$ norms on
$\ch_2$ are equivalent. Thus, for any integer $p\ge 1$, there exists
a constant $c_p$ such that
$$
E\lc \|X^2_{ts}\nhs^{2p}  \rc
\le
c_p\,
(t-s)^{2p\theta}.
$$


\vspace{0.2cm}

\noindent
{\it Step 3:}
We will conclude now thanks to Lemma \ref{lemma:besov}, which
reads here as:
\begin{eqnarray*}
\|X^2  \|_{\ga_2,\eta,-\rho}&\le&
c\lc U_{\ga_2+2/p,p,\eta,-\rho}(X^2)
+ \| \td X^2 \|_{\ga_2,\eta,-\rho}  \rc\\
&=&
c\lc U_{\ga_2+2/p,p,\eta,-\rho}(X^2)
+ \| X^1 X^1\|_{\ga_2,\eta,-\rho}  \rc,
\end{eqnarray*}
for any integer $p\ge 1$. According to the previous step,
it is then easily checked that, for any
$\ga_2<\theta^* = 1/2+\overline \nu+\rho$, and $p$ large enough, the term
$U_{\ga_2+2/p,p,\eta,-\rho}(X^2)$ can be bounded almost surely by
a finite constant. Recall now that we have chosen $\ga=\ko+\rho+\eta$,
$\ko=\bar\nu/2-2\ep$ and $\eta=1/4+\ep$. Thus
$$
\ga+\ko= 2\ko + \rho+\eta= \overline \nu-\eta+1/4+\rho < \theta^*,
$$
and hence $U_{\ga_3+2/p,p,\eta,-\rho}(X^2)<\infty$ for any $\ga_3\le\ga+\ko$.

\smallskip

Let us treat now the term $X^1X^1$. Along the same lines as in Proposition~\ref{prop:reg-n}, it can be shown that
$X^1\in\hcac_2^{\ga}\cl^{\eta,-\rho}$ and $X^1\in\hcac_2^{\ko}\cl^{\eta,\eta}$. Hence,
by composition of operators, we get
$X^1 X^1\in\hcac_{3}^{\ga+\ko}\cl^{\eta,-\rho}$, which means that
$\|X^1X^1\|_{\ga_2,\eta,-\rho}$ is finite for any
$\ga_3\le \ga+\ko$.
 Summing up this short discussion,
we have obtained that
$$
\|X^2\|_{\ga_2,\eta,-\rho}
\mbox{ finite a.s. for any }
\ga_3\le \ga+\ko.
$$
One can proceed then to prove that
$X^2\in\hcac_2^{2\ko}\cl^{\eta,\eta}$ by a slight elaboration of
the computations above. This easy exercise is left to the reader.

\smallskip

The proof for the operator $X^3$ follows the same lines and will not be
reported. Indeed, we prefer to concentrate on the regularity properties of higher order operators in
the more complex situation of Section~\ref{sec:brownian-case-polynomial}.

\end{proof}


We are now able to apply our abstract results to the stochastic heat equation:
\begin{theorem}
\label{th:brown-bilin}
Let $X$ be an infinite dimensional Brownian motion on $\ott\times\ou$,
defined by the covariance function given by (\ref{cov:bm}) and
(\ref{dcp:cov}) with $\nu > 1/3$. Then there exists $\eta>1/4$,  $0<\ka<\gamma<1/2$ such that $\ka<\ko$ and $\ga+3\ka>1$ such that, for any $\psi \in \cb_{\eta}$ the equation
\begin{equation*}
Y(0,\xi) = \psi(\xi), \qquad \partial_t Y(t,\xi)
= \Delta Y(t,\xi) dt + Y(t,\xi) X(dt,d\xi),
\quad
t\in\ott, \xi\in\ou,
\end{equation*}
with periodic boundary conditions, understood as equation (\ref{eq:rough-ev}),
has a unique solution in $\cq_{\ka,\eta,\psi}$.
\end{theorem}

\begin{proof}
Like in the proof of Theorem \ref{exuni:heat-fbm},
the claim is readily checked once
we have shown that $X^n$, $n=1,2,3$ satisfy Hypothesis
\ref{hyp:nm}. This amount to check that there exist $\ko>\ka$ and $\ga
< 1/2$ such
that
$$
\ko < 1/4 - \eta + \overline{\nu}/2, \qquad \ga + 3\ka > 1.
$$
However, thanks to Proposition \ref{reg:m-brownian},
it is enough to take $\eta =1/4+\ep$, $\ko =  \overline{\nu}/2 -2\ep$, $\ka = \ko-\ep$, 
$\rho=1/4-\bar\nu/2$ and $\gamma = \ko+\rho+\eta$  for some small $\ep>0$. 
The condition $\ga+3\ko>1$ can then be read $\ga+3\ko=1/2+3\overline{\nu}/2-4\ep>1$,
which is possible whenever $\nu > 1/3$.

\end{proof}

\subsubsection{The fractional Brownian case}

In order to define an  integration theory for the fractional Brownian
motion beyond the Young case one has to start, like in the Brownian case,
by defining
the operators $X^1$ and $X^2$ in a natural way. We have already seen that
$X^1$ could be understood by  means of Wiener integrals, and for $X^2$, two
reasonable choices for the definition of (\ref{def:levy}) seem
to be the use of either Skorokhod or Stratonovich integrals with respect
to the fractional Brownian motion $X$. However, it turns out that
these two solutions
are equally unsatisfactory, for two different reasons that
we proceed to detail now:
\begin{enumerate}
\item
When one computes moments of random variables of
the second chaos defined by Stratonovich integrals, some trace
terms appear, a classical phenomenon which is explained for
instance in \cite{Nbk} in the general case,
in \cite{Ti} for the stochastic heat equation, or in \cite{PT} for the fBm.
In the current situation, if we want these trace terms to be convergent
for a fractional Brownian motion $X$ defined by (\ref{covstr}) and
(\ref{dcp:cov}), one has to choose $\nu>1/2$, which means
in particular that $Q$ is a bounded function of $\xi\in\ou$. In
other words, we are not allowed, even if $H>1/2$, to consider
a distribution-valued noise in space, which was one of our main aim.
\item
The Skorokhod integral works better as far as convergences and
regularity estimates are concerned. But one of the basic ingredients of our
algebraic manipulations on integrals is the fact that one can write,
under suitable hypothesis:
$$
\ist S_{t  u} dx_u\, b_s
=
\lc \ist S_{t  u} dx_u \rc b_s,
$$
an equality which is known to fail in the Skorokhod case
(see \cite{Nbk} again for further explanations). For instance,
the relation $\td X^2=X^1 X^1$, which is useful in our analysis,
does not hold true when $\hd X^2$ is understood as a Skorokhod integral.
\end{enumerate}
In order to cope with these problems, one can adopt the following
strategy: compute the correction term $P$, understood as a 2-increment
operator-valued process, which allows to write
\begin{equation}\label{m-n:correct}
\td X^2=X^1 X^1 +P,
\end{equation}
when $X^2$ is defined via Skorokhod integration. Notice that, since $X^2$
is an  element of $\ch_2$, the process $P$ is deterministic.

\vspace{0.2cm}

Recall now that the operator $\td$ has been defined as follows: for a suitable
Banach space $V$ and $M\in\hca_k^\infty(V)$,
set
\begin{equation}\label{def:td}
\td M=\hd M- Ma=\der M -aM -Ma.
\end{equation}
Then the operator $\td$ enjoys the same kind of properties as $\td$, and in particular, $\td\td=0$
and $\ker\td|_{\hca_3}=\mbox{Im}\td|_{\hca_2}$. Moreover, relation (\ref{m-n:correct}) can be read
as $\td P=0$. Thus, there exists another process $T\in\hca_2\cl$ such that $\td T=P$. Consider
then an explicit version of $T$ and set $\tilde X^2=X^2-T$. 
Then $\tilde X^2$ is still a Levy area type
process, such that $\td \tilde X^2= X^1 X^1$, which means that hopefully, 
$\tilde X^2$ will enjoy both
algebraic and analytic properties allowing a nice extension of the notion of convolution integral.
However, an open problem is to understand in which sense the integrals defined using this {\it
corrected} Levy area $\tilde X^2$ are useful and/or natural. We plan to report on this possibility
in a further paper.

\subsection{The algebra of a rough path}
Bilinear stochastic equations, in finite or infinite dimensions,
are often handled by means of chaos decomposition
(see e.g. \cite{LR,LS}). In this section, we will try to stress
some relationships  between our pathwise approach and this latter method.

\smallskip

Our basic  Hypothesis
\ref{hyp:nm} states that
\begin{equation}\label{hd:xuxd}
\td X^n= \sum_{k=1}^{n-1} X^k X^{n-k}
\end{equation}
for $n=1,2,3$ and moreover that
\begin{equation}\label{hyp:xuxd}
X^n\in\hcac_2^{\ga+(n-1)\ko}\cl^{\eta,-\rho}, \quad\mbox{ and }X^n\in\hcac_2^{n\ko}\cl^{\eta,\eta},
\end{equation}
with $\ga + 3\ko > 1$, $\eta>1/4$ and $\rho=1/4-\bar\nu/2$. Furthermore, it can be shown, along the same lines as for Theorem
\ref{th:conv-lambda}, that there exists an inverse $\tl$ to $\td$ on
$\hca_{3}^{\mu}\cl^{0,-\rho}\cap\ker(\td)$ for a certain $\mu>1$.

\vspace{0.2cm}

Let us see now how to construct an operator $X^4$ satisfying the operator Chen relation~(\ref{hd:xuxd}): by composition of operators, and invoking Hypothesis (\ref{hyp:xuxd}), it is easily checked
that
$
X^1 X^3 + X^3 X^1 + X^2 X^2 \in \hcac_3^{\ga+3\ko} \LL^{\eta,-\rho}
$.
Furthermore, we have assumed that $\ga+3\ko> 1$, and thus, by analogy with Theorem 
\ref{th:int-rough-lin}, we will set now
$
X^4 := \tl[X^1 X^3 + X^2X^2 + X^3 X^1]
$,
which is well defined as an element of $\ce_2^{\ga+3\ko} \LL^{\eta,-\rho}$ and thus that belongs to $\hcac_2^{4\ko}\LL^{\eta,\eta}$ (since $\ga = \ko+\eta+\rho$). It turns out that this procedure
can be iterated, and we obtain the following:
\begin{proposition}\label{th:rough-path-lifting}
Let $X$ satisfying Hypothesis
\ref{hyp:nm}. Then one can construct a sequence $\{X^n;n\ge 4\}$ out of $X^1,X^2,X^3$, such
that, for any $\ka<\ko=\ga_{n}-\rho$, we have
$
X^n \in \ce_2^{\ga+(n-1)\ko}\LL^{\eta,-\rho}
$,
$$
\|X^n\|_{\hcac_2^{n\ko}\LL^{\eta,-\rho}} \le C (n!)^{-\ko}
$$
and such that the operator Chen relations~(\ref{hd:xuxd}) are satisfied.
\end{proposition}
\begin{proof}
The $X^n$ are constructed by an induction on $n$. Then it is clear that 
$X^n \in \ce_2^{\ga+(n-1)\ko}\LL^{\eta,-\rho}$. Moreover, for $n\ge 4$ we have $n \ko > 1$, so that
$X^n = \tl ( \sum_{k=1}^{n-1} X^{n-k} X^k )$ can be defined
directly as an element of $\LL^{\eta,\eta}$. Then, the same kind of arguments as in the finite dimensional case  \cite{g-ramif} prove that we have the inductive bound
\begin{equation}\label{eq:bnd-factorial}
\|X^n\|_{\hcac_2^{n\ko}\LL^{\eta,\eta}} \le C_{X} (n!)^{-\ko}.
\end{equation}
\end{proof}

In such a setting, the lifted rough path allows to express the It\^o map which sends initial
conditions to solution of the linear equation~(\ref{eq:rough-ev}) by a convergent series of
operators:

\begin{corollary}
Under the conditions of Proposition \ref{th:rough-path-lifting}, there exists an operator $T$, defined as an element of $\ce_2^{\ga+(n-1)\ko}\LL^{\eta,-\rho}$, given by the strongly convergent series
$
T := \sum_{k=1}^\infty X^k
$,
and such that the solution of the linear problem~(\ref{eq:rough-ev})
satisfy the equation
$
\hd y = T y
$, or written in another way
$$
y_t = S_{t  s} y_s + T_{ts} y_s, \qquad (t,s) \in \cs^2.
$$
In particular if we define $X^0_{ts} = S_{t  s}$ and set $\hat T = X^0 +
T$ we have that $\hat T$ is a cocycle of operators:
$$
\hat T_{ts} = \hat T_{tu} \hat T_{us}, \qquad (t,u,s) \in \cs^3.
$$
\end{corollary}
\begin{proof}
The convergence of the series for $T$ in the operator norm follows from the bound (\ref{eq:bnd-factorial}) on $\|X^n\|$. The cocycle property is proven as in finite-dimension. The uniqueness of the solution to the linear problem allows to identify the operator $T$ as the It\^o map for the rough evolution equation. 

\end{proof}

\section{Polynomial non-linearities}
\label{sec:rough-nonlin}
Going back to the general setting explained at Sections \ref{sect:complex} and \ref{sec:circle},
we will consider now an equation of the form
\beq\label{eq:multi-lin}
y_t = S_t \psi + \int_0^t S_{t  u} dx_u M_n(y_u^{\otimes n})
\eeq
where $y$ lives in the Hilbert space $\cb$, where $M_n : \cb^n  \to \cb$ is
some unbounded multilinear operator from the Hilbert tensor $\cb^n =
\cb^{\otimes n}$ to $\cb$, and where we understand $\varphi^{\otimes n} =
\varphi\otimes \cdots \otimes \varphi \in \cb^n$ as the tensor monomial
generated by $\varphi \in \cb$. In fact, for sake of simplicity we assume that $M_n$ is
symmetric and we restrict our discussion to the case $n=2$, letting
$M_2 = B$, the general
situation posing no more conceptual difficulties. Then our general strategy, like in the linear case, will be first to expand equation (\ref{eq:multi-lin}) for a
smooth driving process $x$ in order to guess the appropriate rough-path 
underlying this equation. It will be seen that those expansions involve
some increments indexed by trees. Studying the algebraic and analytic properties
of these increments, we shall obtain a reasonable notion of solution to our 
quadratic equation.

\subsection{Formal expansions and trees}

Let us first simplify a little our setting.
Recall that we wish to solve an equation of the form
\beq\label{eq:quadratic}
y_t = S_t \psi + \int_0^t S_{t  u} dx_u B(y_u^{\otimes 2}),
\eeq
where we specialize our situation in the following way: assume first that
$\cb=L^2(\ou)$, which means that we are back again to the heat equation setting of Section \ref{sec:circle}.
Then $B:\cb^{\otimes 2}\to\cb$ is defined by $[B(\phi\otimes\psi)](\xi)=\phi(\xi)\psi(\xi)$
for $\xi\in\ou$, whenever this expression makes sense in $\cb$. Assume for the
moment that $x\in\hca_1^{1}\cl^{\ka,\ka}$ for $\ka$ large enough.
 We can expand 
equation (\ref{eq:quadratic}) as:
\begin{equation}\label{eq:expansion-multilin}
  \begin{split}
y_t & = S_{ts} y_s + \int_s^t S_{tu} dx_u B((S_{us }y_s)^{\otimes 2})
\\ & \quad
+  2 \int_s^t S_{tu} dx_u B( S_{us }y_s \otimes \int_s^u S_{uv} dx_v
B((S_{vs}y_s)^{\otimes 2}) )    
\\ & \quad
+  2 \int_s^t S_{tu} dx_u B( \int_s^u S_{uv} dx_v
B((S_{vs}y_s)^{\otimes 2}) \otimes \int_s^u S_{uv} dx_v
B((S_{vs}y_s)^{\otimes 2}) )    
\\ & \quad
+  4 \int_s^t S_{tu} dx_u B( S_{us }y_s \otimes \int_s^u S_{uv} dx_v
B((S_{vs}y_s) \otimes \int_s^v S_{vw} dx_w B((S_{ws }y_s)^{\otimes 2})
) )  
\\ & \quad + \text{h. o. iterated integrals}  
  \end{split}
\end{equation}
As we see, iterated integrals appear here in combinations which are not
as easy to handle as in the bilinear case of Section \ref{sec:lin-evolution}. A natural
way to code this kind of expansion is to use planar trees, as explained below.

\vspace{0.2cm}

Without entering too much into formal definitions involving trees, let us mention that
we shall consider planar binary rooted trees $\TT$ of the form:
$$
\Btrees
\btv, \btvva, \btvvb, \btvvv, \btvvva, \btvvvb, \btvvvc, \btvvvd, \text{etc\dots}
$$
which allow to 
give a compact expression of the iterated integrals appearing in the
expansion (\ref{eq:expansion-multilin}). 
Observe that each tree can be constructed from the trivial tree $\tau_0=\bullet$ by using the binary operation $V : \TT \times \TT \to \TT$ consisting in gluing two trees at a newly 
created root, so for example:
$$
\btvvv = V(V(t_0,t_0),V(t_0,t_0)).
$$
For any tree $\tau\in\TT$, we associate the function $d(\tau)$ that counts the number of leaves on the trees, so that  $d(\tau_0)=1$ and $d(V(\tau_1,\tau_2)) = d(\tau_1)+d(\tau_2)$.

\smallskip

Let us see now how to represent expansion (\ref{eq:expansion-multilin}) thanks
to planar trees. Define recursively an operator-valued increment  
$X^\tau \in \hcac_2 \LL(\cb^{d(\tau)}_{\eta}; \cb)$
for $\tau\in \TT$ as
\begin{equation}
\label{eq:X-def}
X^{\tau_0}_{ts} = S_{ts},
\quad\mbox{and}\quad
  X^{V(\tau_1,\tau_2)}_{ts} = \int_s^t S_{tu} dx_u
  B(X^{\tau_1}_{us}\otimes X^{\tau_2}_{us}).
\end{equation}
Notice that $X^{\tau}$ has always to be considered as an operator
acting on $\cb^{d(\tau)}_{\eta}$. For instance,
we understand  that, if $\tau=V(\tau_1,\tau_2)$, we have:
\begin{equation*}
  X^{V(\tau_1,\tau_2)}_{ts}(\varphi_1 \otimes \cdots \otimes \varphi_{d(\tau)}) = 
  \int_s^t S_{tu} dx_u
  B(X^{\tau_1}_{us}(\varphi_1 \otimes\cdots \otimes \varphi_{d(\tau_1)})\otimes X^{\tau_2}_{us}
(\varphi_{d(\tau_1)+1}\otimes \cdots \otimes \varphi_{d(\tau)})).
\end{equation*}
This latter formula justifies also in a sense the use of planar trees, since in general 
$X^{V(\tau_1, \tau_2)} \neq X^{V(\tau_2,\tau_1)} $. In order to illustrate this fact, consider
the simple example where $\tau_1=\tau_0$ and $\tau_2=V(\tau_0,\tau_0)$. Then
$d(V(\tau_1,\tau_2))=3$ and:
$$
 X^{V(\tau_1,\tau_2)}_{ts}(\varphi_1 \otimes \varphi_2 \otimes \varphi_{3})
 =
 \int_s^t S_{tu} dx_u
  B(X^{\tau_0}_{us}(\varphi_1)\otimes X^{V(\tau_0,\tau_0)}_{us}
(\varphi_2 \otimes\varphi_{3})),
$$
while
$$
 X^{V(\tau_2,\tau_1)}_{ts}(\varphi_1 \otimes \varphi_2 \otimes \varphi_{3})
 =
 \int_s^t S_{tu} dx_u
  B( X^{V(\tau_0,\tau_0)}_{us}
(\varphi_{1}\otimes  \varphi_{2}) \otimes X^{\tau_0}_{us}(\varphi_3) ),
$$
which are a-priori clearly different objects.

\smallskip

With these notations in hand, it is now  checked that our previous expansion
(\ref{eq:expansion-multilin}) can be written in a simpler way as:
\btrees
\begin{equation}
(\hd y)_{ts} =  X^{\btv}_{ts}(y^{\otimes 2}_s)
+ 2 X^{\btvva}_{ts}(y^{\otimes 3}_s)
+   X^{\btvvv}_{ts}(y^{\otimes 4}_s)
+  4 X^{\btvvva}_{ts}(y^{\otimes 4}_s)
+ r,
\label{eq:pol-lambda-eq-1}
\end{equation}
where $r \in \hcac_2(\cb)$ is some remainder term, and where we took care to distinguish the various operators obtained by permuting the factors in the $\cb$-tensors. Of course we could have expanded the solution further, and some operators associated to larger trees would have appeared.
However, in a smooth enough situation, the strategy in order to solve (\ref{eq:quadratic}) is now clear: we can use the map $\hl$ to eliminate the remainder from the equation:
\begin{equation}
\hd y =  (1-\hl \hd)[
 X^{\btv}(y^{\otimes 2})
+ 2 X^{\btvva}(y^{\otimes 3})
+   X^{\btvvv}(y^{\otimes 4})
+  4 X^{\btvvva}(y^{\otimes 4})]
\label{eq:pol-lambda-eq-2}
\end{equation}
and try to solve this by fixed-point method. The only condition we need to check is that 
\begin{equation}\label{eq:rel-trees-in-dom-Lambda}
\hd [ X^{\btv}(y^{\otimes 2})
+ 2 X^{\btvva}(y^{\otimes 3})
+   X^{\btvvv}(y^{\otimes 4})
+  4 X^{\btvvva}(y^{\otimes 4})
]
\end{equation}
should be in the domain of $\hl$,
 which means that its time-regularity should be greater than $1$.
The computation of expressions like (\ref{eq:rel-trees-in-dom-Lambda}) requires a little algebraic preparation.

\subsection{Algebraic computations}
To ease some computations we introduce an ``improper'' increment
$E_{ts} = \id$ (improper because is does not vanish as $t=s$), where the $\id$ has
to be understood, according to the context, as the identity operator on the vector
space under consideration. For example, we can write $\hd h = \der h - (S-E)h$. Moreover
we also introduce $e_{ts}=1$ taking values in $\RR$, so that for
example, if $z \in \hcac_2(\cb)$, then $(ze)_{tus} = z_{tu} e_{us} = z_{tu}
$.

\smallskip

It will also be useful to extend the action of $\hd$ to the tensors $\cb^n$ by letting
$\hd z = \der z - (\SS-E) z$,
where $\SS : \cb^n \to \cb^n$ is defined as $\SS = S \otimes \cdots
\otimes S$ for any $n \ge 1$. 
If the reader is unconfortable with
giving the same name at different operators, he can think that $\SS$
is defined on the direct sum $\oplus_{n \ge 1} \cb^n$; furthermore,  we will write explicitly 
$\SS_n$ when the context is insufficient to determine the actual space on which $\SS$
is defined. 
Analogously to the case of $\hd$, the operator
$\td $ defined by (\ref{def:td}) can be allowed to act on $\hcac_1 \LL(\cb^n,\cb)$ as
$\td H = \der H - (\SS-E)H - H(\SS-E)$.


\smallskip


We wish first to understand how the operators $\hd$ and $\td$ act 
on tensor products. More specifically,
we shall need three relations which are summarized in the following:
\begin{lemma}\label{lem:alg-relations-tensors}
The following relations hold true:

\noindent
{\it (1)} Let $z,w\in\cac_1(\cb)$. Then
\begin{equation}
  \label{eq:der-tens-product-C1}
\hd(z \otimes w) = \SS z\otimes \hd w +  \hd z \otimes \SS w + \hd z \otimes \hd w.
\end{equation}

\noindent
{\it (2)} Let $z,w\in\cac_2(\cb)$. Then
\begin{equation}
  \label{eq:der-tens-product}
\hd(z \otimes w) = z e \otimes \hd w + ze \otimes \SS w + \hd z
\otimes w e + \SS z \otimes w e + \SS z\otimes \hd w + \hd z \otimes \SS w
+ \hd z \otimes \hd w  
\end{equation}
Furthermore, it also holds that:
$$
\td (z\otimes z) =  Sz \otimes \td z + \td z \otimes S z + \td z \otimes \td z,
\quad
\hd (z\otimes z) =  Sz \otimes \hd z + \hd z \otimes S z + \hd z \otimes \hd z.
$$

\noindent
{\it (3)} Let  $Z \in \hcac_2 \LL(\cb^{a},\cb^b)$ and $W \in \hcac_2 \LL(\cb^c,\cb^d)$. Then
\begin{multline}
  \label{eq:der-tens-product-op}
\td(Z \otimes W) =   Z \SS \otimes \td W + Z\SS \otimes \SS W + \td Z
\otimes W \SS + \SS Z \otimes W\SS  + \SS Z \otimes \td W \\
+ \td Z \otimes \SS W
+ \td Z \otimes \td W, 
\end{multline}
where an example of notational convention is given by:
$$
(Z\SS \otimes \SS W)_{tus}= 
Z_{tu}\SS_{us} \otimes \SS_{tu}W_{us} \in \LL(\cb^{a+c},\cb^{c+d}).
$$ 
\end{lemma}

\begin{proof}
These relations are easily checked by elementary computations. We include the proof
of the third one for sake of completeness: notice that
$$
\td Z = Z^\diamond - Z E - E Z - (\SS -E) Z - Z (\SS-E) = Z^\diamond - Z \SS - \SS Z, 
$$
where $Z^\diamond_{tus} = Z_{ts}$, and thus
\begin{equation*}
\begin{split}
&\td (Z \otimes W)  = Z^\diamond \otimes W^\diamond -  \SS Z \otimes \SS W - Z\SS \otimes W\SS
\\ & =  (\SS Z + Z \SS + \td Z) \otimes (\SS W + W \SS + \td W) -  \SS Z \otimes \SS W - Z\SS \otimes W\SS,
\end{split}
\end{equation*}
which yields relation (\ref{eq:der-tens-product-op}) by a straightforward expansion.

\end{proof}


We also want to understand how $\hd,\td$ act on the operators $X^\tau$. 
A first relation in this direction is to note that, according to Lemma~\ref{daf},
if $X^\tau\in\hca_2\cl(\cb^{\otimes n};\cb)$ and $h\in\hca_1(\cb^{\otimes n})$,
\beq\label{eq:hd(X-tau-h)}
\hd [X^\tau h] = (\hd X^{\tau})h - X^{\tau} \der h,
\quad\mbox{ i.e. }\quad
\hd [X^\tau h] = (\td X^{\tau})h  - X^{\tau} \hd h.
\eeq
It is thus useful to compute quantities of the form $\td X^{\tau}$. To this purpose,
consider $n\ge 1$ and define  $I: \hcac_2 \LL(\cb^{n},\cb^2) \to \hcac_2 \LL(\cb^n,\cb)$ by
\beq\label{eq:def-map-I}
I(H)_{ts} = \int_s^t S_{tu} dx_u B H_{us}.
\eeq
This kind of expression can be related to our tree-indexed increments, by
noticing for instance that:
$$
X^{\btv} = I(\SS_2), \quad X^{\btvva} = I(X^{\btv} \otimes \SS_1),\quad
X^{\btvvb} = I(\SS_1 \otimes X^{\btv}), 
$$
and generally speaking, (\ref{eq:X-def}) can be read as:
\beq\label{eq:relation-I-map-X-V}
X^{V(\tau_1,\tau_2)}=I(X^{\tau_1}\otimes X^{\tau_2}),
\quad
X^{V(\tau_0,\tau)}=I(\cs\otimes X^\tau),
\quad\mbox{and}\quad
X^{V(\tau,\tau_0)}=I(X^\tau\otimes\cs).
\eeq
Hence, we shall compute differentials of terms of the form $I(H)$:
\begin{lemma}\label{lem:expression-td-I(H)}
Let $H\in\hcac_2^1 \LL(\cb^{n},\cb^2)$.
The following formula holds true for the derivative of $I(H)$:
$$
(\td I(H))_{tus}=I_{tu}(\cs_2) H_{us} + I_{tu}(\td(H)),
\quad\mbox{and}\quad
(\hd I(H))_{tus}=I_{tu}(\cs_2) H_{us} + I_{tu}(\hd(H)).
$$
Furthermore, if we assume that $\td H$ can be decomposed as
$(\td H)_{tus} = \sum_{j\le M} H^{(1,j)}_{tu} H^{(2,j)}_{us}$,
for a given $M\ge 1$, $H^{(1,j)}\in\hcac_2^1 \LL(\cb^{2},\cb^2)$, 
and $H^{(2,j)}\in\hcac_2^1 \LL(\cb^{n},\cb^2)$, then we obtain:
$$
\td I(H) = I(\SS_2) H + \sum_{j\le M} I(H^{(1,j)}) H^{(2,j)}.
$$
\end{lemma}

\begin{proof}
We have:
\begin{align*}
&[\td I(H)]_{tus}= I_{ts}(H)-I_{tu}(H)\cs_{us}-\cs_{tu}I_{us}(H)  \\
&=\int_s^t S_{tw} dx_w BH_{ws} -\int_u^t S_{tw} dx_w BH_{wu} S_{us}
-S_{tu} \int_s^u S_{uw} dx_w BH_{ws}  \\
&=\int_u^t S_{tw} dx_w BH_{ws} - \int_u^t S_{tw} dx_w BH_{wu} S_{us}  \\
&=\int_u^t S_{tw} dx_w B(\cs_{wu}H_{us})
+\int_u^t S_{tw} dx_w \lc BH_{ws}- B(\cs_{wu}H_{us})-BH_{wu} S_{us}\rc  \\
&= I_{tu}(\cs_2) H_{us} + I_{tu}(\td(BH)),
\end{align*}
which proves our first assertion. The second one is now trivially deduced.

\end{proof}

\smallskip

With these preliminaries in hand, we can now compute the action of $\td$ on 
the tree-indexed increments we have met so far, in the following way: 
\begin{lemma}\label{lem:expression-td-X-trees}
Let $x$ be a smooth $\cl(\cb)$-valued path. Then we have:
\begin{eqnarray*}
\td X^{\btv} &=& 0  ,\qquad 
\td X^{\btvva}=I(\SS_2) (X^{\btv}\otimes S)=X^{\btv}(X^{\btv}\otimes S)  \\
\td X^{\btvvv} &=&  X^{\btv} (X^{\btv}\otimes X^{\btv}) +  X^{\btvvb} (X^{\btv} \otimes S) + 
X^{\btvva}  (S \otimes X^{\btv}),
\end{eqnarray*}
and
\begin{eqnarray*}
\td X^{\btvvva} &=& X^{\btv} (X^{\btvva}\otimes S) +  X^{\btvva} (X^{\btv} \otimes \SS_2) \\
\td X^{\btvvvb} &=& X^{\btv} (X^{\btvvb}\otimes S) +  X^{\btvva} (\SS_2 \otimes X^{\btv} )  \\
\td X^{\btvvvc} &=& X^{\btv} (S\otimes X^{\btvva}) +  X^{\btvvb} (X^{\btv} \otimes \SS_2) \\
\td X^{\btvvvd} &=& X^{\btv} (S\otimes X^{\btvvb}) +  X^{\btvvb} (\SS_2 \otimes X^{\btv} ).
\end{eqnarray*}
\end{lemma}

\begin{proof}
All these relations are obtained by elementary computations, and we shall only sketch the proof
for some of them: first of all, invoking Lemma \ref{lem:expression-td-I(H)}, we get:
$$
\td X^{\btv} = \td I(\SS_2) = I(\SS_2) \SS_2 - I(\SS_2) \SS_2 = 0, 
$$
where we used the fact that $\td \SS  = -\SS \SS$. As far as $X^{\btvva}$ is concerned,
we have:
$$
\td X^{\btvva} = \td I(X^{\btv}\otimes S) = I(\SS_2) (X^{\btv}\otimes S),
$$
since $\td (X^{\btv}\otimes S) = 0$ by a direct computation using 
formula~(\ref{eq:der-tens-product-op}).
Similarly, it holds that:
$$
\td X^{\btvvv} = \td I(X^{\btv}\otimes X^{\btv}) = I(\SS_2) (X^{\btv}\otimes X^{\btv}) +  I(S \otimes X^{\btv}) (X^{\btv} \otimes S) +  I( X^{\btv}\otimes S) (S \otimes X^{\btv}),
$$
owing to the fact that
$$
\td (X^{\btv}\otimes X^{\btv}) 
= \SS X^{\btv}\otimes X^{\btv}\SS + X^{\btv}\SS \otimes \SS X^{\btv}.
$$
Now, invoking (\ref{eq:relation-I-map-X-V}), we have $I(S \otimes X^{\btv}) (X^{\btv} \otimes S)=X^{\btvvb}$ and $I( X^{\btv}\otimes S)=X^{\btvva}$, which yields 
$$
\td X^{\btvvv} =  X^{\btv} (X^{\btv}\otimes X^{\btv}) +  X^{\btvvb} (X^{\btv} \otimes S) + 
X^{\btvva}  (S \otimes X^{\btv}),
$$
as claimed in our lemma.


\end{proof}

\smallskip

It is important to note now that all the previous computations have been performed
for a smooth path $x$. However, we shall ask our driving process $x$ to satisfy the
following assumption:
\begin{hypothesis}\label{hyp:x-tree}
We
assume that the path $x$ allows to define some 
increments  $X^\tau$ for any $\tau \in \TT$ such
that $d(\tau) \le 4$. We also suppose that those increments satisfy the relations of Lemma \ref{lem:expression-td-X-trees}, and that the following H\"older regularities hold true: setting $|\tau|=d(\tau)-1$, we have
$X^{\tau} \in \hcac^{|\tau| \kappa_0} \LL(\BB_\eta^{d(\tau)},\BB_\eta)$ and 
$X^{\tau} \in \hcac^{\gamma+(|\tau|-1) \kappa_0} \LL(\BB_\eta^{d(\tau)},\BB_{-\rho})$,
with $\gamma + n\ko > 1$ and $\gamma = \ko + \eta +\rho$, for a given $\eta>1/4$.
\end{hypothesis}

\begin{remark}
Here again, it is important to work in spaces of the form $\BB_\eta$ with $\eta>1/4$. Indeed,
these spaces are algebras, which ensures at least that, whenever $\phi,\psi\in\BB_\eta$,
then $B(\phi,\psi)\in\cb$.
\end{remark}

\subsection{A space of integrable paths}

The general discussion of the bilinear equation requires a deep
understanding of the algebra of $X$. It is not the aim of this paper
to enter into this kind of considerations, and we prefer here to concentrate on a particular case
where $\kappa$ is sufficiently large to stop the expansions at some
low (but non-trivial) order. So here we assume that $\ga+3\kappa > 1$.

In order to solve the fixed-point problem associated to (\ref{eq:quadratic}), we introduce a new space of weakly controlled
paths, denoted by $\QQ_{X,\kappa}$, which enjoys some nice stability properties
under the map $\Gamma: y \mapsto z=S\psi + I(y\otimes y)$.
\begin{definition}
Let $\psi\in\cb_\ka$ be an initial condition, and $x$ a driving noise satisfying 
Hypothesis \ref{hyp:x-tree}, with $\ga+3\ka>1$.
We say that a path $y \in \hcac_1^{\ka}(\cb)\cap\cac_1^{b}(\cb_\ka)$ 
belongs to $\QQ_{X,\kappa}$ if $y_0=\psi$, and
$\hd y$ can be decomposed into:
\begin{equation}
  \label{eq:y-path}
 \hd y = X^{\btv}y^{\btv} + X^{\btvva} y^{\btvva}  + X^{\btvvb} y^{\btvvb} + y^{\sharp},
\end{equation}
where $y^{\btv},y^{\btvva}$ and $y^{\btvvb}$ can be written as:
$$
y^{\btv} = w \otimes w, \quad y^{\btvva} = w^{\btv} \otimes w, \quad
y^{\btvvb} = w\otimes w^{\btv},
$$
and the following regularities hold true:
$$
y \in \hcac_1^\kappa(\cb), \quad w,y \in \hcac_1^{\kappa_1}(\cb), \quad
w^{\btv} \in \hcac_1^{\kappa_2}(\cb^2), \quad y^\sharp \in \hcac_2^{3\kappa_2}(\cb),
$$
where $\ka>\ka_1>\ka_2$, $\ka-\ka_1=\ka_1-\ka_2\equiv \mu$ and $\ga+3\ka_2>1$.
On $\QQ_{X,\kappa}$, we define the semi-norm
$$
\Norm{y}{\QQ_{X,\kappa}} =
\Norm{y}{\hcac_1^{\ka}(\cb)}+
\Norm{w}{\hcac_1^{\ka_1}(\cb)} 
+\Norm{w^{\btv}}{\hcac_1^{\ka_2}(\cb^2)}
+\Norm{y^\sharp}{\hcac_1^{3\ka_2}(\cb)}.
$$
\end{definition}

Note that the constant path $y_t = S_t \psi$ is a controlled path whenever $\psi\in\cb_\ka$, 
and in this case $w,w^{\btv},y^\sharp$ are all identically zero.
Furthermore, the space $\QQ$ satisfies the following useful stability property:
\begin{theorem}
Assume that $x$ satisfies Hypothesis \ref{hyp:x-tree}, where we recall that $\ga+3\kappa > 1$
and $\ka_0>\ka$.
For $y\in\QQ_{\ka,X}$, define $z\equiv\Gamma(y)\in\QQ_{\ka,X}$ by $z_0=\psi$ and
a decomposition of the form:
$$
\hd z = X^{\btv}z^{\btv} + X^{\btvva} z^{\btvva}  + X^{\btvvb} z^{\btvvb} + z^{\sharp},
\quad\mbox{with}\quad
z^{\btv}=w_z\otimes w_z, \quad z^{\btvva} = w_z^{\btv}\otimes w_z,
\quad z^{\btvvb}=w_z\otimes w_z^{\btv},
$$
where $w_z=y$, $w_z^{\btv}=y^{\btv}$, and $z^\sharp \in \hcac^{3\kappa_0}(\cb)$ is
a remainder which can be written as 
$$
z^\sharp=
X^{\btvvv}(y^{\btv}\otimes y^{\btv})+
X^{\btvvva}(y^{\btvva}\otimes y)+ X^{\btvvvb}(y^{\btvvb}\otimes y)
+ X^{\btvvvc}(y\otimes y^{\btvva})+ X^{\btvvvd}(y\otimes y^{\btvvb})-\laa(J),
$$
where $J$ is defined by relation (\ref{eq:def-J}).
Then
\begin{enumerate}
\item
$\Gamma:\QQ_{\ka,X}\to\QQ_{\ka,X}$ is well defined.
\item
$\hd z$ coincides with $I(y\otimes y)$ in the smooth case.
\item
The following estimate holds true: for $0\le a< b \le T$, we have:
\beq\label{eq:norm-Gamma(y)-quadratic}
\Norm{z}{\QQ_{X,\kappa_0}([a,b])} \le C_{X,\psi} \lp 1+
(b-a)^{\mu} \cn^4[y;\,\QQ_{X,\kappa}([a,b])]\rp ,
\eeq
for a positive constant $C_X$ which only depends on the rough path $X$.
\end{enumerate}
\end{theorem}

\begin{proof}
Start from two smooth paths $x$ and $y$. If we apply the $I$ map defined at (\ref{eq:def-map-I})
we obtain, just as in the Young case (\ref{eq:split-1}):
\bean
\hd z = I(y \otimes y) &=& 
I(\hd(y\otimes y)+\SS_2(y\otimes y))
=I(\SS_2) (y\otimes y) + I(\hd(y\otimes y))  \\
&=&
I(\SS_2) (y\otimes y) + I(Sy \otimes \hd y)  + I(\hd y \otimes S y)  + I(\hd y \otimes \hd y),
\eean
where we have used Lemma \ref{lem:alg-relations-tensors}.
Expanding $\hd y$ in this equation and invoking relation (\ref{eq:relation-I-map-X-V}), 
we thus obtain:
\begin{equation*}
  \begin{split}
    \hd z & = X^{\btv}(y\otimes y) + X^{\btvva} (y^{\btv}\otimes y) +
    X^{\btvvb} (y\otimes y^{\btv}) 
\\ & \qquad+ X^{\btvvv}(y^{\btv}\otimes y^{\btv})+
X^{\btvvva}(y^{\btvva}\otimes y)+ X^{\btvvvb}(y^{\btvvb}\otimes y)
+ X^{\btvvvc}(y\otimes y^{\btvva})+ X^{\btvvvd}(y\otimes y^{\btvvb})+
z^{\flat},
  \end{split}
\end{equation*}
where $z^{\flat}$ has to be understood again as a remainder. Now, by our 
standard, argument we shall define $z^{\flat}$ in the non-smooth case by
$z^{\flat}=-\hl (J)$, where $J$ is given by:
\begin{equation}\label{eq:def-J}
  \begin{split}
     J & =  \hd [X^{\btv}(y\otimes y) + X^{\btvva} (y^{\btv}\otimes y) +
    X^{\btvvb} (y\otimes y^{\btv}) 
\\ & \qquad+ X^{\btvvv}(y^{\btv}\otimes y^{\btv})+
X^{\btvvva}(y^{\btvva}\otimes y)+ X^{\btvvvb}(y^{\btvvb}\otimes y)
+ X^{\btvvvc}(y\otimes y^{\btvva})+ X^{\btvvvd}(y\otimes y^{\btvvb})].
  \end{split}
\end{equation}

\smallskip

In order for this equation to be well defined we still need to check that
$J$  belongs to
$\hcac^{\ga+3\kappa}(\BB_{-\rho})$, which is in the domain of $\hl$ since $\ga+3\kappa>1$.
Let us then compute $J$: owing to (\ref{eq:hd(X-tau-h)}), we have
\begin{eqnarray*}
\hd  [ X^{\btv}(y\otimes y)] &=& [\td X^{\btv}](y\otimes y)-X^{\btv}[\hd(y\otimes y)]  \\
&=&
- X^{\btv}(S y\otimes \hd y) 
- X^{\btv}(\hd y\otimes S y)
- X^{\btv}(\hd y\otimes \hd y),
\end{eqnarray*}
thanks to relation (\ref{eq:der-tens-product-C1}) and Lemma \ref{lem:expression-td-X-trees}.
The other terms can be computed along the same lines, and here is a sample of what
is obtained:
\bean
\hd [ X^{\btvva} (y^{\btv}\otimes y)]&=&
X^{\btv}(X^{\btv} y^{\btv}\otimes S y) - X^{\btvva} (\hd y^{\btv}\otimes S y)
- X^{\btvva} (S y^{\btv}\otimes \hd y)
-X^{\btvva} (\hd y^{\btv}\otimes \hd y)  \\
\hd [ X^{\btvvva}(y^{\btvva}\otimes y)] &=& 
X^{\btv} (X^{\btvva} y^{\btvva}\otimes S y) 
+  X^{\btvva} (X^{\btv} \otimes \SS_2) (y^{\btvva}\otimes y)
- X^{\btvvva} \hd (y^{\btvva}\otimes y),
\eean
and
\begin{multline*}
    \hd [ X^{\btvvv}(y^{\btv}\otimes y^{\btv})] =  X^{\btv}
    (X^{\btv}y^{\btv} \otimes X^{\btv}y^{\btv}) + X^{\btvvb}
    (X^{\btv}y^{\btv} \otimes \SS_2 y^{\btv}) + X^{\btvva} (\SS_2
    y^{\btv} \otimes X^{\btv}y^{\btv}) \\  
    - X^{\btvvv}(\SS_2
    y^{\btv}\otimes \hd y^{\btv}) - X^{\btvvv}(\hd y^{\btv}\otimes \SS_2
    y^{\btv}) - X^{\btvvv}(\hd y^{\btv}\otimes \hd y^{\btv}).
\end{multline*}
Now, by gathering all the terms we have obtained, 
we obtain that $J=\sum_{k=1}^4J_k$, with:
\bean
J_1&=&
X^{\btv}\Big[ 
S y\otimes (- \hd y + X^{\btv}y^{\btv} + X^{\btvva} y^{\btvva}  + X^{\btvvb} y^{\btvvb}   ) 
+(- \hd y + X^{\btv}y^{\btv} + X^{\btvva}
y^{\btvva}  + X^{\btvvb} y^{\btvvb}   )\otimes Sy  \\
&&\hspace{6cm}- \hd y^{\btv} \otimes \hd y^{\btv}+ (X^{\btv}\otimes X^{\btv})(y^{\btv}\otimes y^{\btv})
\Big]  \\
J_2&=&
X^{\btvva}\Big[
- S y^{\btv} \otimes \hd y - \hd y^{\btv} \otimes Sy 
+ (\SS_2 \otimes X^{\btv})(y^{\btv}\otimes y^{\btv})
-\hd  y^{\btv} \otimes \hd y] \\
&&\hspace{6cm}+ (X^{\btv}\otimes \SS_2)(y^{\btvva} \otimes y)  + (S \otimes X^{\btv}\otimes S)(y^{\btvvb} \otimes y)\Big] 
\eean
and 
\bean
J_3&=&
X^{\btvvb}\Big[- S y \otimes \hd y^{\btv} 
- \hd y \otimes S y^{\btv} 
+ (X^{\btv} \otimes \SS_2)(y^{\btv}\otimes y^{\btv}) -\hd  y \otimes \hd y^{\btv}\\
&&\hspace{6cm}+ (\SS_2 \otimes X^{\btv})(y \otimes y^{\btvvb})  
+ (S \otimes X^{\btv}\otimes S)(y \otimes y^{\btvva})\Big]  \\
J_4&=&
- X^{\btvvv}\hd (y^{\btv} \otimes y^{\btv}) 
- X^{\btvvva}\hd(y^{\btvva}\otimes y)- X^{\btvvvb}\hd(y^{\btvvb}\otimes y)
- X^{\btvvvc}\hd(y\otimes y^{\btvva})-  X^{\btvvvd}\hd(y\otimes y^{\btvvb}).
\eean
Furthermore, notice that, using equation (\ref{eq:y-path}) for the increments of $y$, the quantity
$J_1$ can be simplified into:
$$
J_1=
X^{\btv}\Big[ 
-S y\otimes y^{\sharp}
-y^{\sharp}\otimes Sy  
- \hd y^{\btv} \otimes \hd y^{\btv}+ (X^{\btv}\otimes X^{\btv})(y^{\btv}\otimes y^{\btv})
\Big].
$$

\smallskip

We are now left with the cumbersome task which consists in analyzing the regularity
of all the terms we have produced so far. 
We shall just focus on one particular example, namely
$X^{\btvva}(\hd y^{\btv}\otimes Sy)$,
leaving the other ones to the patient reader. However, invoking again Lemma 
\ref{lem:alg-relations-tensors}, we have 
$$
X^{\btvva}(\hd y^{\btv}\otimes Sy)=
X^{\btvva}(\hd w\otimes \hd w\otimes Sy)+X^{\btvva}(S w\otimes \hd w\otimes Sy)
+X^{\btvva}(\hd w\otimes S w\otimes Sy).
$$
Among the three terms in the right hand side of this relation, we shall analyze
the first one, the other ones being similar: recall that
$X^{\btvva}\in\cac_2^{\ga+2\ka_0}\cl(\cb_\eta^3,\cb_{-\rho})$, 
$\hd w\in\cac_2^{\ka}(\cb_\eta)$ and $Sy\in\cac_2^0(\cb_\eta)$. Thus
$X^{\btvva}(\hd w\otimes \hd w\otimes Sy)\in\cac_3^{\ga+2\ka_0+\ka}(\cb_{-\rho})$,
which is enough regularity to apply the $\Lambda$-map. The other terms in the 
decomposition of $J$ can be treated similarly, which ends the proof of our first
assertion.

\smallskip

Our second claim being immediate from the construction of our integral, let us say a 
few words about the last one. Here again, many terms have to be estimated, and we
shall focus on a representative example, namely the term $w_z^{\btv}=y^{\btv}=w\otimes w$.
In fact the quantity $\cn[w\otimes w;\hca_1^{\ka_2}(\cb^2)]$ 
has to be estimated, and recall that,
according to Lemma \ref{lem:alg-relations-tensors}, the following relation holds true:
$$
\hd (w\otimes w) =  Sw \otimes \hd w + \hd w \otimes S w + \hd w \otimes \hd w.
$$
Thus, since $w\in\hca_1^{\ka_1}(\cb)$ and $S_t$ is a bounded operator on $\cb$
for any positive $t$, we obtain:
\bean
\cn[w\otimes w;\hca_1^{\ka_2}(\cb^2)]
&\le&  c \lp  1+(b-a)^{\mu}\, \cn^2[w;\, \hca_1^{\ka_1}(\cb^2) ]\rp  \\
&\le&  c \lp  1+(b-a)^{\mu}\, \cn^2[y;\cq_{X,\ka}([a,b])]\rp.
\eean
The other terms defining $\hd z$ can be treated along the same lines, which proves
relation (\ref{eq:norm-Gamma(y)-quadratic}).

\end{proof}

We can turn now to the main goal of this section, which is to
get an existence and uniqueness result for equation (\ref{eq:quadratic}):
\begin{theorem}\label{thm:lin-rough2}
Assume that $x$ allows to define some incremental operators $X^\tau$
for any $\tau\in\TT$ such that $d(\tau)\le 4$, and that these increments satisfy
Hypothesis \ref{hyp:x-tree}, for
$\ga,\ko,\ka,\eta$ such that $\ka < \ko < \ga$, $\ga+3\ka > 1$ and $\eta>1/4$. Then there exists a
strictly positive $T_0=T_0(X^\tau;\, d(\tau)\le 4)$ such that equation (\ref{eq:quadratic})
admits a unique solution $y\in\cq_{\ka,X}([0,T_0];\cb)$.
\end{theorem}

\begin{proof}
The proof of this result is very similar to those of Theorems \ref{th:young-spde} and \ref{thm:lin-rough},
and we shall omit the details here. Just notice that inequality (\ref{eq:norm-Gamma(y)-quadratic})
allows to construct an invariant ball for the map $\Gamma$ in $\cq_{\ka,X}([0,T_0];\cb)$, 
whenever $T_0$ is small enough. The contraction argument can then be written in a standard way.

\end{proof}


%%%%%%%%%%%%%%%%%%%%%%%%%%%%%%%%%%%%%%%%%%


\subsection{The Brownian case}\label{sec:brownian-case-polynomial}
In this section, we investigate the behavior of the operators $X^{\tau}$
 defined above, when $x=X$  is an infinite dimensional Brownian
motion, defined at Section \ref{sec:circle}. Our aim is of course to show that, under certain
conditions, $X$ satisfies Hypothesis \ref{hyp:x-tree}. To this purpose, for the remainder of the 
section, we will mainly consider some applications on the space 
$\BB_\eta$ for $ \eta = 1/4+\eps$ and a small $\eps > 0$. Let us also introduce an additional notation: for the remainder of the article, we will write $A \lesssim B$ for two real quantities $A$ and $B$ when $A\le c \, B$ for a universal constant $c$.
\begin{proposition}\label{prop:reg-M-brownian}
Let $X$ be an infinite dimensional Brownian motion defined by the
covariance structure (\ref{cov:bm}), with $Q$ given by (\ref{dcp:cov})
for $\nu\in(1/3,1/2]$. Recall also that, for a planar binary tree, we have set $|\tau|=d(\tau)-1$.
Let $X^{\tau}$  be the incremental operator
given by~(\ref{eq:X-def})
where the stochastic integrals have to be understood in the It\^o sense.
Then, almost surely,
$X^{\tau}\in\hcac_2^{\ko |\tau|}\cl_{\mbox{\tiny HS}}(\cb^{d(\tau)}_{\eta};\cb_{\eta})$ 
for 
$
\tau = \btv, \btvva, \btvvb, \btvvv, \btvvva, \btvvvb, \btvvvc, \btvvvd
$
and 
for any $\ko$ satisfying  $0<\ko<1/4-\eta+\overline\nu/2$. Moreover
$X^{\tau}\in\hcac_2^{\ga+\ko (|\tau|-1)-1/4}\cl_{\mbox{\tiny HS}}(\cb^{d(\tau)}_{\eta};\BB_{-\rho})$ 
for $\ga = \ko+\eta+\rho<1/2$. Theorem \ref{thm:lin-rough2} can then be applied in this situation.
\end{proposition}

\begin{proof}
In all the cases the line of the proof is the same: we obtain an $L^2$ estimate
on the Hilbert-Schmidt norm of $X^\tau$, which by Gaussian tools can
be boosted to an $L^p$ bound for any $p$. Applying Lemma~\ref{lemma:besov}, the result is then easily deduced. 

\smallskip

Admitting for the moment the results of Lemma~\ref{lemma:ltwo-bounds} below, let us give
some details about our method: since our incremental processes always belong to a finite chaos 
of the infinite-dimensional Brownian motion $X$, it is easily deduced from 
Lemma~\ref{lemma:ltwo-bounds} that
$$
\expect \lc\|X^\tau_{ts}\|_{HS,\LL(\cb^{d(\tau)}_{\eta};\cb_{\eta})}^p\rc \lesssim (t-s)^{p|\tau|(\kappa_0 +\eps +(1-1/|\tau|)/2)}
 \lesssim (t-s)^{p|\tau|(\kappa_0 +\eps)}
$$
for  any  $0<\ka_0<1/4-\eta+\overline\nu/2< \eta/2$. Moreover, it also holds that:
$$
\|\td X^\tau_{tus}\| \le \sum_{\tau^1,\tau^2} \|X^{\tau^1}_{tu}\| \|X^{\tau^1}_{us}\| \lesssim
 (t-u)^{|\tau^2|\kappa_0} (u-s)^{|\tau^1|\kappa_0} \lesssim  (t-s)^{|\tau|\kappa_0},
$$
where $\tau^1,\tau^2$ denote the trees appearing in the expansion for $\td X^\tau_{tus}$ given
at Lemma \ref{lem:expression-td-X-trees}, for which we have always $|\tau^1|,|\tau^2|\ge 1$
and $|\tau^1|+|\tau^2|=|\tau|+1$. So it is clear that using the extended G-R-R Lemma \ref{lemma:besov}, we obtain
$$
\|X^\tau_{tus}\| \lesssim (t-s)^{|\tau|\kappa_0},
$$
for any $\tau$ such that $d(\tau) \le 4$. Finally, observe that the conditions $\ga<1/2$, $\ko=\bar\nu/2-\ep$ and $3\ko=\ga>1$ force us to choose $\nu>1/3$.

\end{proof}

An easy consequence of the last estimations is an existence result for a Brownian SPDE in the rough-path sense:

\begin{theorem}
Let $X$ be an infinite dimensional Brownian motion on $\ott\times\ou$,
defined by the covariance function given by (\ref{cov:bm}) and
(\ref{dcp:cov}) with $\nu > 1/3$. Then there exists $\eta>1/4$,  $0<\ka<\gamma<1/2$ satisfying $\ka<\ko$ and $\ga+3\ka>1$ such that, for any $\psi \in \cb_{\eta}$ the equation
\begin{equation*}
Y(0,\xi) = \psi(\xi), \qquad \partial_t Y(t,\xi)
= \Delta Y(t,\xi) dt + Y(t,\xi)^2 X(dt,d\xi),
\quad
t\in\ott, \xi\in\ou,
\end{equation*}
with periodic boundary conditions, understood as equation (\ref{eq:rough-ev}),
has a unique \emph{local} solution in $\cq_{\ka,\eta,\psi}$ up to a time $T_*$ which depends on the initial condition and on the operators $X^\tau$, $|\tau|\le 3$.
\end{theorem}

\begin{proof}
Like in the proof of Theorem \ref{th:brown-bilin} the proof amounts to check the validity of Hypothesis~\ref{hyp:x-tree} in the light of Prop.~\ref{prop:reg-M-brownian}.
\end{proof}


\bigskip

The rest of the paper is dedicated to the $L^2$ estimations for the operators $X^\tau$. In fact we will obtain 
a slightly stronger result than the one we claimed at Proposition 
\ref{prop:reg-M-brownian}.
\begin{lemma}
\label{lemma:ltwo-bounds}
For the trees considered at Proposition 
\ref{prop:reg-M-brownian}, we have the following $L^2$ bounds:
$$
E\lc \|X^\tau_{ts}\|_{HS,\LL(\cb^{d(\tau)}_{\eta};\cb_{\eta})}^2\rc \lesssim (t-s)^{|\tau|\Delta-1/2},
$$
where $\Delta = 1-2\eta+\nu-\eps$ for some arbitrary small $\eps>0$.
\end{lemma}
\begin{proof}
 It is conceptually easy to generalize the arguments of
Prop.~\ref{reg:m-brownian} to reduce the problem to an estimation of a
mixed sum (over eigenvalues of $A_o$) and integral over time variables
(after contraction of the stochastic integrals). This long and tedious task is left to the reader.
We prefer to give a diagrammatical algorithm which allows 
to go from the  kernel on $L^2$  (associated to each operator) to a simple sum estimation. 
This will be detailed in the next two subsections.

\end{proof}

\subsection{Diagrammatica}\label{sec:diagrammatica}
We will first show, at a heuristic level and on a simple example, how to pass from an
incremental operator to a graph for the computation of Hilbert-Schmidt norms.

\smallskip

\noindent
{\bf (1)} {\it Case of the operator} $X^{\btv}.$
Consider first the operator $X^{\btv}:\cb_\eta^2\to\cb_{\eta}$. 
Recalling the notations of Section \ref{sec:circle}, an orthonormal basis for 
$\cb_\eta$ is given by $\{\tilde e_j;\, j\ge 1\}$, where $\tilde e_j=\la_j^{-\eta}e_j$.
Furthermore, the particular form (\ref{dcp:cov}) we have assumed on the covariance
function $Q$ implies that $x$ can be decomposed as
$$
x_u=\sum_{p\in\Z} \la_p^{-\nu/2} e_p \, \beta_u^p,
\quad u\ge 0,
$$
where $\{\beta^p;\, p\in \Z\}$ is a sequence of independent Brownian motions.
Hence, setting $\langle\cdot ,\cdot\rangle$ for the inner product in $L^2(S)$,
the matrix elements of $X^{\btv}$ are given by:
\begin{align*}
&[X_{ts}^{\btv}]_{i,jk} =\left\langle e_i , X_{ts}^{\btv}( \tilde e_j\otimes  \tilde e_k)\right\rangle
= \left\langle e_i ;\,  \ist S_{tu} dx_u  (S_{us} \tilde e_j) (S_{us}  \tilde e_k) \right\rangle  \\
&= \ist \left\langle S_{tu} e_i ;\,   dx_u  (S_{us} \tilde e_j) (S_{us}  \tilde e_k) \right\rangle 
= \sum_{p\in\Z}\la_j^{-\eta}\la_k^{-\eta}\la_p^{-\nu/2}\ist d\beta^p_u
\langle   S_{tu} e_i,  e_p  (S_{us} e_j) (S_{us}  e_k) \rangle \\
&=\sum_{p;\, i=p+j+k}\la_j^{-\eta}\la_k^{-\eta}\la_p^{-\nu/2} 
 \ist  e^{-\lambda_i(t-u)-\lambda_j(u-s)-\lambda_k(u-s)} \, d\beta^p_u,
\end{align*}
Thus, the 
Hilbert-Schmidt norm of $X_{ts}^{\btv}$ in $\LL(\cb^2;\cb_{\eta})$ can be written as
$$
\|X_{ts}^{\btv}\|_{HS,\LL(\cb_{\eta}^2;\cb_{\eta})}^2 = \sum_{i,j,k\in\Z}
 \la_i^{2\eta} \, |[X_{ts}^{\btv}]_{i,jk}|^2.
$$
From this simple computation, the following rules pop out:
\begin{itemize}
\item
Some multiple sums (involving terms of the form $\la_i^\al$) with constraints on the indices appear, 
due to the fact that $\{e_i;\, i\in \Z\}$ is the trigonometric basis of $L^2(S)$.
\item
Some contractions in the sums take place, because of the Brownian stochastic integrals.
\end{itemize}
With the above considerations in mind, we 
can associate to  $X_{ts}^{\btv}$ the following graphical representation:
$$
X_{ts}^{\btv} = \int_s^t d\beta^p_u  \,\,{\parbox[c][40mm]{30mm}{
 \begin{fmfgraph*}(60,60)\setlength{\unitlength}{1cm}
\fmfstraight
   \fmftop{i1} \fmfbottom{o1,o2}
\fmfright{r3,r2,r1}
\fmflabel{$i$}{i1}
\fmflabel{$j$}{o1}
\fmflabel{$k$}{o2}
\fmflabel{$p$}{r1}
\fmflabel{$u$}{v1}
   \fmf{plain,tension=2.0}{v1,i1}
   \fmf{plain}{o2,v1}
   \fmf{plain}{o1,v1}
\fmffreeze
\fmf{dashes}{v1,r1}
   \fmfdot{v1}
 \end{fmfgraph*}
}}
$$
where the solid lines represent factors of $S$. This is a bookkeeping
device for the relation between the various indices and time
parameters. The computation of $\expect
[\|X_{ts}^{\btv}\|_{\HS,\LL(\cb_{\eta}^2;\cb_{\eta})}^2]$ corresponds to putting side
by side two specular copies of this graph and connecting the
corresponding top and bottom lines (to compute the HS norm), while contracting in all the
allowed ways the dashed lines (to compute the contractions of the
stochastic integrals). Doing so we obtain the graph
\fmfcmd{%
  vardef bar (expr p, len, ang) =
   ((-len/2,0)--(len/2,0))
      rotated (ang + angle
        direction length(p)/2 of p)
      shifted point length(p)/2 of p
  enddef;
  style_def tfermion expr p =
   draw_plain p;
   ccutdraw bar (p, 4mm,  60);
   ccutdraw bar (p, 4mm, -60)
  enddef;
  style_def tfermioni expr p =
   draw_double p;
   ccutdraw bar (p, 4mm,  60);
   ccutdraw bar (p, 4mm, -60)
  enddef;}
$$
\expect
\|X_{ts}^{\btv}\|_{\HS,\LL(\cb_{\eta}^2;\cb_{\eta})}^2 = \int_s^t du
    \parbox[c][30mm]{30mm}{%\fbox
    {
 \begin{fmfgraph*}(100,100)
\setlength{\unitlength}{1cm}
   \fmftop{i1} \fmfbottom{o1}
  \fmf{phantom,tension=5.5}{v1,i1}
   \fmf{phantom,tension=5.5}{v2,o1}
   \fmf{tfermioni,left=1.5}{v1,v2}
   \fmf{tfermioni,left=0.5}{v1,v2}
   \fmf{dashes,right=0.5}{v1,v2}
   \fmf{tfermion,right=1.5}{v1,v2}
    \fmfdot{v1,v2}
\fmflabel{$u$}{v1}
\fmflabel{$u$}{v2}
 \end{fmfgraph*}
}}
$$
where we use the following convention: solid lines correspond to
factors of $S$, time parameters are attached to vertices, crossed
solid lines correspond to factors of $A^{2\eta}S$ (coming form
contraction of output lines), crossed double lines correspond to factors of $A^{-2\eta}S$ (coming form
contraction of input lines), dashed lines are associated to factors of
$Q$ (coming from the It\^o contraction of the noise). We fix an
orientation for each edge of the graph and associate an index to each
oriented edge. To each vertex corresponds a constraint that the sum of
indexes of incoming edges minus indexes of outgoing edges should be
zero.
According to these rules the formula of the mean squared norm is then
$$
\expect
\lc\|X_{ts}^{\btv}\|_{\HS ,\LL(\cb_{\eta}^2;\cb_{\eta})}^2\rc = \int_s^t du  \sum_{i+j+k+l=0} \lambda_i^{2\eta}  \lambda_j^{-\nu} \lambda_k^{-2\eta}
\lambda_l^{-2\eta}  e^{-2\lambda_i(t-u)-2\lambda_k(u-s)-2\lambda_l(u-s)},
$$
and the reader can check that this is indeed the expression for the
mean value of the HS norm of $X_{ts}^{\btv}$.

\smallskip

Consider now the expression
$$
A=\sum_{i+j+k+l=0} \lambda_i^{2\eta}  \lambda_j^{-\nu} \lambda_k^{-2\eta}
\lambda_l^{-2\eta}  e^{-2\lambda_i(t-u)-2\lambda_k(u-s)-2\lambda_l(u-s)}.
$$
We trivially have
$$
A\le \sum_{i+j+k+l=0} \lambda_i^{2\eta}  \lambda_j^{-\nu} \lambda_k^{-2\eta}
\lambda_l^{-2\eta}  e^{-2\lambda_i(t-u)}.
$$
Furthermore, setting $q=k+l$, applying Lemma~\ref{lemma:asym-sum}, and recalling
that we have chosen $\eta>1/4$, we can bound $A$ as:
$$
A
\lesssim \sum_{i+j+q=0} \lambda_i^{2\eta}  \lambda_j^{-\nu} \lambda_q^{-2\eta} 
e^{-2\lambda_i(t-u)},
$$
where we used the relation
$
\sum_{k+l=q} \lambda_k^{-2\eta} \lambda_l^{-2\eta}\lesssim \lambda_q^{-2\eta}.
$
Moreover, assuming that $\nu \le 2\eta$ we can use again Lemma~\ref{lemma:asym-sum} to get
$$
A
\lesssim \sum_{i} \lambda_i^{2\eta-\nu}  e^{-2\lambda_i(t-u)} \lesssim (t-u)^{-a} \sum_i  \lambda_i^{2\eta-\nu-a},
$$
and choosing $a=1/2+2\eta-\nu+\eps$ so that the sum is convergent, we obtain
$\lesssim (t-u)^{\Delta-3/2}$, with $\Delta=1-2\eta+\nu-\ep$.
So we proved the graphical equation
\begin{equation}
\label{eq:graph-1}
  \begin{split}
     %
    \parbox[c][30mm]{30mm}{%\fbox
    {
 \begin{fmfgraph*}(100,100)
\setlength{\unitlength}{1cm}
   \fmftop{i1} \fmfbottom{o1}
  \fmf{phantom,tension=5.5}{v1,i1}
   \fmf{phantom,tension=5.5}{v2,o1}
   \fmf{tfermioni,left=1.5}{v1,v2}
   \fmf{tfermioni,left=0.5}{v1,v2}
   \fmf{dashes,right=0.5}{v1,v2}
   \fmf{tfermion,right=1.5}{v1,v2}
    \fmfdot{v1,v2}
\fmflabel{$u$}{v1}
\fmflabel{$u$}{v2}
 \end{fmfgraph*}
}}
\lesssim
(t-u)^{\Delta-3/2}
\end{split}
\end{equation}
and hence, if we suppose that $\Delta-3/2>-1$, i.e $\Delta>1/2$, we have obtained that
$$
\expect
\lc \|X_{ts}^{\btv}\|_{\HS ,\LL(\cb_{\eta}^2;\cb_{\eta})}^2\rc \lesssim \int_s^t du  (t-u)^{\Delta-3/2} \lesssim (t-s)^{\Delta-1/2},
$$
which is the desired bound for Lemma \ref{lemma:ltwo-bounds}. 
Let us say a few words about the condition $\Delta>1/2$: if $\eta=1/4+\hep$, then
one has $\Delta=1/2+\nu-\ep+\hep$. This means that the condition $\Delta>1/2$
can be met as soon as $\nu>0$, which simply rules out the possibility of considering
a space-time white noise at this stage.

\smallskip

\noindent
{\bf (2)} {\it Case of the operator} $X^{\btvva}.$
With the same kind of considerations as for $X^{\btv}$, it can be shown that the 
the matrix elements of the operator $X^{\btvva}$ are given by:
\begin{multline*}
[X_{ts}^{\btvva}]_{i,jkl} = \sum_{p+j+n=i} \sum_{q+k+l=n}
\lambda_j^{-2\eta}\lambda_k^{-2\eta}\lambda_l^{-2\eta}
\lambda_p^{-\nu/2} \lambda_{q}^{-\nu/2} \\
\times \int_s^t e^{-\lambda_i(t-u)} d\beta^p_u
e^{-\lambda_j(u-s)} \int_s^u e^{-\lambda_n(u-v)} d\beta^{q}_v e^{-\lambda_k(v-s)}  e^{-\lambda_l(v-s)}.
\end{multline*}
Thus, its
Hilbert-Schmidt norm in $\LL(\cb_{\eta}^3;\cb_{\eta})$ can be written as
$$
\|X_{ts}^{\btvva}\|_{\HS ,\LL(\cb_{\eta}^3;\cb_{\eta})}^2 = \sum_{i,j,k,l}
\lambda_i^{2\eta} |[X_{ts}^{\btvva}]_{i,jkl}|^2,
$$
and the following graphical representation can be associated to this last expression: 
$$
X_{ts}^{\btvva} = \int_s^t d\beta^p_u \int_s^u d\beta^q_v \,\,{\parbox[c][40mm]{30mm}{
 \begin{fmfgraph*}(60,60)\setlength{\unitlength}{1cm}
\fmfstraight
   \fmftop{i1} \fmfbottom{o1,o2,o3}
\fmfright{r3,r2,r1}
\fmflabel{$i$}{i1}
\fmflabel{$j$}{o1}
\fmflabel{$k$}{o2}
\fmflabel{$l$}{o3}
\fmflabel{$p$}{r1}
\fmflabel{$q$}{r2}
\fmflabel{$u$}{v1}
\fmflabel{$v$}{v2}
   \fmf{plain,tension=2.0}{v1,i1}
   \fmf{plain}{o3,v1}
   \fmf{phantom}{o1,v1}
\fmffreeze
\fmf{dashes}{v1,r1}
\fmf{dashes}{v2,r2}
   \fmf{plain}{o1,v2}
   \fmf{plain}{o2,v2}
   \fmf{plain}{v2,v1}
   \fmfdot{v1,v2}
%\fmflabel{$S$}{}
 \end{fmfgraph*}
}}
$$
Thus, for the Hilbert-Schmidt norm of $X_{ts}^{\btvva}$, we obtain the graph
%\newcommand{\ml}[1]{$\scriptstyle #1$}
$$ \expect \lc\|X_{ts}^{\btvva}\|_{\HS ,\LL(\cb_{\eta}^3;\cb_{\eta})}^2\rc 
= \int_s^t du  \int_s^u dv
%\fbox
{\parbox[c][30mm]{50mm}{\hskip -1cm %\fbox
{
 \begin{fmfgraph*}(200,250)
\setlength{\unitlength}{1cm}
   \fmftop{i1,i2} \fmfbottom{o1,o2}
\fmflabel{$u$}{v1}
\fmflabel{$u$}{v2}
\fmflabel{$v$}{v3}
\fmflabel{$v$}{v4}
   \fmf{phantom,tension=1.5}{v1,i1}
   \fmf{phantom,tension=1.5}{v2,o1}
   \fmf{tfermion,right=1}{v1,v2}
   \fmf{dashes}{v1,v2}
   \fmf{tfermioni,left=1}{v1,v2}
   \fmfdot{v1,v2}
   \fmf{phantom,tension=1.5}{v3,i2}
   \fmf{phantom,tension=1.5}{v4,o2}
   \fmf{tfermioni,left=1}{v3,v4}
   \fmf{dashes}{v3,v4}
   \fmf{tfermioni,right=1}{v3,v4}
   \fmf{plain,left=0.2}{v1,v3}
   \fmf{plain,right=0.2}{v2,v4}
   \fmfdot{v3,v4}
%\fmflabel{$S$}{}
 \end{fmfgraph*}
}}}
$$
and the corresponding formula:
\begin{multline*}
 \expect \lc\|X_{ts}^{\btvva}\|_{\HS ,\LL(\cb_{\eta}^3;\cb_{\eta})}^2\rc
 =
\int_s^t du \int_s^u dv \sum_{i+j+k+l=0} \lambda_i^{2\eta}  \lambda_j^{-\nu} \lambda_k^{-2\eta}  \\
\times \sum_{n+m+o=l} \lambda_m^{-\nu} \lambda_n^{-2\eta} \lambda_o^{-2\eta} e^{-2\lambda_i(t-u)-2\lambda_l(u-v)-2\lambda_n(v-s)-2\lambda_o(v-s)-2\lambda_k(u-s)}.
\end{multline*}

\smallskip

The strategy to control this expression is now straightforward: bounding the exponential and performing the time integrations gives, for two positive constants $a,b$ such that
$a+b<2$,
$$
\expect \lc\|X_{ts}^{\btvva}\|_{\HS ,\LL(\cb_{\eta}^3;\cb_{\eta})}^2\rc
\lesssim (t-s)^{2-(a+b)}\sum_{i+j+k+l=0} \lambda_i^{2\eta-a} \lambda_j^{-\nu} \lambda_k^{-2\eta}
\lambda_l^{-b} \sum_{n+m+o=l} \lambda_m^{-\nu} \lambda_n^{-2\eta} \lambda_o^{-2\eta}.
$$
Using the fact that $2\eta > 1/2$  we can reduce this to the  bound
$$
\expect \lc\|X_{ts}^{\btvva}\|_{\HS ,\LL(\cb_{\eta}^3;\cb_{\eta})}^2\rc
\lesssim (t-s)^{2-(a+b)}\sum_{i+j+l+k=0} \lambda_i^{2\eta-a} \lambda_j^{-\nu} \lambda_k^{-2\eta}
\lambda_l^{-b} \sum_{n+m=l} \lambda_m^{-\nu} \lambda_n^{-2\eta}.
$$
Assuming moreover that $\nu \le 2\eta$ we get
$$
\expect \lc\|X_{ts}^{\btvva}\|_{\HS ,\LL(\cb_{\eta}^3;\cb_{\eta})}^2\rc
\lesssim (t-s)^{2-(a+b)}\sum_{i+j+l+k=0} \lambda_i^{2\eta-a} \lambda_j^{-\nu} \lambda_k^{-2\eta}
\lambda_l^{-\nu -b}.
$$
At this point choose $b = 2\eta -\nu$, so that $\nu+b=2\eta$ and
$$
\expect \lc\|X_{ts}^{\btvva}\|_{\HS ,\LL(\cb_{\eta}^3;\cb_{\eta})}^2\rc
\lesssim (t-s)^{2-(a+b)}\sum_{i+k+l=0} \lambda_i^{2\eta-a} \lambda_l^{2\eta} \lambda_k^{-2\eta}.
$$
Hence this sum is finite if we choose $a=1/2+2\eta-\nu +\eps$, and we get
$$
\expect \lc\|X_{ts}^{\btvva}\|_{\HS ,\LL(\cb_{\eta}^3;\cb_{\eta})}^2\rc \lesssim (t-s)^{2\Delta-1/2}.
$$

\smallskip

Before proceeding to the estimation of the other more complex operators, let us make a useful observation.
Consider the following subgraph on the left of the previous graph:
\begin{equation*}
  \begin{split}
\int_s^u dv
     %
    \parbox[c][30mm]{30mm}{%\fbox
    {
 \begin{fmfgraph*}(150,250)
\setlength{\unitlength}{1cm}
   \fmftop{i2,i3} \fmfbottom{o2,o3}
  \fmf{phantom,tension=1.5}{v3,i3}
   \fmf{phantom,tension=1.5}{v4,o3}
   \fmf{tfermioni,left=1}{v3,v4}
   \fmf{dashes}{v3,v4}
   \fmf{tfermioni,right=1}{v3,v4}
   \fmf{plain,left=0.2,label=$i$}{v1,v3}
   \fmf{plain,right=0.2,label=$i$}{v2,v4}
   \fmfdot{v3,v4}
\fmflabel{$u$}{v1}
\fmflabel{$u$}{v2}
\fmflabel{$v$}{v3}
\fmflabel{$v$}{v4}
   \fmf{phantom,tension=1.5}{v1,i2}
   \fmf{phantom,tension=1.5}{v2,o2}
   \fmf{phantom,right=.5}{v1,v2}
   \fmf{phantom,left=.5}{v1,v2}
   \fmf{phantom}{v1,v2}
   \fmfdot{v1,v2}
 \end{fmfgraph*}
}}
\end{split}
\end{equation*}
After reduction of the crossed double lines (carrying the factors due to $A^{-\eta}$) 
by an iterated application of Lemma \ref{lemma:asym-sum}, we obtain the following expression which corresponds to a bound for this graph:
$$
\lesssim \int_s^u dv \lambda_i^{-\nu} e^{-2\lambda_i(u-v)} \lesssim  \int_s^u \frac{dv}{(u-v)^b} \lambda_i^{-\nu-b},
$$
so that choosing $b=2\eta-\nu$, we get an estimate of the form
$\lesssim   (v-s)^{\Delta} \lambda_i^{-2\eta}$. Summarizing these considerations,
we have obtained the graphical equation:
\begin{equation}\label{eq:graph-ineq}
  \begin{split}
\int_s^u dv
     %
    \parbox[c][30mm]{60mm}{%\fbox
    {
 \begin{fmfgraph*}(150,250)
\setlength{\unitlength}{1cm}
   \fmftop{i2,i3} \fmfbottom{o2,o3}
  \fmf{phantom,tension=1.5}{v3,i3}
   \fmf{phantom,tension=1.5}{v4,o3}
   \fmf{tfermioni,left=1}{v3,v4}
   \fmf{dashes}{v3,v4}
   \fmf{tfermioni,right=1}{v3,v4}
   \fmf{plain,left=0.2}{v1,v3}
   \fmf{plain,right=0.2}{v2,v4}
   \fmfdot{v3,v4}
\fmflabel{$u$}{v1}
\fmflabel{$u$}{v2}
\fmflabel{$v$}{v3}
\fmflabel{$v$}{v4}
   \fmf{phantom,tension=1.5}{v1,i2}
   \fmf{phantom,tension=1.5}{v2,o2}
   \fmf{phantom,right=.5}{v1,v2}
   \fmf{phantom,left=.5}{v1,v2}
   \fmf{phantom}{v1,v2}
   \fmfdot{v1,v2}
 \end{fmfgraph*}
}}
\lesssim
(u-s)^{\Delta}
    \parbox[c][30mm]{30mm}{%\fbox
    {
 \begin{fmfgraph*}(50,50)
\setlength{\unitlength}{1cm}
   \fmftop{i1} \fmfbottom{o1}
   \fmf{tfermioni,left=0.5}{i1,o1}
   \fmfdot{i1,o1}
\fmflabel{$u$}{o1}
\fmflabel{$u$}{i1}
 \end{fmfgraph*}
}}
\end{split}
\end{equation}
which we will use multiple times below. 


\subsection{More complex graphs}\label{sec:complex-graphs}
The tools we have introduce so far will allow us to treat the two remaining cases
we are left with, namely $X^{\btvvva}$ and $X^{\btvvv}$.

\smallskip

\noindent
{\bf (1)} {\it Case of the operator} $X^{\btvvva}.$
By using the same kind of arguments as in the previous subsection, we obtain a representation of the form:
$$
X_{ts}^{\btvvva} = \int_s^t d\beta^p_u \int_s^u d\beta^q_v \int_s^v d\beta^r_w \,\,{\parbox[c][40mm]{30mm}{
 \begin{fmfgraph*}(60,60)\setlength{\unitlength}{1cm}
\fmfstraight
   \fmftop{i1} \fmfbottom{o1,o2,o3,o4}
\fmfright{r4,r3,r2,r1}
\fmflabel{$p$}{r1}
\fmflabel{$q$}{r2}
\fmflabel{$r$}{r3}
   \fmf{plain,tension=2.0}{v1,i1}
   \fmf{phantom}{o1,v1}
   \fmf{phantom}{o4,v1}
\fmffreeze
\fmf{dashes}{v1,r1}
\fmf{dashes}{v2,r2}
\fmf{dashes}{v3,r3}
   \fmf{plain}{o1,v1}
   \fmf{plain}{o2,v2}
   \fmf{plain}{v2,v1}
   \fmf{plain}{o3,v3}
   \fmf{plain}{o4,v3}
   \fmf{plain}{v3,v2}
   \fmfdot{v1,v2,v3}
%\fmflabel{$S$}{}
 \end{fmfgraph*}
}}
$$
Thus,
for the computation of $\expect [\|X_{ts}^{\btvvva}\|^2_{\HS ,\LL(\cb_{\eta}^4;\cb_{\eta})}]$, 
we obtain the graph:
$$
 \expect  \|X_{ts}^{\btvvva}\|_{\HS ,\LL(\cb_{\eta}^3;\cb_{\eta})}^2  = 
%\\ & \qquad 
\int_s^t du
    \int_s^u dv \int_s^u dw
    % \fbox
    {\parbox[c][50mm]{50mm}{
 \begin{fmfgraph*}(200,250)
\setlength{\unitlength}{1cm}
   \fmftop{i1,i2,i3} \fmfbottom{o1,o2,o3}
  \fmf{phantom,tension=1.5}{v3,i3}
   \fmf{phantom,tension=1.5}{v4,o3}
   \fmf{tfermioni,left=1}{v3,v4}
   \fmf{dashes}{v3,v4}
   \fmf{tfermioni,right=1}{v3,v4}
   \fmf{plain,left=0.2}{v1,v3}
   \fmf{plain,right=0.2}{v2,v4}
   \fmfdot{v3,v4}
\fmfleft{l1}
%\fmffreeze  
 \fmf{phantom,tension=1.5}{v5,i1}
   \fmf{phantom,tension=1.5}{v6,o1}
   \fmf{tfermioni,left=1}{v5,v6}
   \fmf{dashes}{v5,v6}
   \fmf{tfermion,right=1}{v5,v6}
   \fmf{plain,right=0.2}{v1,v5}
   \fmf{plain,left=0.2}{v2,v6}
   \fmfdot{v5,v6}
   % \fmflabel{$S$}{}
\fmflabel{$v$}{v1}
\fmflabel{$v$}{v2}
\fmflabel{$w$}{v3}
\fmflabel{$w$}{v4}
\fmflabel{$u$}{v5}
\fmflabel{$u$}{v6}
   \fmf{phantom,tension=1.5}{v1,i2}
   \fmf{phantom,tension=1.5}{v2,o2}
   \fmf{tfermioni,right=.5}{v1,v2}
   \fmf{dashes,left=.5}{v1,v2}
   \fmf{phantom}{v1,v2}
   \fmfdot{v1,v2}
 \end{fmfgraph*}
}}
$$
Now, invoking repeatedly relation (\ref{eq:graph-ineq}), 
we can iteratively reduce the above graph to obtain a bound for 
$\expect [\|X_{ts}^{\btvvva}\|^2_{\HS ,\LL(\cb_{\eta}^4;\cb_{\eta})}]$ of the form:
\begin{equation*}
  \begin{split}
  \lesssim
\int_s^t du (u-s)^{2\Delta}
     %
    \parbox[c][30mm]{30mm}{%\fbox
    {
 \begin{fmfgraph*}(100,100)
\setlength{\unitlength}{1cm}
   \fmftop{i1} \fmfbottom{o1}
  \fmf{phantom,tension=5.5}{v1,i1}
   \fmf{phantom,tension=5.5}{v2,o1}
   \fmf{tfermioni,left=1.5}{v1,v2}
   \fmf{tfermioni,left=0.5}{v1,v2}
   \fmf{dashes,right=0.5}{v1,v2}
   \fmf{tfermion,right=1.5}{v1,v2}
    \fmfdot{v1,v2}
\fmflabel{$u$}{v1}
\fmflabel{$u$}{v2}
 \end{fmfgraph*}
}}
\lesssim \int_s^t du (t-u)^{\Delta-3/2} (u-s)^{2\Delta} \lesssim (t-s)^{3\Delta-1/2},
\end{split}
\end{equation*}
which is again what is needed for our Lemma \ref{lemma:ltwo-bounds}.

\smallskip

\noindent
{\bf (2)} {\it Case of the operator} $X^{\btvvv}$.
Using the same conventions as before,
the operator $X^{\btvvv}$ can be represented as:
$$
X_{ts}^{\btvvv} = \int_s^t d\beta^p_u \int_s^u d\beta^q_v \int_s^u d\beta^r_w \,\,{\parbox[c][40mm]{30mm}{
 \begin{fmfgraph*}(60,60)\setlength{\unitlength}{1cm}
\fmfstraight
   \fmftop{i1} \fmfbottom{o1,o2,o3,o4}
\fmfright{r4,r3,r2,r1}
\fmflabel{$p$}{r1}
\fmflabel{$q$}{r2}
\fmflabel{$r$}{r3}
   \fmf{plain,tension=2.0}{v1,i1}
   \fmf{phantom}{o1,v1}
   \fmf{phantom}{o4,v1}
\fmffreeze
\fmf{dashes}{v1,r1}
\fmf{dashes}{v2,r2}
\fmf{dashes}{v3,r3}
   \fmf{plain}{o1,v2}
   \fmf{plain}{o2,v2}
   \fmf{plain}{v2,v1}
   \fmf{plain}{o3,v3}
   \fmf{plain}{o4,v3}
   \fmf{plain}{v3,v1}
   \fmfdot{v1,v2,v3}
%\fmflabel{$S$}{}
 \end{fmfgraph*}
}}
$$
Now our current situation is slightly different from the previous ones, since in the triple
Brownian integral above, the last two are not iterated. This means that we have to handle
some sums of the form $E[(\sum_{\al}Y_\al Z_\al)^2]$ for some centered Gaussian
random variables $(Y_\al,Z_\al)_\al$ forming a Gaussian vector. The standard way to compute such sums
is to write:
\bean
E\Big[\Big(\sum_{\al}Y_\al Z_\al\Big)^2\Big]&=&\sum_{\al,\beta} E[Y_\al Z_\al Y_\beta Z_\beta]  \\
&=& \sum_{\al,\beta} E[Y_\al Z_\al] E[Y_\beta Z_\beta]  
+E[Y_\al Y_\beta] E[Z_\al Z_\beta]  +E[Y_\al Z_\beta] E[Y_\beta Z_\al].  
\eean
By extrapolating this elementary consideration to our situation, this implies that,
in the computation of $\expect \|X_{ts}^{\btvvv}\|^2_{\HS ,\LL(\cb_{\eta}^4;\cb_{\eta})}$, three different kind of contractions are involved. Hence, we also obtain three different graphs:
\begin{equation}\label{eq:complex-graphs-1}
  \begin{split}
    \expect & \|X_{ts}^{\btvvv}\|_{\HS ,\LL(\cb_{\eta}^3;\cb_{\eta})}^2  = 
\\ & \qquad \int_s^t du
    \int_s^u dv \int_s^u dw
  %  \fbox
    {\parbox[c][30mm]{50mm}{
 \begin{fmfgraph*}(200,250)
\setlength{\unitlength}{1cm}
   \fmftop{i1,i2,i3} \fmfbottom{o1,o2,o3}
  \fmf{phantom,tension=1.5}{v3,i3}
   \fmf{phantom,tension=1.5}{v4,o3}
   \fmf{tfermioni,left=1}{v3,v4}
   \fmf{dashes}{v3,v4}
   \fmf{tfermioni,right=1}{v3,v4}
   \fmf{plain,left=0.2}{v1,v3}
   \fmf{plain,right=0.2}{v2,v4}
   \fmfdot{v3,v4}
\fmfleft{l1}
%\fmffreeze  
 \fmf{phantom,tension=1.5}{v5,i1}
   \fmf{phantom,tension=1.5}{v6,o1}
   \fmf{tfermioni,left=1}{v5,v6}
   \fmf{dashes}{v5,v6}
   \fmf{tfermioni,right=1}{v5,v6}
   \fmf{plain,right=0.2}{v1,v5}
   \fmf{plain,left=0.2}{v2,v6}
   \fmfdot{v5,v6}
   % \fmflabel{$S$}{}
\fmflabel{$u$}{v1}
\fmflabel{$u$}{v2}
\fmflabel{$v$}{v3}
\fmflabel{$v$}{v4}
\fmflabel{$w$}{v5}
\fmflabel{$w$}{v6}
   \fmf{phantom,tension=1.5}{v1,i2}
   \fmf{phantom,tension=1.5}{v2,o2}
   \fmf{tfermion,right=.5}{v1,v2}
   \fmf{dashes,left=.5}{v1,v2}
 %  \fmf{tfermioni,left=1}{v1,v2}
   \fmfdot{v1,v2}
 \end{fmfgraph*}
}}
\\
 & \quad + \int_s^t du
    \int_s^u dv \int_s^u dw
    % \fbox
    {\parbox[c][40mm]{50mm}{
 \begin{fmfgraph*}(200,250)
\setlength{\unitlength}{1cm}
   \fmftop{i1,i2,i3} \fmfbottom{o1,o2,o3}
  \fmf{phantom,tension=1.5}{v3,i3}
   \fmf{phantom,tension=1.5}{v4,o3}
   \fmf{tfermioni,left=1}{v3,v4}
   \fmf{dashes,tension=0,right=.5}{v3,v5}
   \fmf{phantom}{v3,v4}
   \fmf{phantom}{v5,v6}
   \fmf{tfermioni,right=1}{v3,v4}
   \fmf{plain,left=0.2}{v1,v3}
   \fmf{plain,right=0.2}{v2,v4}
   \fmfdot{v3,v4}
\fmfleft{l1}
%\fmffreeze  
 \fmf{phantom,tension=1.5}{v5,i1}
   \fmf{phantom,tension=1.5}{v6,o1}
   \fmf{tfermioni,left=1}{v5,v6}
   \fmf{dashes,tension=0,left=.5}{v4,v6}
   \fmf{tfermioni,right=1}{v5,v6}
   \fmf{plain,right=0.2}{v1,v5}
   \fmf{plain,left=0.2}{v2,v6}
   \fmfdot{v5,v6}
   % \fmflabel{$S$}{}
\fmflabel{$u$}{v1}
\fmflabel{$u$}{v2}
\fmflabel{$w$}{v3}
\fmflabel{$v$}{v4}
\fmflabel{$w$}{v5}
\fmflabel{$v$}{v6}
   \fmf{phantom,tension=1.5}{v1,i2}
   \fmf{phantom,tension=1.5}{v2,o2}
   \fmf{tfermion,right=.5}{v1,v2}
   \fmf{dashes,left=.5}{v1,v2}
 %  \fmf{tfermioni,left=1}{v1,v2}
   \fmfdot{v1,v2}
 \end{fmfgraph*}
}}
\\
 & \quad + \int_s^t du
    \int_s^u dv \int_s^u dw
    % \fbox
    {\parbox[c][30mm]{50mm}{
 \begin{fmfgraph*}(200,250)
\setlength{\unitlength}{1cm}
   \fmftop{i1,i2,i3} \fmfbottom{o1,o2,o3}
  \fmf{phantom,tension=1.5}{v3,i3}
   \fmf{phantom,tension=1.5}{v4,o3}
   \fmf{tfermioni,left=1}{v3,v4}
   \fmf{dashes,tension=0}{v3,v6} %%%%%%%
   \fmf{phantom}{v3,v4}
   \fmf{phantom}{v5,v6}
   \fmf{tfermioni,right=1}{v3,v4}
   \fmf{plain,left=0.2}{v1,v3}
   \fmf{plain,right=0.2}{v2,v4}
   \fmfdot{v3,v4}
\fmfleft{l1}
%\fmffreeze  
 \fmf{phantom,tension=1.5}{v5,i1}
   \fmf{phantom,tension=1.5}{v6,o1}
   \fmf{tfermioni,left=1}{v5,v6}
   \fmf{dashes,tension=0}{v4,v5} %%%%%%%%
   \fmf{tfermioni,right=1}{v5,v6}
   \fmf{plain,right=0.2}{v1,v5}
   \fmf{plain,left=0.2}{v2,v6}
   \fmfdot{v5,v6}
   % \fmflabel{$S$}{}
\fmflabel{$u$}{v1}
\fmflabel{$u$}{v2}
\fmflabel{$v$}{v3}
\fmflabel{$w$}{v4}
\fmflabel{$w$}{v5}
\fmflabel{$v$}{v6}
   \fmf{phantom,tension=1.5}{v1,i2}
   \fmf{phantom,tension=1.5}{v2,o2}
   \fmf{tfermion,right=.5}{v1,v2}
   \fmf{dashes,left=.5}{v1,v2}
 %  \fmf{tfermioni,left=1}{v1,v2}
   \fmfdot{v1,v2}
 \end{fmfgraph*}
}}
\end{split}
\end{equation}
Observe that the first of those graphs already appeared in the study of  $X^{\btvvva}$. 
We will then focus on the two other ones.

\smallskip

The analysis of the second graph above can be started  by reducing again the crossed double lines, which gives a  new graph of the form:
\beq\label{eq:complex-graph-2-2}
  \int_s^t du  \int_s^u dv \int_s^u dw
    % \fbox
    {\parbox[c][30mm]{30mm}{
 \begin{fmfgraph*}(150,200)
\setlength{\unitlength}{1cm}
   \fmftop{i1,i2,i3} \fmfbottom{o1,o2,o3}
  \fmf{phantom,tension=1.5}{v3,i3}
   \fmf{phantom,tension=1.5}{v4,o3}
   %\fmf{tfermioni,left=1}{v3,v4}
   \fmf{dashes,tension=0,right=.5}{v3,v5}
   \fmf{phantom}{v3,v4}
   \fmf{phantom}{v5,v6}
   \fmf{tfermioni}{v3,v4}
   \fmf{plain,left=0.2}{v1,v3}
   \fmf{plain,right=0.2}{v2,v4}
   \fmfdot{v3,v4}
\fmfleft{l1}
%\fmffreeze  
 \fmf{phantom,tension=1.5}{v5,i1}
   \fmf{phantom,tension=1.5}{v6,o1}
   \fmf{tfermioni}{v5,v6}
   \fmf{dashes,tension=0,left=.5}{v4,v6}
  % \fmf{tfermioni,right=1}{v5,v6}
   \fmf{plain,right=0.2}{v1,v5}
   \fmf{plain,left=0.2}{v2,v6}
   \fmfdot{v5,v6}
   % \fmflabel{$S$}{}
\fmflabel{$u$}{v1}
\fmflabel{$u$}{v2}
\fmflabel{$w$}{v3}
\fmflabel{$v$}{v4}
\fmflabel{$w$}{v5}
\fmflabel{$v$}{v6}
   \fmf{phantom,tension=1.5}{v1,i2}
   \fmf{phantom,tension=1.5}{v2,o2}
   \fmf{tfermion,right=.5}{v1,v2}
   \fmf{dashes,left=.5}{v1,v2}
 %  \fmf{tfermioni,left=1}{v1,v2}
   \fmfdot{v1,v2}
 \end{fmfgraph*}
}}
\eeq
However, at this point, 
we cannot proceed as in the previous cases, with a sequence of reduction of subgraphs,
in order to prove the convergence. Indeed, in the current situation, some irreducible triangular structures like the following appear, at the top and bottom of the graph
(\ref{eq:complex-graph-2-2}):
\beq\label{eq:complex-graph-2-3}
   \int_s^u dw  {\parbox[c][50mm]{30mm}{
 \begin{fmfgraph*}(80,80)
\setlength{\unitlength}{1cm}
   \fmftop{i1} \fmfbottom{o1,o2}
\fmflabel{$i$}{i1}
\fmflabel{$j$}{o1}
\fmflabel{$k$}{o2}
\fmf{dots,tension=5.5}{i1,v1}
\fmf{plain,label=$q$}{v2,v1}
\fmf{plain,label=$q+i$}{v3,v1}
\fmf{dashes,label=$q-j$}{v2,v3}
\fmf{dots,tension=5.5}{o1,v2}
\fmf{dots,tension=5.5}{o2,v3}
\fmfdot{v1,v2,v3}
\fmflabel{$u$}{v1}
\fmflabel{$w$}{v2}
\fmflabel{$w$}{v3}
 \end{fmfgraph*}
}}
\eeq
where the dotted lines stand for the remaining part of the graph, and where we put explicit indexes on outgoing lines and on the edges of the subgraph. Note that by the rules we have
imposed on the graph, the constraint $i+j+k=0$ (we consider the dotted lines directed inwards)
holds true. The contribution of the triangular structure is thus given by
$$
 \delta_{i+j+k=0}  \int_s^u dw \sum_{q} \lambda_{q+k}^{-\nu} e^{-\lambda_{q}(u-w)-\lambda_{q-i}(u-w)},
$$
and we can bound this last expression by:
$$
\lesssim \delta_{i+j+k=0}  \int_s^u \frac{dw}{(u-w)^b} \sum_{q} \lambda_{q+k}^{-\nu} \lambda_{q}^{-b/2} \lambda_{q-i}^{-b/2},
$$
for some $b\in(0,1)$. The latter sum is finite when $\nu+b > 1/2$, which means that, by
choosing  $b = 1/2-\nu+\eps$, the triangular structure yields a bound of order
$$
\lesssim \delta_{i+j+k=0} (u-s)^{1-b}.
$$
Summarizing the previous discussion, it is now easily seen that the structure
(\ref{eq:complex-graph-2-3}) behaves like a simple vertex, up to an appropriate factor of $(u-s)$:
$$
   \int_s^u dw  {\parbox[c][30mm]{40mm}{
 \begin{fmfgraph*}(80,80)
\setlength{\unitlength}{1cm}
   \fmftop{i1} \fmfbottom{o1,o2}
\fmflabel{$i$}{i1}
\fmflabel{$j$}{o1}
\fmflabel{$k$}{o2}
\fmf{dots,tension=5.5}{i1,v1}
\fmf{plain,label=$q$}{v2,v1}
\fmf{plain,label=$q+i$}{v3,v1}
\fmf{dashes,label=$q-j$}{v2,v3}
\fmf{dots,tension=5.5}{o1,v2}
\fmf{dots,tension=5.5}{o2,v3}
\fmfdot{v1,v2,v3}
\fmflabel{$u$}{v1}
\fmflabel{$w$}{v2}
\fmflabel{$w$}{v3}
 \end{fmfgraph*}
}}
 \lesssim (u-s)^{\Delta} {\parbox[c][30mm]{25mm}{
 \begin{fmfgraph*}(50,50)
\setlength{\unitlength}{1cm}
   \fmftop{i1} \fmfbottom{o1,o2}
\fmflabel{$i$}{i1}
\fmflabel{$j$}{o1}
\fmflabel{$k$}{o2}
\fmf{dots,tension=2.5}{i1,v1}
\fmf{dots,tension=2.5}{o1,v1}
\fmf{dots,tension=2.5}{o2,v1}
\fmfdot{v1}
\fmflabel{$u$}{v1}
 \end{fmfgraph*}
}}
$$
Using this fact, we can reduce our graph (\ref{eq:complex-graph-2-2}) to the following
simpler structure:
$$
 \lesssim  \int_s^t du (u-s)^{2\Delta}  {\parbox[c][30mm]{35mm}{
 \begin{fmfgraph*}(100,200)
\setlength{\unitlength}{1cm}
   \fmftop{i1} \fmfbottom{o1}
\fmf{phantom,tension=3}{i1,v1}
\fmf{phantom,tension=3}{o1,v2}
\fmf{tfermioni,left=1.5}{v1,v2}
\fmf{tfermioni,left=0.5}{v1,v2}
\fmf{tfermion,right=0.5}{v1,v2}
\fmf{dashes,right=1.5}{v1,v2}
\fmfdot{v1,v2}
\fmflabel{$u$}{v1}
\fmflabel{$u$}{v2}
 \end{fmfgraph*}
}}
\lesssim (t-s)^{3\Delta-1/2},
$$
where the last bound has been obtained similarly to (\ref{eq:graph-1}).

\smallskip

Finally, let us associate a bound to the 
third graph in (\ref{eq:complex-graphs-1}): first, notice that this third graph
can be reduced to
\beq\label{eq:complex-graph-2-4}
\int_s^t du
    \int_s^u dv \int_s^u dw
   %  \fbox
    {\parbox[c][50mm]{50mm}{
 \begin{fmfgraph*}(200,200)
\setlength{\unitlength}{1cm}
   \fmftop{i1,i2,i3} \fmfbottom{o1,o2,o3}
  \fmf{phantom,tension=1.5}{v3,i3}
   \fmf{phantom,tension=1.5}{v4,o3}
 %  \fmf{tfermioni,left=1}{v3,v4}
   \fmf{dashes,tension=0,left=0}{v3,v6} %%%%%%%
   \fmf{phantom}{v3,v4}
   \fmf{phantom}{v5,v6}
   \fmf{tfermioni}{v3,v4}
   \fmf{plain,left=0.2}{v1,v3}
   \fmf{plain,right=0.2}{v2,v4}
   \fmfdot{v3,v4}
\fmfleft{l1}
%\fmffreeze  
 \fmf{phantom,tension=1.5}{v5,i1}
   \fmf{phantom,tension=1.5}{v6,o1}
%   \fmf{tfermioni,left=1}{v5,v6}
   \fmf{dashes,tension=0,right=1.3}{v4,v5} %%%%%%%%
   \fmf{tfermioni}{v5,v6}
   \fmf{plain,right=0.2}{v1,v5}
   \fmf{plain,left=0.2}{v2,v6}
   \fmfdot{v5,v6}
   % \fmflabel{$S$}{}
\fmflabel{$u$}{v1}
\fmflabel{$u$}{v2}
\fmflabel{$v$}{v3}
\fmflabel{$w$}{v4}
\fmflabel{$w$}{v5}
\fmflabel{$v$}{v6}
   \fmf{phantom,tension=1.5}{v1,i2}
   \fmf{phantom,tension=1.5}{v2,o2}
   \fmf{tfermion,right=.5}{v1,v2}
   \fmf{dashes,left=.5}{v1,v2}
 %  \fmf{tfermioni,left=1}{v1,v2}
   \fmfdot{v1,v2}
 \end{fmfgraph*}
}}
\eeq
Furthermore, this graph contains the irreducible subgraph:
$$
    \int_s^u dv \int_s^u dw
    % \fbox
    {\parbox[c][50mm]{50mm}{
 \begin{fmfgraph*}(200,200)
\setlength{\unitlength}{1cm}
   \fmfleft{i1} \fmfright{o1}
\fmftop{t1,t2}
\fmfbottom{b1,b2}
\fmf{phantom,tension=1}{v1,t1}
\fmf{phantom,tension=1}{v3,t2}
\fmf{phantom,tension=1}{v2,b1}
\fmf{phantom,tension=1}{v4,b2}
\fmf{dots,tension=2,label=$i$}{i1,vi1}
\fmf{dots,tension=2,label=$i$}{o1,vo1}
   \fmf{plain,label=$j$}{v1,vi1}
   \fmf{plain,label=$k$}{v2,vi1}
   \fmf{tfermioni,label=$j-l$}{v3,v1}
   \fmf{tfermioni,label=$k-q$}{v2,v4}
   \fmf{plain,label=$j-l+q$}{vo1,v3}
   \fmf{plain,label=$k-q+l$}{vo1,v4}
   \fmf{dashes,tension=0,left=0.3,label=$q$}{v2,v3}
   \fmf{dashes,tension=0,left=-0.3,label=$l$}{v4,v1}
   \fmfdot{v1,v2,v3,v4,vi1,vo1}
\fmflabel{$v$}{v1}
\fmflabel{$w$}{v2}
\fmflabel{$v$}{v3}
\fmflabel{$w$}{v4}
\fmflabel{$u$}{vi1}
\fmflabel{$u$}{vo1}
 \end{fmfgraph*}
}}
$$
which corresponds to the expression
\begin{multline*}
D\equiv
    \int_s^u dv \int_s^u dw \\ 
    \times \sum_{j+k=i}\sum_{q,l} \frac{e^{-\lambda_j(u-v)-\lambda_k(u-w)-\lambda_{j-l+q}(u-v)-\lambda_{k-q+l}(u-w)-\lambda_{j-l}(w+v-2s)-\lambda_{k-q}(w+v-2s)}}{\lambda_l^\nu \lambda_q^\nu \lambda_{j-l}^{2\eta} \lambda_{k-q}^{2\eta}}.
\end{multline*}
Introducing an additional parameter $b$ and bounding the exponential terms as usual gives:
$$
D
\lesssim    \int_s^u \frac{dv}{(u-v)^b} \int_s^u \frac{dw}{(u-w)^b} \sum_{j+k=i}\sum_{q,l} \frac{1}{\lambda_j^{b} \lambda_k^{b} \lambda_l^\nu \lambda_q^\nu \lambda_{j-l}^{2\eta} \lambda_{k-q}^{2\eta}}.
$$
Now, using Lemma~\ref{lemma:asym-sum}, we can bound the sums over $q$ and $l$ in order
to obtain:
$$
D \lesssim    \int_s^u \frac{dv}{(u-v)^b} \int_s^u \frac{dw}{(u-w)^b} \sum_{j+k=i} \frac{1}{\lambda_j^{b+\nu} \lambda_k^{b+\nu}}.
$$
Thus, choosing $b=2\eta-\nu$, we end up with:
$$
D \lesssim   (u-s)^{2\Delta} \lambda_i^{-2\eta},
$$
which means that we have obtained the graphical inequality:
$$
    \int_s^u dv \int_s^u dw
    % \fbox
    {\parbox[c][50mm]{50mm}{
 \begin{fmfgraph*}(150,150)
\setlength{\unitlength}{1cm}
   \fmfleft{i1} \fmfright{o1}
\fmftop{t1,t2}
\fmfbottom{b1,b2}
\fmf{phantom,tension=1}{v1,t1}
\fmf{phantom,tension=1}{v3,t2}
\fmf{phantom,tension=1}{v2,b1}
\fmf{phantom,tension=1}{v4,b2}
\fmf{phantom,tension=2}{i1,vi1}
\fmf{phantom,tension=2}{o1,vo1}
   \fmf{plain}{v1,vi1}
   \fmf{plain}{v2,vi1}
   \fmf{tfermioni}{v3,v1}
   \fmf{tfermioni}{v2,v4}
   \fmf{plain}{vo1,v3}
   \fmf{plain}{vo1,v4}
   \fmf{dashes,tension=0,left=0}{v2,v3}
   \fmf{dashes,tension=0,left=-0}{v4,v1}
   \fmfdot{v1,v2,v3,v4,vi1,vo1}
\fmflabel{$v$}{v1}
\fmflabel{$w$}{v2}
\fmflabel{$v$}{v3}
\fmflabel{$w$}{v4}
\fmflabel{$u$}{vi1}
\fmflabel{$u$}{vo1}
 \end{fmfgraph*}
}}
\lesssim
(u-s)^{2\Delta}
    \parbox[c][30mm]{30mm}{%\fbox
    {
 \begin{fmfgraph*}(50,50)
\setlength{\unitlength}{1cm}
   \fmftop{i1} \fmfbottom{o1}
   \fmf{tfermioni,left=0.5}{i1,o1}
   \fmfdot{i1,o1}
\fmflabel{$u$}{o1}
\fmflabel{$u$}{i1}
 \end{fmfgraph*}
}}
$$
Plugging this representation in the complete graph (\ref{eq:complex-graph-2-4}), we obtain:
$$
\lesssim \int_s^t du (u-s)^{2\Delta}
  \parbox[c][30mm]{30mm}{%\fbox
    {
 \begin{fmfgraph*}(100,100)
\setlength{\unitlength}{1cm}
   \fmftop{i1} \fmfbottom{o1}
  \fmf{phantom,tension=5.5}{v1,i1}
   \fmf{phantom,tension=5.5}{v2,o1}
   \fmf{tfermioni,left=1.0}{v1,v2}
  % \fmf{tfermioni,left=0.5}{v1,v2}
   \fmf{dashes,right=0.0}{v1,v2}
   \fmf{tfermion,right=1.0}{v1,v2}
    \fmfdot{v1,v2}
\fmflabel{$u$}{v1}
\fmflabel{$u$}{v2}
 \end{fmfgraph*}
}}
\lesssim (t-s)^{3\Delta -1/2}.
$$

Going back to Lemma \ref{lemma:ltwo-bounds}, we should still treat the case 
of $X^\tau$ for $\tau = \btvvvb, \btvvvc, \btvvvd$. But these estimates are now mere
variations of the previous ones, and are left to the reader for sake of conciseness.





%%%%%%%%%%%%%%%%%%%%%%%%%%%%%%%%%%%%%%%%%%%%

\newcommand{\SortNoop}[1]{} \providecommand{\urlbib}[1]{\texttt{\small #1}}
\begin{thebibliography}{99}

\bibitem{Ad}
R.~A. Adams.
\newblock {\em Sobolev spaces}.
\newblock Academic Press, 1975.

\bibitem{CQ}
L. Coutin and Z. Qian.
\newblock Stochastic rough path analysis and fractional Brownian motion.
\newblock {\em Probab. Theory Relat. Fields} 122:108-140, 2002.

\bibitem{Da}R. Dalang. 
\newblock{Extending martingale measure stochastic integral with applications to spatially homogeneous S.P.D.E's},
\newblock {\em Electron. J. Probab.} 4, Paper No.6, 29 p., 1999 (electronic).

\bibitem{DPZ}
G.~Da~Prato and J.~Zabczyk.
\newblock {\em Stochastic equations in infinite dimensions}, volume~44 of {\em
  Encyclopedia of Mathematics and its Applications}.
\newblock Cambridge University Press, Cambridge, 1992.

\bibitem{engel00a}
K.-J. Engel and R.~Nagel.
\newblock {\em One-parameter semigroups for linear evolution equations}, volume
  194 of {\em Graduate Texts in Mathematics}.
\newblock Springer-Verlag, New York, 2000.

\bibitem{fattorini99a}
H.~O. Fattorini.
\newblock {\em Infinite-dimensional optimization and control theory}, volume~62
  of {\em Encyclopedia of Mathematics and its Applications}.
\newblock Cambridge University Press, Cambridge, 1999.

\bibitem{FP}
D. Feyel and A. de La Pradelle.
\newblock Curvilinear integrals along enriched paths.
\newblock {\em Electron. J. Probab.} 11:860-892, 2006.

\bibitem{FV}
P. Friz and N. Victoir.
\newblock {\em Multidimensional dimensional processes seen as rough paths}.
\newblock Cambridge University Press, to appear.

\bibitem{rough}
M.~Gubinelli.
\newblock Controlling rough paths.
\newblock {\em Jour. Funct. Anal.} 216:86-140, 2004.

\bibitem{g-ramif}
M.~Gubinelli.
\newblock Ramification of rough paths. Preprint, 2006.

\bibitem{g2}
M.~Gubinelli. 
\newblock Rough integrals in higher dimensions.
\newblock In preparation.

\bibitem{glt}
M.~Gubinelli, A.~Lejay and S.~Tindel.
\newblock Young integrals and SPDEs
\newblock {\em Pot. Anal.} 25:307--326, 2006.

\bibitem{Lej}
A. Lejay.
\newblock An introduction to rough paths.
\newblock In {\em S{\'e}minaire de Probabilit{\'e}s 37}, volume 1832 of
  {\em Lecture Notes in Mathematics}, pages 1--59. Springer-Verlag Heidelberg,
  2003.

\bibitem{Ledoux-Lyons}
M.~Ledoux, T.~Lyons, and Z.~Qian.
\newblock L\'evy area of {W}iener processes in Banach spaces.
\newblock {\em Ann. Probab.}, 30(2):546--578, 2002.

\bibitem{LS}
J. Le\'on and J. San Martin.
\newblock Linear stochastic differential equations driven by a fractional Brownian motion with Hurst parameter less than $1/2$.
\newblock To appear at {\em Stoch. And Stoch. Reports}.

\bibitem{LR}
S. Lototsky and B. Rozovsky.
\newblock Wiener Chaos Solutions of Linear Stochastic Evolution Equations. 
\newblock {\em Annals of Probability} 34(2):638-662, 2006.

\bibitem{LyonsBook}
T.~Lyons and Z.~Qian.
\newblock {\em System control and rough paths}.
\newblock Oxford University Press, 2002.

\bibitem{Lyons}
T.~ Lyons.
\newblock Differential equations driven by rough signals.
\newblock {\em Rev. Mat. Iberoamericana}, 14(2):215--310, 1998.

\bibitem{Maslo-Nualart}
B. Maslowski and D. Nualart.
\newblock Evolution equations driven by a fractional {B}rownian motion.
\newblock {\em J. Funct. Anal.}, 202(1):277--305, 2003.

\bibitem{Nbk}
D. Nualart.
\newblock {\em Malliavin calculus and related topics}.
\newblock Springer, 1995.

\bibitem{paz}
A.~Pazy.
\newblock {\em Semigroups of linear operators and applications to partial
  differential equations}, volume~44 of {\em Applied Mathematical Sciences}.
\newblock Springer-Verlag, New York, 1983.

\bibitem{PT}
V. P\'erez-Abreu and C. Tudor.
\newblock Transfer principle for stochastic fractional integral.
\newblock {\em Bol. Soc. Mat. Mexicana} 8:55-71, 2002.

\bibitem{qt}
Ll. Quer and S. Tindel.
\newblock 
The 1-d stochastic wave equation driven by a fractional Brownian motion.
\newblock To appear at {\em Stoch. Processes  Appl.}  117(10):1448--1472, 2007.

\bibitem{Ti}
S. Tindel.
\newblock 
Stochastic parabolic equations with anticipative initial condition
\newblock  {\em Stoch. and Stoch. Reports} 62:1-20, 1997.

\bibitem{TTV}
S.~Tindel, C.~A. Tudor, and F.~Viens.
\newblock Stochastic evolution equations with fractional {B}rownian motion.
\newblock {\em Probab. Theory Related Fields}, 127(2):186--204, 2003.

\bibitem{Wa}
J.~B. Walsh.
\newblock An introduction to stochastic partial differential equations.
\newblock In {\em \'Ecole d'\'et\'e de probabilit\'es de Saint-Flour,
  XIV---1984}, volume 1180 of {\em Lecture Notes in Math.}, pages 265--439.
  Springer, Berlin, 1986.

\bibitem{Young}
L.~C. Young.
\newblock An inequality of H\"older type, connected with Stieljes integration.
\newblock {\em Acta Math.}, 67:251--282, 1936.

\bibitem{Za}
M.~Z{\"a}hle.
\newblock Integration with respect to fractal functions and stochastic
  calculus. {I}.
\newblock {\em Probab. Theory Related Fields}, 111(3):333--374, 1998.

\end{thebibliography}
\end{fmffile}
\end{document}
